\documentclass[11pt]{article}

\usepackage[margin=1in]{geometry}
\usepackage{amsmath,amssymb,amsthm}
\usepackage{mathtools}
\usepackage{enumitem}
\usepackage[hidelinks]{hyperref}
\usepackage[expansion=false]{microtype}

\newtheorem{theorem}{Theorem}[section]
\newtheorem{lemma}[theorem]{Lemma}
\newtheorem{proposition}[theorem]{Proposition}
\newtheorem{corollary}[theorem]{Corollary}
\newtheorem{conjecture}[theorem]{Conjecture}
\theoremstyle{definition}
\newtheorem{definition}[theorem]{Definition}
\newtheorem{hypothesis}[theorem]{Hypothesis}
\newtheorem{remark}[theorem]{Remark}

\DeclareMathOperator{\lcm}{lcm}
\newcommand{\NN}{\mathbb{N}}
\newcommand{\ZZ}{\mathbb{Z}}
\newcommand{\QQ}{\mathbb{Q}}
\newcommand{\PP}{\mathcal{P}}
\newcommand{\der}{{}'}
\newcommand{\dd}{\,\mathrm{d}}

\title{Abridged: On the equation $n''=n$ and a problem of Erd\H{o}s and Barbeau \\
       on products of prime--reciprocal sums}
\author{Vico Bonfioli\thanks{\texttt{vicobonfioli@gmail.com}.
  See the acknowledgments for a candid note on the use of AI assistance in this work.}}
\date{\today}

\begin{document}
\maketitle

\begin{abstract}
Erd\H{o}s Problem~\#307 (posed by Barbeau, 1976; recorded in Erd\H{o}s--Graham, 1980) asks
whether there exist two finite sets $P,Q$ of primes with
$\bigl(\sum_{p\in P}\tfrac1p\bigr)\bigl(\sum_{q\in Q}\tfrac1q\bigr)=1$.
Writing $a=\prod P$, $b=\prod Q$, a solution exists if and only if the \emph{arithmetic derivative}
has a two--cycle $a'=b,\ b'=a$ with $a\neq b$; equivalently iff $n''=n$ has a solution that is not a
fixed point $p^{p}$. That reformulation, and the equivalence with the derivative, are due to
Ufnarovski--\AA hlander \cite{UfnarovskiAhlander2003}, whose Conjecture~4 asserts in exactly these
terms that $\bigl(\sum1/p_i\bigr)\bigl(\sum1/q_j\bigr)=1$ has no solution in distinct primes. Our
starting point is that \#307 \emph{is} that conjecture --- an identification neither literature
records --- which ties the problem to the Ufnarovski--\AA hlander dynamics: their conjecture that
every $n''=n$--cycle has period one would force \#307 to have answer \textsc{no}, while a
\textsc{yes} would refute it. Exploiting the equivalence we prove a \emph{rigidity lemma} (each prime--reciprocal
sum is automatically in lowest terms, so $N_P=D_Q$, $N_Q=D_P$, and $P\cap Q=\varnothing$) and a
quantitative \emph{barrier}: building on Rosen's bound $|P\cup Q|\ge 59$, any solution has
$\min(\prod P,\prod Q)\ge 2.09\times10^{56}$, for \emph{both} members. This is the first valid
unconditional bound of its kind. The prior claims are Kovi\v{c}'s: his Proposition~17 is sound but
conditional on the smaller member being odd, while Proposition~18 --- the only unconditional
cardinality and size bound in the literature, giving $r+s\ge34$ and
$\max\{m,n\}\ge\prod_{i\le17}p_i$ --- rests on an invalid step, deriving $p_1q_s<rs$ by multiplying
an upper bound against a lower one, where the available products give only $p_1q_1<rs<p_rq_s$ and no
cardinality bound at all; we give an explicit model and machine--check both points. Against the
strongest surviving unconditional bound, that paper's computer search $x>10^4$, the gain is $52$
orders of magnitude. A finite verification over the complete list of
$49{,}961$ admissible $59$--prime supports then closes the $59$--prime level: any solution in fact
has $|P\cup Q|\ge 60$ and $\min(\prod P,\prod Q)>3.50\times10^{57}$; a quadratic--reciprocity
certificate, requiring no factorisation, then proves the canonical tail family (the first $59$
primes plus any $q$) empty, and together with a factoring pass $46{,}483$ of the $49{,}961$
one--new--prime families at level $60$ ($93\%$; the residual is factoring--hard). We add a parity dichotomy (a hypothetical all--odd
cycle would need $\ge 1412$ primes and have $\min(\prod P,\prod Q)>10^{2519}$), a proof that longer
cycles are strictly harder ($k=3$ already forces $361{,}139$ primes), and a function--field theorem:
over $\mathbb{F}_q[t]$ the derivative has \emph{no} cycles, so the obstruction over $\mathbb{Z}$ is
purely archimedean. A local congruence anatomy of the cycle equation, the Erd\H{o}s--Wintner fact
that the abundance density $0.0420$ is a genuine constant, and a Th\=abit/Euler--type constructive
calculus (the \emph{Frame
Rule} and a divisor--trick) that faithfully parametrises solutions but, by an exact identity
$\Delta_0=M_0N(1-s_{M_0}s_N)$, provably does \emph{not} reduce the difficulty below \#307 itself
(the natural ``make $\Delta_0$ small'' objective is \#307 in disguise). No fixed--shape family can
help either: for families of \emph{any} growth rate, not merely polynomial ones, a product identity
holding infinitely often forces every member constant, which closes the Th\=abit/Euler programme
(primes $3\cdot2^{n}-1$ and relatives) that is the classical technique for problems of this kind. Recast as a bilinear breeder it yields a conditional reduction
$(\mathrm{H})\Rightarrow$\textsc{yes} to an explicit prime--density hypothesis, which a quantitative
ledger shows is not a feasible finite search. On the counting side we prove that both Pythagorean layers have density zero
unconditionally \emph{on the squarefree stratum}, which is where every tool in the proof lives and
where a \#307 solution lies by rigidity; the unrestricted statement is not proved. Both Pythagorean layers have density zero on the squarefree stratum by a square--sieve argument; that argument, its conditional rate and the three obstructions to the sharp exponent $N^{1/2}$ are given in the full version and are not reproduced here. The $34$ reciprocity-immune families at
level $60$ are upgraded from probable to proved by Atkin--Morain certificates, archived for
independent verification. The rigidity result, the barrier, and the arithmetic core of the counting
argument are formalised in Lean~4 (mathlib), extremality proved rather than assumed; the only
non--logical input is a verified evaluation of the first $59$ primes. The problem remains open. The heuristic case for existence is weaker than it appears: the constant in
the expected-count heuristic cannot be measured even in principle, since the barrier places the
entire relevant region above $10^{112.9}$ while the largest value of $\sigma(a)\sigma(a')$ observed
below $10^{7}$ is $0.5535$.
\end{abstract}


\begin{center}\small\textit{Abridged version.} This note omits the deviation ladder, the plus--layer
analysis and the square sieve, together with their proofs, and cites their statements by the
numbering of the full version, which is archived with the Lean development and the computational
scripts. Nothing proved here depends on a proof appearing only there.\end{center}

\section{Introduction}\label{sec:intro}

Let $P,Q$ be finite, nonempty sets of (distinct) primes. Erd\H{o}s Problem~\#307
\cite{ErdosProblems307}, attributed to Barbeau's 1976 ``Computer challenge corner'' problem
\#477 \cite{Barbeau1976} and listed in Erd\H{o}s--Graham \cite{ErdosGraham1980}, asks:
\begin{quote}
do there exist finite sets of primes $P,Q$ with
$\displaystyle\Bigl(\sum_{p\in P}\frac1p\Bigr)\Bigl(\sum_{q\in Q}\frac1q\Bigr)=1$\,?
\end{quote}
The problem is verifiable in the sense that a single example would settle it. No example is known.
The published partial facts are a remark of J.~Rosen on MathOverflow \cite{MO320838}, that any
solution needs at least $59$ primes in $P\cup Q$, and the problem page \cite{ErdosProblems307},
which asserts $|P\cup Q|\ge60$.

The relation between the two is worth stating precisely, since this paper proves the second. The
page's justification reads, in full, that $P$ and $Q$ are disjoint and
$\sum_{p\in P\cup Q}1/p\ge2$, ``whence $|P\cup Q|\ge60$''. That implication does not hold as
displayed: the reciprocals of the first $59$ primes already sum to $2.00235>2$, so a mass of $2$ is
attainable on $59$ primes and the stated hypotheses yield $|P\cup Q|\ge59$, which is Rosen's bound.
Reaching $60$ requires excluding the $49{,}961$ admissible $59$--element supports one at a time,
which is what Proposition~\ref{prop:close59} does and what \texttt{erdos307\_sixty} machine-checks.
We therefore claim no priority for the \emph{statement} $|P\cup Q|\ge60$, which predates this work on
the problem page, and we do claim the proof. The known
``solutions'' of Cambie and others solve only the \emph{weakened} version in which $1$ is permitted
as an element, e.g.\ $(1+\tfrac15)(\tfrac12+\tfrac13)=1$; these do not bear on the prime version.

Expanding the product exhibits \#307 as a rigidly shaped Egyptian fraction for $1$:
$1=\sum_{p\in P}\sum_{q\in Q}1/(pq)$, whose denominators are the \emph{rectangular grid} $P\times Q$
of semiprimes; a \emph{rank--one}, multiplicatively separable representation. The unrestricted
companion (representing $1$ by reciprocals of distinct two--prime--factor numbers, with no product
structure imposed) is by contrast \emph{solved}: Johnson (1978) gave a $48$--term decomposition,
Watanabe \cite{Watanabe2020} $47$--term ones, and Graham's survey \cite{GrahamEgyptian} records exactly
this contrast, noting after Barbeau that the product form ``is not known.'' Thus \#307 is precisely the
rank--one restriction of a solved problem. Nor is the density--zero nature of the prime denominators the
obstruction in itself: Bloom \cite{BloomShifted} represents every positive rational by distinct unit
fractions drawn from the (equally sparse) shifted primes $\{p-1\}$. What resists is the rank--one
rigidity (the demand that the representation \emph{factor} as a product of two prime--reciprocal
sums) which the barrier (Theorem~\ref{thm:barrier}) shows costs at least $59$ primes, hence of order
$870$ semiprime terms against Johnson's $47$.

Our starting point is an identification of \#307 with a question about the \emph{arithmetic
derivative} $n\mapsto n'$, the unique map with $0'=1'=0$, $p'=1$ for primes $p$, and the Leibniz
rule $(ab)'=a'b+ab'$. (The map goes back to J.~Mingot Shelly (1911) and appeared as a 1950 Putnam
problem; Barbeau \cite{Barbeau1961} gave the first systematic study, with the modern account in
Ufnarovski--\AA hlander \cite{UfnarovskiAhlander2003} and OEIS \cite{OEIS_A003415}.) For
$n=\prod_i p_i^{e_i}$ one has $n'=n\sum_i e_i/p_i$. We show (Section~\ref{sec:bridge}) that \#307
is \emph{equivalent} to the existence of a two--cycle of $n\mapsto n'$ that is not a fixed point.
The two--cycle reformulation itself is not new, and it is older than is generally recorded.
Ufnarovski--\AA hlander \cite{UfnarovskiAhlander2003} prove that a period--two point of $n\mapsto n'$
has squarefree members with disjoint supports, and deduce that the period--two case of their
Conjecture~3 is equivalent to their Conjecture~4: $\bigl(\sum_{i\le k}1/p_i\bigr)
\bigl(\sum_{j\le l}1/q_j\bigr)=1$ has no solution in distinct primes. That is \#307, in the
derivative literature, in 2003. It was restated ``in disguise'' by Seva \cite{Seva323194} in 2019
(via $Q\mapsto Q\sum_{p\mid Q}1/p$, which agrees with the arithmetic derivative on squarefree
integers, where the question lives), who is the first to name it as \#307, and obtained
independently by Bado \cite{Bado2026}; see Section~\ref{sec:prior}. What is new here is the
identification of the two problems --- \#307 is Ufnarovski--\AA hlander's Conjecture~4 --- and the
resulting two--way bridge, importing their period--one conjecture and exporting our bounds.
Remarkably, Barbeau authored foundational papers on both sides (1961 and 1976
\cite{Barbeau1961,Barbeau1976}).

The bridge is productive in both directions.
\begin{itemize}[leftmargin=1.4em]
\item \emph{Importing} from the derivative literature: the Ufnarovski--\AA hlander theorem that any
two--cycle has squarefree members with disjoint supports \cite{UfnarovskiAhlander2003} makes the
``finite sets of distinct primes'' hypothesis of \#307 automatic, and links \#307 to the
Ufnarovski--\AA hlander conjecture that all derivative cycles have period one.
\item \emph{Exporting} to that literature: our barrier (Theorem~\ref{thm:barrier}) is, we believe,
the first valid unconditional bound of its kind. Kovi\v{c} \cite{Kovic2012} states three bounds on an $n''=n$ solution. A computer
search gives $x>10^4$, unconditionally. Proposition~17 gives $>3\cdot5\cdots29\approx3.23\times10^9$,
but only when the \emph{smaller} member is odd; Kovi\v{c} notes explicitly that no estimate follows
otherwise. Proposition~18 claims the unconditional $r+s\ge34$ and $\max\{m,n\}\ge\prod_{i=1}^{17}p_i
\approx1.92\times10^{21}$; its proof does not establish them, for the reason set out in
Section~\ref{sec:prior}. Against the strongest \emph{valid unconditional} bound, $10^4$, our
$\min(\prod P,\prod Q)\ge2.09\times10^{56}$ is a gain of $52$ orders of magnitude; against the
conditional Proposition~17 it is $47$, in the case where that proposition applies. On cardinality
there is no prior bound to improve: $|U|\ge60$ appears to be the first, and it re-derives
Kovi\v{c}'s claimed $r+s\ge34$ as a corollary by an unrelated route.
\end{itemize}

\paragraph{Results and organisation.}
Section~\ref{sec:rigidity} proves the Rigidity Lemma (Theorem~\ref{thm:rigidity}): the fraction
$\sum_{p\in S}1/p$ is automatically in lowest terms, which forces, for any solution, the exact
identities $N_P=D_Q$, $N_Q=D_P$, disjointness $P\cap Q=\varnothing$, and $D_Q=s\,D_P$ where
$s=\sum_{p\in P}1/p$. The disjointness core was anticipated by R.~Israel \cite{MO320838}; the full
rigidity is what powers everything later. Section~\ref{sec:bridge} establishes the bridge.
Section~\ref{sec:barrier} proves the barrier, the parity dichotomy, and that longer cycles are
strictly harder (so two--cycles are extremal). Section~\ref{sec:anatomy} gives a local congruence
anatomy, a heuristic model, and a proof (via Erd\H{o}s--Wintner) that the relevant abundance density
is a genuine constant. Section~\ref{sec:ff} adds a function--field lens: over $\mathbb{F}_q[t]$ the
derivative has \emph{no} cycles, so the obstruction in $\mathbb{Z}$ is purely archimedean.
Section~\ref{sec:frame} develops the Frame Rule and divisor--trick,
and, crucially, records the exact identity that renders the natural constructive objective
circular. Section~\ref{sec:breeder} recasts the rule as a bilinear breeder, records a conditional
reduction $(\mathrm{H})\Rightarrow$ \textsc{yes} to an explicit prime--density hypothesis, and shows
quantitatively why $(\mathrm{H})$ does not amount to a feasible finite search.
Section~\ref{sec:computations} documents the computations, and Section~\ref{sec:lean} the
Lean~4 formalisation. We are explicit throughout about what is proved, what is heuristic, and what
is borrowed; the problem remains open.

\paragraph{Notation.}
For a finite set $S$ of distinct primes put
\[
D_S \coloneqq \prod_{p\in S} p, \qquad N_S \coloneqq \sum_{p\in S} \frac{D_S}{p},
\qquad\text{so}\qquad \sum_{p\in S}\frac1p = \frac{N_S}{D_S}.
\]
We call $N_S$ the \emph{cofactor sum}. For a solution we write $a=D_P$, $b=D_Q$,
$s=N_P/D_P=\sum_{p\in P}1/p$ and $t=N_Q/D_Q=\sum_{q\in Q}1/q$, so $st=1$, and (without loss of
generality) $s>1>t$. We write $\omega(n)$ for the number of distinct primes dividing $n$.

\section{The Rigidity Lemma}\label{sec:rigidity}

\begin{lemma}[Rigidity / lowest terms]\label{thm:rigidity}
For a finite set $S$ of distinct primes, $\gcd(N_S,D_S)=1$. Equivalently, $\sum_{p\in S}1/p$ is
already in lowest terms, and $v_p\!\bigl(\sum_{r\in S}1/r\bigr)=-1$ for every $p\in S$.
\end{lemma}

\begin{proof}
The prime divisors of $D_S=\prod_{r\in S}r$ are exactly the elements of $S$, so it suffices to show
$p\nmid N_S$ for each $p\in S$. Fix $p\in S$ and reduce $N_S=\sum_{r\in S}D_S/r$ modulo $p$. For
$r\ne p$ the term $D_S/r=\prod_{u\in S\setminus\{r\}}u$ still contains the factor $p$, so
$p\mid D_S/r$. The single surviving term is $D_S/p=\prod_{u\in S\setminus\{p\}}u$, whence
\[
N_S \equiv \prod_{u\in S\setminus\{p\}} u \pmod p.
\]
Each $u$ in this product is a prime distinct from $p$, so $p\nmid u$; as $p$ is prime it does not
divide the product. Thus $N_S\not\equiv0\pmod p$, i.e.\ $p\nmid N_S$. As $p$ ranged over all prime
divisors of $D_S$, $\gcd(N_S,D_S)=1$. The valuation statement is the same computation:
$v_p(N_S/D_S)=v_p(N_S)-v_p(D_S)=0-1=-1$.
\end{proof}

\begin{remark}
Distinctness and primality are both essential: $v_p$ of the cofactor of a repeated prime need not be
$-1$, and the final step uses Euclid's lemma. The case $|S|=1$ is consistent ($N_{\{p\}}=1$).
\end{remark}

\begin{theorem}[Solution structure]\label{thm:structure}
Suppose $P,Q$ are finite sets of distinct primes with $\dfrac{N_P}{D_P}\cdot\dfrac{N_Q}{D_Q}=1$.
Then
\[
N_P=D_Q,\qquad N_Q=D_P,\qquad P\cap Q=\varnothing,
\]
$N_P$ is squarefree with prime support exactly $Q$ (and $N_Q$ squarefree with support $P$), and
$D_Q=s\,D_P$ where $s=N_P/D_P$.
\end{theorem}

\begin{proof}
Cross-multiplying, $N_PN_Q=D_PD_Q$. By Lemma~\ref{thm:rigidity} both $N_P/D_P$ and $N_Q/D_Q$ are in
lowest terms. Since $N_Q/D_Q=D_P/N_P$ and the lowest--terms representative of a positive rational is
unique, comparing $N_Q/D_Q$ with $D_P/N_P$ gives $N_Q=D_P$ and $D_Q=N_P$.

\emph{Disjointness} (the Israel mechanism \cite{MO320838}). If $r\in P\cap Q$ then $r\mid D_P$ and,
by $N_P=D_Q=\prod Q$, also $r\mid N_P$, but $\gcd(N_P,D_P)=1$ forces $r\nmid N_P$, a contradiction.
(Equivalently $v_r(s)=v_r(t)=-1$, so $v_r(st)=-2\neq0=v_r(1)$.) Hence $P\cap Q=\varnothing$.

Finally $N_P=D_Q=\prod_{q\in Q}q$ is a product of distinct primes, hence squarefree with prime
support $Q$; symmetrically for $N_Q$. And $s\,D_P=(N_P/D_P)D_P=N_P=D_Q$.
\end{proof}

The disjointness in Theorem~\ref{thm:structure} means a solution is genuinely a partition of
$U\coloneqq P\cup Q$, and $D_PD_Q=\prod_{p\in U}p$. The valuation form $v_p(\sum 1/p_i)=-1$ underlying
disjointness is due to R.~Israel \cite{MO320838}, and appears also in Seva \cite{Seva323194}; the
cross--matching form ($N_P=D_Q$, $N_Q=D_P$) was obtained independently and concurrently by Bado
\cite[Thm.~A]{Bado2026} (Section~\ref{sec:prior}).

\section{The bridge to the arithmetic derivative}\label{sec:bridge}

\begin{lemma}\label{lem:der-cofactor}
For squarefree $m$, the arithmetic derivative equals the cofactor sum: $m'=\sum_{p\mid m}m/p=N_{\{p:\,p\mid m\}}$. Moreover $\gcd(m,m')=1$.
\end{lemma}

\begin{proof}
For $m=\prod_{p\mid m}p$ squarefree, $m'=m\sum_{p\mid m}1/p=\sum_{p\mid m}m/p$, which is exactly the
cofactor sum of the prime set of $m$. Coprimality is Lemma~\ref{thm:rigidity}.
\end{proof}

\begin{theorem}[Bridge]\label{thm:bridge}
The following are equivalent.
\begin{enumerate}[label=\textup{(\arabic*)},leftmargin=2.2em]
\item Erd\H{o}s \#307 has a solution: finite prime sets $P,Q$ with
$\bigl(\sum_{p\in P}1/p\bigr)\bigl(\sum_{q\in Q}1/q\bigr)=1$.
\item The arithmetic derivative has a two--cycle: integers $a\ne b$ with $a'=b$ and $b'=a$.
\item The equation $n''=n$ has a solution $n$ that is not a fixed point (i.e.\ $n'\ne n$, so $n$ is
not of the form $p^{p}$).
\end{enumerate}
\textup{(}The implication $(2)\Rightarrow(1)$ uses the Ufnarovski--\AA hlander squarefreeness
theorem for cycles; all other implications, and everything used in the barrier of
Section~\textup{\ref{sec:barrier}}, are unconditional; see Remark~\textup{\ref{rem:cited}}.\textup{)}
\end{theorem}

\begin{proof}
$(1)\Rightarrow(2)$. Let $a=D_P,b=D_Q$. By Lemma~\ref{lem:der-cofactor},
$a'=N_P$ and $b'=N_Q$, and by Theorem~\ref{thm:structure}, $N_P=D_Q=b$ and $N_Q=D_P=a$. Thus
$a'=b$, $b'=a$, and $a\ne b$ since $s\ne1$ (a prime--reciprocal sum is never $1$, being
$N/D$ in lowest terms with $D>1$).

$(2)\Rightarrow(1)$. Let $a'=b,b'=a$ with $a\ne b$. By the Ufnarovski--\AA hlander analysis of
two--cycles (cf.\ \cite{UfnarovskiAhlander2003,Kovic2012}), both $a$ and $b$ are squarefree; their
prime supports $P,Q$ are then disjoint since $\gcd(a,b)=\gcd(a,a')=1$. By
Lemma~\ref{lem:der-cofactor}, $a'=N_P$ and $b'=N_Q$; the
cycle gives $N_P=b=D_Q$ and $N_Q=a=D_P$, so
$\bigl(\sum_{p\in P}1/p\bigr)\bigl(\sum_{q\in Q}1/q\bigr)=\frac{N_P}{D_P}\frac{N_Q}{D_Q}
=\frac{D_Q}{D_P}\frac{D_P}{D_Q}=1$.

$(2)\Leftrightarrow(3)$. A two--cycle $a\ne b$ gives $a''=a$ with $a'=b\ne a$, so $a$ is a non--fixed
solution of $n''=n$. Conversely a solution of $n''=n$ with $n'\ne n$ yields the two--cycle
$(n,n')$. The fixed points of $n\mapsto n'$ are exactly the numbers $p^{p}$
\cite{Barbeau1961,UfnarovskiAhlander2003}.
\end{proof}

\begin{remark}[Status of the cited input]\label{rem:cited}
Part $(2)\Rightarrow(1)$ uses that a derivative two--cycle has squarefree, support--disjoint
members. Squarefreeness is the content of the Ufnarovski--\AA hlander analysis of cycles (a
non--squarefree $n$ has $p\mid\gcd(n,n')$ for some $p$, and along $n\mapsto n'$ such common factors
propagate, precluding return); disjointness then follows from $\gcd(a,a')=1$. We use it as a cited
theorem. Even \emph{without} it, the implication $(1)\Rightarrow(2)$ and the entire barrier below are
unconditional, because Rigidity already delivers squarefree, disjoint $P,Q$ from a \#307 solution.

\textup{Formalised.} \texttt{Erdos307.bridge\_forward} is $(1)\Rightarrow(2)$ and cites nothing;
the distinctness $a\ne b$ is \emph{derived} there, not assumed
\textup{(}\texttt{dprod\_ne\_dprod}\textup{)}. Support--disjointness is available separately as
\texttt{disjoint\_primeFactors\_of\_cycle}, also from Rigidity; it is not consumed by
\texttt{bridge\_forward}, which does not need it.
\texttt{bridge\_backward} is $(2)\Rightarrow(1)$ with the Ufnarovski--\AA hlander squarefreeness
appearing as an explicit hypothesis, and \texttt{two\_cycle\_iff\_double\_fixed} is
$(2)\Leftrightarrow(3)$ for an arbitrary self--map, on \emph{no} axioms at all. $0$ \texttt{sorry}
\textup{(\texttt{lean/Erdos307/Bridge.lean})}.
\end{remark}

\begin{corollary}\label{cor:ua}
If the Ufnarovski--\AA hlander conjecture holds in the weak form ``every cycle of $n\mapsto n'$ has
period one,'' then Erd\H{o}s \#307 has answer \textsc{no}. Conversely, a \textsc{yes} to \#307
exhibits a period--two cycle and refutes that statement.
\end{corollary}

\section{The barrier}\label{sec:barrier}

\begin{lemma}[Strict AM--GM]\label{lem:amgm}
If $s,t>0$ with $st=1$ and $s\ne1$, then $s+t>2$.
\end{lemma}
\begin{proof}
$s+t-2=s+1/s-2=(s-1)^2/s>0$ since $s\ne1$.
\end{proof}

\begin{lemma}[The number $59$, after Rosen \cite{MO320838}]\label{lem:59}
Any finite set $U$ of distinct primes with $\sum_{p\in U}1/p>2$ has $|U|\ge 59$. Moreover, among
$k$--element prime sets, $\sum 1/p$ is maximised by the $k$ smallest primes.
\end{lemma}
\begin{proof}
Replacing any prime of $U$ by a smaller prime not in $U$ increases $\sum1/p$, so the maximum over
$k$--element sets is at the first $k$ primes. The reciprocal sum of the first $58$ primes is
$1.99874\ldots<2$ and of the first $59$ primes is $2.00235\ldots>2$, so a sum exceeding $2$ needs at
least $59$ primes.
\end{proof}

\begin{theorem}[Barrier]\label{thm:barrier}
For any solution $(P,Q)$ of \#307, with $U=P\cup Q$:
\begin{enumerate}[label=\textup{(\alph*)},leftmargin=2.2em]
\item $|U|\ge 59$ and $\displaystyle\prod_{p\in U}p \ge \Pi_{59}\coloneqq\prod_{i=1}^{59}p_i
\approx 8.77\times10^{112}$;
\item $\displaystyle\min(\textstyle\prod P,\prod Q)\ \ge\ \sqrt{\Pi_{59}/T_{59}} \ =\
2.0929\ldots\times10^{56}$, where $T_{59}=\sum_{i=1}^{59}1/p_i=2.00235\ldots$. In particular both
$\prod P$ and $\prod Q$ exceed $2.09\times10^{56}$.
\end{enumerate}
Consequently every $n''=n$ solution has both members $>2.09\times10^{56}$. The strongest valid
unconditional bound previously available is the computer search $x>10^4$ of \cite{Kovic2012}, so this
is a gain of $52$ orders of magnitude; against the conditional Proposition~17 of that paper
($3.23\times10^9$, when the smaller member is odd) it is $47$. The unconditional $r+s\ge34$ and
$\max\{m,n\}\ge1.92\times10^{21}$ claimed in Proposition~18 of \cite{Kovic2012} are not established
by their proof (Section~\ref{sec:prior}); the bound above does not rely on them and in fact implies
the first.
\end{theorem}

\begin{proof}
(a) By $st=1$, $s\ne1$ and Lemma~\ref{lem:amgm}, $\sum_{p\in U}1/p=s+t>2$, so $|U|\ge 59$ by
Lemma~\ref{lem:59}; since $P,Q$ are disjoint (Theorem~\ref{thm:structure}),
$\prod_{p\in U}p=D_PD_Q$, and a product of $\ge 59$ distinct primes is at least the product of the
$59$ smallest, $\Pi_{59}$.

(b) Take $D_P\le D_Q$ to be the minimal side. By Theorem~\ref{thm:structure}, $D_Q=s\,D_P$ with
$s>1$, hence
\[
s\,D_P^2 \;=\; D_P\,D_Q \;=\; \prod_{p\in U}p \;=\!:\, R,
\qquad\text{so}\qquad D_P^2=\frac{R}{s}.
\]
Since $s\le s+t=T(U)\coloneqq\sum_{p\in U}1/p$ (as $t>0$),
\[
D_P^2=\frac{R}{s}\ \ge\ \frac{R}{T(U)}.
\]
It remains to bound $R/T(U)$ below over all admissible $U$. Among $k$--element prime sets the
quotient $\bigl(\prod U\bigr)\big/\bigl(\sum_{p\in U}1/p\bigr)$ is minimised by the $k$ smallest
primes (smallest numerator and, by Lemma~\ref{lem:59}, largest denominator simultaneously), and
passing from $59$ to more primes only increases it (the product grows multiplicatively while the
reciprocal sum grows by a single small term). Hence the global minimum over admissible $U$ is at
$U=\{p_1,\dots,p_{59}\}$, giving $R/T(U)\ge \Pi_{59}/T_{59}$. Therefore
\[
D_P^2 \ \ge\ \frac{\Pi_{59}}{T_{59}} \;=\; \frac{\Pi_{59}^2}{N_{59}}\;=\;4.3806\ldots\times10^{112},
\qquad D_P\ \ge\ 2.0929\ldots\times10^{56},
\]
where $N_{59}=T_{59}\,\Pi_{59}$ is the integer cofactor sum of the first $59$ primes.
\end{proof}

\begin{remark}[Why the \emph{minimum} is bounded, not only the maximum]
The product identity $D_PD_Q\ge\Pi_{59}$ alone gives only $\max(\prod P,\prod Q)\ge
\sqrt{\Pi_{59}}\approx2.96\times10^{56}$. The stronger statement~(b), a lower bound on the
\emph{smaller} side, genuinely uses rigidity: it is the identity $D_Q=s\,D_P$ (equivalently
$D_P^2=R/s$) together with $s\le T(U)$ that converts the product bound into a bound on $D_P$. We
verified numerically that $\min_U R(U)/T(U)$ is attained exactly at the first $59$ primes (the
minimising configuration is unique), so the constant $2.09\times10^{56}$ is sharp for this method.
\end{remark}

\begin{proposition}[Closing the $59$--prime level]\label{prop:close59}
No solution of \#307 has $|P\cup Q|=59$. Hence $|P\cup Q|\ge60$,
$\prod_{p\in U}p\ \ge\ \Pi_{60}>2.46\times10^{115}$, and
\[ \min\bigl(\textstyle\prod P,\prod Q\bigr)\ \ge\ \sqrt{\Pi_{60}/T_{60}}\ =\
3.5053\ldots\times10^{57}. \]
\end{proposition}

\begin{proof}[Proof (finite verification)]
Let $U=P\cup Q$ with $|U|=59$; as in Theorem~\ref{thm:barrier}(a), $T(U)>2$. The candidate list is
finite and complete: \textup{(i)} every prime $p\le167$ lies in $U$; omitting $p$ caps $T(U)$ at
$T_{60}-1/p<2$, since the $59$ largest prime reciprocals avoiding $p$ are those of the first $60$
primes less $p$, and $1/p>T_{60}-2=0.00590\ldots$ exactly for $p\le167$; \textup{(ii)} every element
$p'\ge281$ satisfies $T_{58}+1/p'>2$ (the other $58$ reciprocals sum to at most $T_{58}$), so
$p'\le787$. Thus $U$ consists of the $39$ primes up to $167$ plus twenty primes from
$(167,787]$, and exactly $49{,}961$ such sets have $T(U)>2$. For each, a solution supported on $U$
forces both $N'+2N=(a+b)^2$ and $N'-2N=(a-b)^2$ to be perfect squares, where $N=\prod U$ and
$N'=a^2+b^2$ (rigidity; the $k=0$ case of Proposition~\ref{prop:pairform}); one integer test that
simultaneously excludes all $2^{58}$ splits of $U$. Exact--integer verification shows $N'+2N$ is a
non--square for \emph{every} one of the $49{,}961$ candidates, by two independent implementations
with different enumeration orders and prunings (\texttt{code/close59.py} and
\texttt{hunt/close59.rs} in the repository). Hence no $59$--element support closes, and
$|U|\ge60$. The bound on the minimal side repeats the proof of Theorem~\ref{thm:barrier}(b) with the
minimum of $R/T$ now over $|U|\ge60$, attained at the first $60$ primes:
$D_P^2\ge\Pi_{60}/T_{60}=\Pi_{60}^2/N_{60}$, where $N_{60}$ is the cofactor sum of the first $60$
primes.
\end{proof}

\begin{corollary}[The coprime variant inherits the closed level]\label{cor:coprime60}
Consider the weakened variant in which the elements of $P$ and $Q$ are merely pairwise coprime
integers $\ge2$ \textup{(}the open form of the variant, $1\notin P\cup Q$\textup{)}. Any solution
has $|P\cup Q|\ge 60$ elements, and $\min(\prod P,\prod Q)>3.50\times10^{57}$.
\end{corollary}
\begin{proof}
Rigidity holds verbatim for pairwise--coprime families (the cofactor--sum argument uses only
coprimality), so again $T(U)=s+t>2$. The least prime factors $\ell(e)$ of the elements are distinct,
and $\sum_e 1/e\le\sum_e1/\ell(e)$. Suppose $|U|=59$ and some element $e$ is not prime, so
$e\ge\ell(e)^2$. \textup{(}Formalised: \texttt{Erdos307.coprime\_loss} is the per--element bound,
\texttt{loss\_antitone} its monotonicity, and \texttt{loss\_at\_277} the numeric consequence
\textup{(}\texttt{lean/Erdos307/Coprime60.lean}\textup{)}; the slack itself is the level--59
census.\textup{)} If $\ell(e)\le277$, the loss is
$1/\ell-1/e\ge(\ell-1)/\ell^2\ge276/277^2=0.00359\ldots$, exceeding the maximal available slack
$T_{59}-2=0.00235\ldots$; if $\ell(e)\ge281$, then $e\ge281^2$ while the other $58$ reciprocals sum
to at most $T_{58}$, giving $T(U)\le T_{58}+281^{-2}=1.99875\ldots<2$. Either way $T(U)<2$, a
contradiction: at $59$ elements every element is prime, and Proposition~\ref{prop:close59} applies
verbatim. The product bound transfers because $\prod_e e\ge\prod_e\ell(e)$ and
$T(U)\le\sum_e1/\ell(e)$, so the extremal ratio bound at the first $60$ primes still applies.
\end{proof}

\begin{remark}[Verification tier, and why the ladder stops here]\label{rem:tier60}
Theorem~\ref{thm:barrier} is machine--checked in Lean~4, and so is
Proposition~\ref{prop:close59}: the theorem \texttt{erdos307\_sixty} formalises the entire
argument; the forced--primes and element bounds, the Pythagorean plus--square, and a pruned
depth--first search over the pool whose \emph{completeness is proved, not assumed}
(\texttt{dfs\_sound}), the search itself running under \texttt{native\_decide}; on top of the
two independent unverified programs. The level cannot be climbed further by finite means: already at $|U|=60$ the family
$U=\{p_1,\dots,p_{59}\}\cup\{q\}$ qualifies for \emph{every} prime $q>277$, and on it the two square
conditions read $x^2=Aq+\Pi_{59}$, $y^2=Bq+\Pi_{59}$ with $A,B=N_{59}\pm2\Pi_{59}$, i.e.\ the
Pell--type system $Bx^2-Ay^2=-4\Pi_{59}^2$ with $113$--digit coefficients. We verified that no
accessible modulus obstructs it (mod $32$, and all odd $\ell<4000$: the conditions pin
$q\equiv5\pmod 8$ and thin each $\ell\mid\Pi_{59}$ to about half of its residue classes and each
$\ell\nmid\Pi_{59}$ to about a quarter, but never to zero), consistent with the local--solubility
theme of Section~\ref{sec:anatomy}. The obstruction, however, lives at the primes \emph{of} $A$
and $B$, and it is detectable \emph{without factoring}: Proposition~\ref{prop:tailkill} below
closes this family for \emph{every} $q$ by a single Jacobi--symbol computation. (An exact sweep to
$q\le10^{12}$, \texttt{hunt/tail\_sweep.rs}, $1.25\times10^{11}$ candidates, none with
$Aq+\Pi_{59}$ even a single square, which corroborates it independently.) The tail family also exposes the self--similarity of the ladder:
writing $a=qm$, so that $m\,b=\Pi_{59}$ partitions the first $59$ primes, its two closing
conditions collapse to the exact equation
\[ D(m)\,D(b)+m^2\;=\;\Pi_{59} \]
together with the divisibility $m\mid D(b)$ (then $q=D(b)/m$, required prime); verified
exhaustively on a $12$--prime toy universe. Equivalently, the deficit $1-\sigma(m)\sigma(b)$ must
equal $m^2/\Pi_{59}$ \emph{exactly}: the circularity $\Delta_0=M_0N(1-s_{M_0}s_N)$ of
Section~\ref{sec:frame} in its purest form, with the deviation forced to be the perfect square
$m^2$. Closing even this one family is thus a $2^{58}$--fold exact--coincidence search of the same
type as \#307 itself (expected count ${\sim}2^{58}/\Pi_{59}\approx10^{-95}$): each rung of the
ladder past $60$ contains the whole problem. The one case with \emph{no} free large prime (a
$60$--set drawn entirely from small primes) is instead a finite search, and empty in the
minimal range: a direct both--squares enumeration of all $60$--prime sets with $T>2$ supported in
the first $70$ primes ($4{,}011{,}289$ sets; \texttt{hunt/close60.rs}) finds none passing even the
plus--square, so no level--$60$ solution has all its primes $\le349$. (The first $68$ primes give
$1{,}258{,}448$ sets, also empty, and reproduce independently.) By
Corollary~13.5 of the full version this verdict is decisive rather than provisional: at these masses a
support passing the plus-- and minus--squares, with the parity of Proposition~13.1(2) of the full version,
\emph{is} a two--cycle, so a survivor would have been a solution and not a candidate. The
enumeration extends with compute.
\end{remark}

\begin{proposition}[The canonical tail family is empty: a reciprocity kill]\label{prop:tailkill}
No solution of \#307 has support $U=\{p_1,\dots,p_{59}\}\cup\{q\}$, for any prime $q$.
\end{proposition}
\begin{proof}
A solution supported on $U$ forces $x^2=Aq+\Pi_{59}$ with $x=a+b$ and $A=N_{59}+2\Pi_{59}$
(Remark~\ref{rem:tier60}). Every prime $\ell\mid A$ exceeds $277$: for $\ell\le277$ one has
$\ell\mid\Pi_{59}$, so $\ell\mid A$ would force $\ell\mid N_{59}$, contradicting
$\gcd(N_{59},\Pi_{59})=1$ (rigidity). Hence $\ell\nmid\Pi_{59}$, and reduction mod $\ell$ gives
$x^2\equiv\Pi_{59}\pmod\ell$, requiring the Legendre symbol $(\Pi_{59}\mid\ell)\neq-1$. But the
Jacobi symbol satisfies
\[ \left(\tfrac{\Pi_{59}}{A}\right)\;=\;-1 \]
(a finite gcd--speed computation, script \texttt{code/tailkill.py}; no factorisation of the
$114$--digit $A$ is needed), so some $\ell\mid A$ of odd multiplicity has $(\Pi_{59}\mid\ell)=-1$.
The family is empty, for every $q$.

The companion symbol at $B=N_{59}-2\Pi_{59}$ equals $-1$ as well; necessarily: for any
admissible base $S$ (writing $D=D_S$, $N=N_S$, $A_S,B_S=N\pm2D$) one has $A_S\equiv B_S\pmod 8$
(as $4D\equiv8\bmod{16}$) and $A_S\equiv B_S\equiv N\pmod p$ for every odd $p\mid D$, while the
quadratic--reciprocity sign difference carries the exponent $\tfrac{p-1}{2}\cdot\tfrac{A_S-B_S}{2}$
with $\tfrac{A_S-B_S}{2}=2D$ even; hence $(D\mid A_S)=(D\mid B_S)$ always; one binary
certificate per tail family. Running it over all $49{,}961$ admissible bases of
Proposition~\ref{prop:close59} proves $31{,}219$ of the one--new--prime families at level $60$
empty ($62.5\%$, curiously near $5/8$); the remaining $18{,}742$ are Jacobi--inconclusive: their
$-1$--Legendre primes, if any, occur to an even count, visible only through a factorisation of
$A_S$. A factoring attack on those (trial division and Brent--Pollard $\rho$ over a Montgomery
core, \texttt{hunt/survivor\_kill.rs}; the found factors are kept as a pairwise--coprime basis, so
a kill certificate (an odd--multiplicity basis piece $P$ with $(D_S\mid P)=-1$) needs no
primality testing at all), pushed to $\rho$--budget $10^6$ with a Pollard $p-1$ stage, closes a
further $15{,}264$: in total $46{,}483$ of the $49{,}961$ one--new--prime families at level $60$
are proven empty ($93.0\%$). The $3{,}478$ still open resist all of this for a structural reason,
not a want of effort: each is Baillie--PSW \emph{composite} yet has no prime factor below ${\sim}
10^{12}$, so its $A_S$ (or $B_S$) is a \emph{balanced} semiprime of two ${\sim}57$--digit primes,
splitting one is factoring at cryptographic scale. Two cheap tests confirm no shortcut survives:
a product--tree batch--GCD over all $37{,}484$ values $A_S,B_S$ finds only shared primes below
$10^9$ (already found by $\rho$, so \emph{no} free kills), and Baillie--PSW leaves exactly
$34$ families in which $A_S$ and $B_S$ are both \emph{probable} primes; these are immune
\emph{conditionally on that primality} (a prime $A_S$ with $(D_S\mid A_S)=+1$ has no inert factor,
hence no congruence obstruction for any $q$). The conditionality is not pedantry and not
removable by more of the same computation: Baillie--PSW is rigorous only in its \emph{composite}
verdicts, so the $46{,}483$ kills are unconditional while the $34$ immunities are not. Upgrading them
would require primality \emph{certificates} for thirty-four $114$--digit integers (ECPP or a
Pocklington factorisation of $A_S-1$). The elementary route does not reach: immunity needs
\emph{both} $A_S$ and $B_S$ prime, and partial factorisation of $n-1$ (trial division to
$3\times10^6$ together with the fully split part) leaves $\log F/\log n$ below $1/3$ on at least one
side of every one of the $34$, so neither Pocklington ($F>\sqrt n$) nor
Brillhart--Lehmer--Selfridge ($F>n^{1/3}$) applies to any family: $0$ of $34$ are certifiable this
way (\texttt{code/immune\_certify.py}). This is now settled. Running the Adleman--Pomerance--Rumely
--Cohen--Lenstra test on all $68$ integers proves every one of them prime in $8$ seconds of
wall clock, so the $34$ immunities upgrade from \textup{[C]} to \textup{[P]} and are
\emph{unconditional} (\texttt{code/immune\_prove.gp}, PARI/GP $2.17.4$, whose \texttt{isprime} is a
proof and not a probable--prime test). The run reproduces the whole chain independently: $49{,}961$
admissible bases, $31{,}219$ Jacobi kills, and, among the $81$ bases whose $A_S$ and $B_S$ are both
prime, exactly the $34$ with $(D_S\mid A_S)=+1$; the other $47$ carry $(D_S\mid A_S)=-1$ and were
already killed. The primality is moreover backed by
\emph{independently checkable} evidence rather than a verdict: \texttt{code/immune\_certify.gp}
emits Atkin--Morain \textup{ECPP} certificates for all $68$ integers, archived at
\texttt{certs/immune\_ecpp.txt} ($512$\,KB) so a third
party can check them without trusting the prover. That file is a \texttt{primecertexport} dump and
is not machine--readable: PARI's \texttt{read} parses its header as a polynomial, so
\texttt{primecertisvalid} cannot consume it, and earlier versions of this paper and of the
repository documented that command in error. The working check is
\texttt{code/immune\_cert\_verify.gp}, which extracts the $68$ asserted subjects and re--proves each
by \textup{ECPP} in $1.3$ seconds: $68/68$ prime, $0$ failures. The structural half of the campaign is now
formal: the identity $(D_S\mid A_S)=(D_S\mid B_S)$, which is what makes one binary test decide a
family rather than two, is proved in Lean~4 as \texttt{Erdos307.jacobiSym\_A\_eq\_B}
\textup{(\texttt{lean/Erdos307/Campaign.lean})}, on the three standard axioms and with no
\texttt{native\_decide}; it follows from periodicity, $A_S-B_S=4D_S$ being exactly the period of
$J(a\mid\cdot)$. The primality itself is not formalised and cannot presently be: mathlib's only
route to a large prime is \texttt{lucas\_primality}, which needs a complete factorisation of
$n-1$, and that is precisely what the $0$ of $34$ above rules out. Note the two counts answer different questions and both are correct: $81$ is the
number of prime pairs, $34$ the number of \emph{immune} families among them. A complete kill/immune split of the residual is in any case beyond feasible
computation, and the honest state of the single--tail sector is \emph{$46{,}483$ empty
\textup{[P]}, $34$ immune and now unconditionally \textup{[P]} by the \textup{APRCL} run above,
$3{,}478$ factoring--hard}.
\end{proof}

\begin{proposition}[Sector decomposition of level $60$]\label{prop:sectors}
Let $U$ be a $60$--element prime support with $T(U)>2$. Then exactly one of the following holds:
\textup{(i)} $T(U)-2\ge 1/\max U$, and $U=S\cup\{q\}$ for an admissible base $S$
\textup{(}one of the $49{,}961$ of Proposition~\ref{prop:close59}\textup{)}; the
\emph{single--tail sector}, or \textup{(ii)} $T(U)-2<1/\max U$, and, writing $p>q$ for the two
largest elements and $W$ for the rest, $T(W\cup\{p\})<2$ and $T(W\cup\{q\})<2$; the
\emph{pair sector}, whose members escape every single--tail family.
\end{proposition}
\begin{proof}
$U$ lies in the tail family of $U\setminus\{r\}$ iff $T(U)-1/r\ge2$, which is easiest for
$r=\max U$; (i) is that case ($T=2$ being impossible by rigidity, the base is admissible). In
case (ii) no removal reaches $2$, and a fortiori the two stated $59$--subsets stay below $2$.
\end{proof}

\begin{remark}[The campaign; completeness of the certificate; the arity barrier]\label{rem:campaign}
The reciprocity certificate captures \emph{every} $q$--independent congruence obstruction of a
tail family: at a prime $\ell\mid A_SB_S$ the system degenerates to the rank--one congruences
$x^2\equiv D_S$ or $y^2\equiv D_S\pmod\ell$, while at $\ell\nmid 2A_SB_S$ the conditions define
conics in $(x,q)$ or $(y,q)$ with $\sim\ell$ points, deleting at most half of the residue classes
of $q$ and never all (the $2$--adic layer pins $q\equiv5\bmod8$ without emptiness); by CRT and
Dirichlet, admissible $q$ survive any finite set of such conditions. Hence a tail family admits a
congruence kill iff $(D_S\mid\ell)=-1$ for some $\ell\mid A_SB_S$, which the Jacobi symbol
detects for $31{,}219$ bases, and partial factorisation (\texttt{hunt/survivor\_kill.rs}) for a
further $15{,}264$: in total $46{,}483$ of the $49{,}961$ single--tail families at level $60$ are
proven empty ($93.0\%$), $3{,}478$ remaining open (each a balanced ${\sim}57$--digit semiprime, factoring--hard). The weapon is intrinsically \emph{arity--one}, and provably fails on the pair sector. With
$s=p+q$, $t=pq$, a pair--sector member requires the two simultaneous squares $x^2=A_Wt+D_Ws$ and
$y^2=B_Wt+D_Ws$ (the would--be third condition $s^2-4t=\square$ is automatic, being $(p-q)^2$),
together with the reciprocity--linked constraints $(D_Wp\mid q)=(D_Wq\mid p)=1$ (mod $q$, a
solution with $q\mid a$ has $N_U\equiv D_Wp$, forcing $D_Wp$ to be a square). At a prime
$\ell\mid A_W$ the condition $x^2\equiv D_Ws\pmod\ell$ constrains only $s$, which the free second
prime absorbs; the pair system is in fact locally soluble at \emph{every} modulus (a theorem,
Proposition~\ref{prop:pairlocal} below, corroborated by a direct scan over all $\ell\le600$) so,
unlike the single--tail family, the pair sector admits \emph{no} congruence obstruction at all. It
is a genuine two--parameter surface per base, $\#307$ in miniature (Remark~\ref{rem:oneside});
reciprocity cannot touch it, and only a direct both--squares search or a descent argument can.
\end{remark}

\begin{proposition}[The pair sector admits no congruence kill]\label{prop:pairlocal}
Let $W$ be any pair--sector base and $A=A_W$, $B=B_W$, $D=D_W$, so that a pair--sector solution
requires $x^2=A\,\alpha\beta+D(\alpha+\beta)$ and $y^2=B\,\alpha\beta+D(\alpha+\beta)$ with
$\alpha,\beta$ the two new primes. For every prime power $\ell^j$ this system has a solution in
units $\alpha,\beta$ and integers $x,y$ at a nonsingular point. Consequently, by CRT and Dirichlet,
no finite system of congruence conditions empties any pair--sector family: the arity barrier of
Remark~\ref{rem:campaign} is a theorem, and the certificate programme provably ends at arity one.

\textup{Formalised, in part.} \texttt{Erdos307.no\_pinning\_isUnit} is the step that makes the
tail--kill mechanism unavailable, \texttt{collapse\_identity} is regime~\textup{(2)}'s
observation that the two square conditions are one, \texttt{regime\_one\_first} and
\texttt{regime\_one\_second} are the explicit unit of regime~\textup{(1)}, and
\texttt{targets\_agree\_mod\_eight} is the mod--$8$ collapse of regime~\textup{(3)}. The lift from
$\ell$ to $\ell^j$ that regime~\textup{(1)} needs is \texttt{square\_lifts\_to\_prime\_powers},
proved from \texttt{hensel\_all\_powers}: if $\ell$ is an odd prime not dividing $c$ and $c$ is a
square modulo $\ell$, it is a square modulo every power of $\ell$. Weil's bound for the point count
in regime~\textup{(2)} is the one remaining input, and it is \emph{cited}, not pending: the quantity
$\delta=u^2-4m$ is a quartic in $y$, so the count is Hasse--Weil for $v^2=\delta(y)$, and Mathlib has
elliptic--curve infrastructure but no point--count bound. Formalising Hasse's theorem is a
contribution to Mathlib rather than to this problem, so the star here records a citation. $0$
\texttt{sorry} \textup{(\texttt{lean/Erdos307/PairLocal.lean})}.
\end{proposition}
\begin{proof}
First, no $\ell$ pins either target to a constant: that needs $\ell\mid A$ (or $B$) \emph{and}
$\ell\mid D$, whence $\ell\mid N$, contradicting $\gcd(N,D)=1$ (rigidity); so the tail--kill
mechanism is unavailable, and solubility is decided in three regimes.

\emph{(1) $\ell$ odd, $\ell\mid D$.} This covers every odd $\ell\le103$: at level $60$, omitting a
prime $p$ caps $T(U)$ at $T_{61}-1/p$, and $1/p>T_{61}-2=0.00944\ldots$ exactly for $p\le103$, so
all such primes lie in $U$, hence in $W$ (the two adjoined primes are the largest). Take $\beta=1$,
$x=1$, $\alpha=(1-D)/(A+D)\bmod\ell^j$, a unit since $A+D\equiv N\not\equiv0\pmod\ell$: the first
equation holds by construction, and the second target $(B+D)\alpha+D\equiv N\cdot N^{-1}=1\pmod\ell$
is a unit congruent to a square, so $y$ exists modulo $\ell^j$.

\emph{(2) $\ell$ odd, $\ell\nmid2D$} (in particular every $\ell\ge107$ outside $W$). The two
conditions are secretly \emph{one}: setting $m=(x^2-y^2)/4D$ and $u=(x^2-Am)/D$, the relation
$A-B=4D$ gives the polynomial identity $Bm+Du=y^2$. So any $(x,y)$ with $xy\ne0$, $x^2\ne y^2$ and
$\delta:=u^2-4m$ a square delivers $\alpha,\beta$ as the roots of $z^2-uz+m$ (nonzero, as
$\alpha\beta=m\ne0$). For each admissible $x$, $\delta$ is a quartic in $y$ and not a square in
$\overline{\mathbb F}_\ell[y]$ (a factorisation $(Q-R)(Q+R)=16D^3(x^2-y^2)$ would force
$\ell\mid4D$), so Weil's bound gives $\tfrac12\ell^2+O(\ell^{3/2})$ admissible pairs; positive
well below $\ell=107$. The point is nonsingular ($xy\ne0$; the Jacobian contains
$\operatorname{diag}(2x,2y)$), so it lifts to every $\ell^j$.

\emph{(3) $\ell=2$.} Both targets are odd and agree modulo $8$ (as $\alpha\beta$ is odd,
$D(\alpha+\beta)\equiv0\bmod4$, and $(A-B)\alpha\beta=4D\alpha\beta\equiv0\bmod8$), and a finite
table over $(A\bmod8,\,D\bmod8)$ shows that for every base some achievable pair
$(\alpha\beta,\alpha+\beta)\bmod8$ makes the common target $\equiv1\pmod8$; a $2$--adic square,
soluble modulo $2^j$ for every $j$.

All three regimes, the collapse identity, the threshold $1/103>T_{61}-2>1/107$, and the mod--$32$
layer were verified computationally on three bases and every $\ell\le600$
(\texttt{code/pairlocal.py}; the measured density of admissible $(x,y)$ in regime~(2) is
$0.497$--$0.503$ against the predicted $\tfrac12$).
\end{proof}

\begin{remark}[Anatomy of the certificate; the empirical $5/8$ law]\label{rem:fiveeighths}
Writing $D=2D_1$, the reciprocity signs in the certificate cancel
\textup{(}$\tfrac{A-1}2+\tfrac{N-1}2=N-1+2D_1$ is even\textup{)}, leaving the clean product
\[ \left(\tfrac{D}{A}\right)\;=\;\left(\tfrac{2}{A}\right)\left(\tfrac{D_1}{N}\right),\qquad
\left(\tfrac{2}{A}\right)=+1\iff N\equiv3,5\ (\mathrm{mod}\ 8), \]
and one layer down $(N\mid D_1)=\prod_{p\mid D_1}\bigl(\tfrac{D/p}{p}\bigr)$, since
$N\equiv D/p\pmod p$: the cofactor structure of rigidity recurs inside its own certificate. This product evaluates in closed form: with $D/p=2D_1/p$ and quadratic
reciprocity,
\[ (N\mid D_1)=(2\mid D_1)\!\!\prod_{\{p,q\}\subset D_1}\!\!(-1)^{\frac{p-1}2\frac{q-1}2}
=(-1)^{\,s+\binom{t}{2}},\quad s=\#\{p\mid D_1:p\equiv3,5\ (8)\},\ \ t=\#\{p\mid D_1:p\equiv3\ (4)\}, \]
a R\'edei genus character of $D_1$, \emph{independent of $N$}. This is forced, not merely observed:
$(N\mid D_1)=\prod_{p\mid D_1}(N\mid p)$, and cofactor survival gives $N\equiv D/p=2D_1/p\pmod p$, so
$(N\mid p)=(2\mid p)\,(D_1/p\mid p)$. The first factors multiply to $(-1)^{s}$ since $(2\mid p)=-1$
exactly for $p\equiv3,5\pmod 8$; and, each unordered pair $\{p,q\}$ occurring once in each of the two
cofactors,
\[ \prod_{p\mid D_1}\Bigl(\tfrac{D_1/p}{p}\Bigr)=\prod_{\{p,q\}\subset D_1}\Bigl(\tfrac qp\Bigr)
\Bigl(\tfrac pq\Bigr)=\prod_{\{p,q\}\subset D_1}(-1)^{\frac{p-1}2\frac{q-1}2}=(-1)^{\binom t2}, \]
by quadratic reciprocity, since $\tfrac{p-1}2$ is odd exactly for $p\equiv3\pmod4$. (Checked
independently, prime by prime, on $20{,}000$ bases: $0$ violations.) The inner symbol is thus not a coin but a genus invariant; its empirical fairness
($49.9\%$, $50.0\%$ over the $49{,}961$ bases) is that character's equidistribution; and
$(D_1\mid N)$ follows by the reciprocity sign $(-1)^{\frac{D_1-1}2\frac{N-1}2}$, hence depends on $N$
only through $N\bmod4$. In fact the entire certificate closes into a finite functional. Let
$v(S)=(n_1,n_3,n_5,n_7)\bmod4$ count the odd primes of $S$ in the residue classes $1,3,5,7$ modulo
$8$. Every factor above is a function of $v$ alone: $s\equiv n_3+n_5$, $t=n_3+n_7$ (with
$\binom t2$ needing only $t\bmod4$), $D_1\equiv3^{n_3}5^{n_5}7^{n_7}\pmod8$,
$S(D_1)\equiv(n_1+n_5)-(n_3+n_7)\pmod4$, and (by the mod--$8$ law of
Proposition~13.14 of the full version, $D_1'\equiv D_1S(D_1)\pmod8$) also
$N=D_1+2D_1'\equiv D_1+2\bigl(D_1S(D_1)\bmod4\bigr)\pmod8$. Hence
\[ (D\mid A)\;=\;K\bigl(v(S)\bigr): \]
\emph{the kill is decided by sixteen residue counts modulo $4$.} The realized vectors are exactly
the coset $n_1+n_3+n_5+n_7\equiv58\equiv2\pmod4$ (all $64$ of its classes occur), $K$ is constant
on each, and the functional reproduces the direct $113$--digit Jacobi computation on every one of
the $49{,}961$ bases (\texttt{code/kill\_classvector.py}). The five--eighths law is then an exact
finite statement, and its mechanism is the \emph{reciprocity switch}: for $N\equiv1\pmod4$ the
reciprocity sign is inert and the kill reduces to the genus coin; exact conditional rate
$12{,}699/24{,}975=0.5085$; while for $N\equiv3\pmod4$ the sign couples the parity of $t$ into
the product and the rate jumps to $18{,}520/24{,}986=0.7412$ (by class modulo $8$:
$0.509,\,0.741,\,0.508,\,0.741$ at $N\equiv1,3,5,7$; the $3/4$ classes are $N\equiv3\pmod4$, not
the $(2\mid A)$ classes). The blend $\tfrac12\cdot\tfrac{24975}{49961}+\tfrac34\cdot\tfrac{24986}{49961}
=0.62503$ against the exact $31{,}219/49{,}961=0.62487$ is the whole of the ``coincidence''. The
correlation is an ensemble property, a random--subset proxy gives ${\approx}0.45/0.55$; and
the functional yields no second certificate: it recomputes the same symbol.
\end{remark}

\begin{remark}[The kill--program is one--sided]\label{rem:oneside}
Every kill is a proof that no solution hides in that family, but if \#307 has a solution, the
solution's own family is unkillable, so no accumulation of kills can overshoot the truth: the
program's asymptote \emph{is} the existence question. Together with the self--similarity of
Remark~\ref{rem:tier60} (each rung past $60$ contains the whole problem), this locates the
campaign precisely: it corners where a solution could live, one certified region at a time, and
can terminate only if \#307 has answer \textsc{no}, which the heuristics disfavour.
\end{remark}

\begin{remark}[The reciprocity kills are the local factors of the expected count]\label{rem:killcount}
The divergent naive count of Section~\ref{sec:anatomy} and the reciprocity sieve of
Proposition~\ref{prop:tailkill} are the \emph{same} computation. Decompose the expected number of
solutions over tail families $S\cup\{q\}$: the summand for a base $S$ is
$\kappa_S\,\bigl(D_S\sqrt{\mathrm{mass}^2-4}\bigr)^{-1}\sum_q q^{-1}$, where
$\kappa_S=\prod_\ell(\text{local solubility density at }\ell)$ is exactly the reciprocity data.
A reciprocity--killed base carries an inert $\ell\mid A_S$ with $(D_S\mid\ell)=-1$, so
$N_S+2D_S\equiv D_S\pmod\ell$ is a non--residue for \emph{every} $q$, forcing $\kappa_S=0$; the
family drops from the count entirely. An immune base ($A_S,B_S$ both prime, both residues) has no
inert prime, so $\kappa_S>0$ and its summand is a divergent $\sum_q q^{-1}$ (verified on an explicit
$114$--digit immune base: local density $1$ at every tested prime, partial sums growing without
bound). Hence the campaign of Remark~\ref{rem:campaign} is a Selberg--sieve evaluation of the
singular series: the kills remove ${\sim}93\%$ of the count's mass, and the residual
$\log$--divergence (the entire ``weakly favours \textsc{yes}'' signal) is carried by the $34$
conditionally immune families together with the pair sector. (Conditionally, because their immunity
rests on the Baillie--PSW primality of $A_S,B_S$; the pair sector's contribution, by contrast, is
unconditional, being a theorem, Proposition~\ref{prop:pairlocal}.) This does not touch the reliability of the
count (which ignores the $N(r)\le1$ rigidity), but it pins down \emph{where} the heuristic's
\textsc{yes} lives.
\end{remark}

\begin{proposition}[Local completeness of the reciprocity sieve]\label{prop:localcomplete}
For an admissible base $S$ and an odd prime $\ell$, the admissible tail residues
$Q_\ell=\{q\bmod\ell:(A_Sq+D_S\mid\ell)\ne-1\text{ and }(B_Sq+D_S\mid\ell)\ne-1\}$ fall into exactly three
regimes. \textup{(i)} $\ell\mid D_S$, which covers every $\ell\le167$ uniformly over the $49{,}961$
admissible bases, since $1/p\ge(T_{59}-2)+\tfrac1{281}$ forces every prime $p\le167$ into every base:
then $A_S\equiv B_S\equiv N_S\not\equiv0\pmod\ell$, the two conditions collapse to the single condition
$(N_Sq\mid\ell)\ne-1$, and $|Q_\ell|=(\ell+1)/2\ni0$. \textup{(ii)} $\ell\mid A_SB_S$: one condition is
the constant $(D_S\mid\ell)$, so $Q_\ell=\emptyset$ iff $(D_S\mid\ell)=-1$ (exactly the certificate of
Proposition~\ref{prop:tailkill}) and otherwise $|Q_\ell|\ge(\ell-1)/2$. \textup{(iii)}
$\ell\nmid A_SB_SD_S$ (hence $\ell\ge173$): the complete character sums give
$|Q_\ell|\ge(\ell-3)/4-2\ge40$ (the linear sums vanish, the quadratic sum is $-\chi(A_SB_S)$, boundary
terms $\le2$). \textup{(}Regimes~\textup{(i)} and~\textup{(ii)} are formalised:
\texttt{Erdos307.regime\_one\_collapse}, \texttt{regime\_one\_same\_argument} and
\texttt{regime\_two\_constant} \textup{(}\texttt{lean/Erdos307/LocalComplete.lean}\textup{)}.
Regime~\textup{(iii)}'s counting step is \texttt{regime\_three\_count}, which takes the two complete
character--sum evaluations as hypotheses. Those evaluations, $\sum_q\chi(aq+b)=0$ for $a\ne0$ and
$\sum_q\chi((aq+b)(cq+d))=-\chi(ac)$ for $ad\ne bc$, are classical and are \emph{cited}, not pending:
Mathlib does not carry them, and supplying them is a contribution to Mathlib rather than to this
problem. The star on this proposition records a citation, not unfinished work.\textup{)}
Consequently $Q_\ell\ne\emptyset$ forces $|Q_\ell|\ge2$ at every modulus, the surviving
local density is positive (e.g.\ $\prod_{\ell\le300}|Q_\ell|/\ell\approx2.2\times10^{-19}>0$ on the
canonical base), and by CRT \emph{no finite system of congruences (no covering) can annihilate an
unkilled family}: the binary certificate of Proposition~\ref{prop:tailkill} is the complete local theory
of the tail sector, there is no ``second kill layer,'' and the one--sidedness of
Remark~\ref{rem:oneside} is structural rather than an artefact of method. (All three regimes verified
exhaustively on the canonical base: collapse and $|Q_\ell|=(\ell+1)/2$ at
$\ell\in\{3,5,7,101,277\}$; the Weil floor at every prime $281\le\ell\le600$; $A_S,B_S$ carry no factor
below $10^6$, consistent with the balanced--semiprime wall of Remark~\ref{rem:campaign}.)
\end{proposition}

\begin{proposition}[The Pythagorean tail prime is bounded]\label{prop:tailbound}
Let $S$ be a finite set of primes, $D=\prod S$, $N=\sum_{p\in S}D/p$, $A=N+2D$, $B=N-2D$, and suppose
$D<N$. Let $q$ be a prime with $q\nmid 2D$ for which $Aq+D$ and $Bq+D$ are both perfect squares. Then
there are integers $m,c$ with $mc=4D$ and
\[ q^2m^2-4Nq+c^2-4D=0, \qquad \bigl(qm^2-2N\bigr)^2=4\bigl(AB+Dm^2\bigr), \]
so that $q=2(N\pm k)/m^2$ with $k^2=AB+Dm^2$ and $m\mid 4D$. Moreover $m$ is \emph{even}, say
$m=2\alpha$ with $\alpha\mid D$ and $\beta=D/\alpha$, and then $q$ is a root of
\[ \alpha^2q^2-Nq+(\beta^2-D)=0, \qquad\text{so}\qquad
   q=\frac{N+\sqrt{N^2-4D^2+4\alpha^2D}}{2\alpha^2}, \]
a quantity strictly decreasing in $\alpha$. Hence the maximum is at $\alpha=1$ and
\[ q\ \le\ \tfrac12\Bigl(N+\sqrt{N^{2}-4D^{2}+4D}\Bigr)\ =\ tD+O(1),
   \qquad t=\tfrac12\bigl(\sigma+\sqrt{\sigma^{2}-4}\bigr),\ \ \sigma=\sigma(S)=N/D. \]
The largest admissible tail is therefore the base times the very ratio $b/a$ that the mass forces.
At level $59$ this is $9.207\times10^{112}=1.0497\,D$, a factor $7.63$ below $4N$.
\end{proposition}

\begin{proof}
Write $x^2=Aq+D$ and $y^2=Bq+D$ with $x,y\ge0$. Since $A-B=4D$ \emph{exactly},
$x^2-y^2=(A-B)q=4Dq$, that is $(x+y)(x-y)=4Dq$. The prime $q$ divides exactly one of the two factors:
were it to divide both it would divide $2x$ and $2y$, hence $x$ and $y$ as $q$ is odd, hence
$q^2\mid 4Dq$ and $q\mid 4D$, against $q\nmid2D$. Write the factor divisible by $q$ as $qm$ and the
other as $c$, so that $mc=4D$ and, in either case, $2x=qm+c$. Multiplying by $m$,
\[ 2xm\;=\;qm^2+mc\;=\;qm^2+4D, \]
and squaring against $x^2=(N+2D)q+D$ gives
$4m^2\bigl((N+2D)q+D\bigr)=q^2m^4+8Dqm^2+16D^2$. Substituting $16D^2=m^2c^2$ and cancelling $m^2>0$
leaves $q^2m^2-4Nq+c^2-4D=0$; completing the square and using $AB=N^2-4D^2$ gives the second display.
For the parity, $x+y$ and $x-y$ differ by $2y$ and so have equal parity; were both odd their
product would be odd, whereas it is $4Dq$. Hence both are even, and as $q$ is odd, $m$ is even.
Finally $m^2\ge4$ and $c^2\ge0$ give $4Nq=q^2m^2+c^2-4D\ge 4q^2-4D$, that is $Nq\ge q^2-D$; if
$q\ge N+1$ then $q^2-Nq-D\ge q-D>0$ because $D<N<q$, a contradiction, so $q\le N$. (Using only
$m\ge1$ here gives $q\le4N$; the parity removes the factor four.) For the sharp form, $m=2\alpha$ and
$c=2\beta$ with $\alpha\beta=D$ turn $q^2m^2-4Nq+c^2-4D=0$ into $\alpha^2q^2-Nq+\beta^2-D=0$, whose
larger root is $\bigl(N+g(u)\bigr)/(2u)$ with $u=\alpha^2$ and $g(u)^2=N^2-4D^2+4uD$. Its derivative
in $u$ has the sign of $2uD-g(N+g)$, and $4uD=g^2-N^2+4D^2$ turns that into
$-\bigl((g+N)^2-4D^2\bigr)/2$, negative since $N>2D$. So the root decreases in $\alpha$ and is
largest at $\alpha=1$.
\end{proof}

\begin{proposition}[What a barrier improvement must supply]\label{prop:barthreshold}
Let $(a,b)$ be a solution with $a<b$ and put $t=b/a$, so that $t=\sigma(a)$ and, by
$\sigma(a)\sigma(b)=1$,
\[ \sigma(U)\;=\;\sigma(a)+\sigma(b)\;=\;t+\tfrac1t. \]
For $n$ with $T_n>2$ let $t_n>1$ be the root of $t+1/t=T_n$. Since $\sigma(U)\le T_{|U|}$,
\[ \sigma(a)>t_n\ \Longrightarrow\ |U|\ \ge\ n+1, \]
and this is the only route by which a mass argument can raise the barrier. Numerically
\[
\begin{array}{l|cccccc}
n & 59 & 60 & 61 & 62 & 63 & 64\\\hline
T_n & 2.0023502 & 2.0059089 & 2.0094424 & 2.0128554 & 2.0161127 & 2.0193282\\
t_n & 1.0496677 & 1.0798804 & 1.1020081 & 1.1199914 & 1.1352477 & 1.1490254
\end{array}
\]
\end{proposition}

So the barrier level is decided entirely by how close $\sigma(a)$ may lie to $1$, and each further
level demands a definite gap: $|U|\ge61$ needs $\sigma(a)>1.0799$, a full $8\%$ above $1$.

That is why the mass mechanism is exhausted here rather than merely difficult. The only
unconditional lower bound on the gap is $\sigma(a)-1=(a'-a)/a\ge1/a$, which is worthless at
$a>10^{56}$; and $\sigma$ genuinely comes close to $1$ from above on very few primes, the
Sylvester-type set $\{2,3,7,41\}$ giving $\sigma=1.00058\ldots$ with four. Nothing in the cycle
equations forbids $a'-a=1$. A barrier past $60$ must therefore come from a mechanism that bounds
$\sigma(a)$ away from $1$ without appealing to the size of $U$ --- which is what
Theorem~\ref{thm:barrier} does in the opposite direction --- and we know of none.

The same reduction locates the extreme case. Writing $c=b-a$, the identity
$\sigma(U)-2=(t-1)^2/t=c^2/D$ is the minus layer $D'-2D=c^2$, and $\gcd(c,D)=\gcd(b-a,ab)=1$ because
$\gcd(a,b)=1$. Since an admissible support contains every prime $\le167$
\textup{(}Remark~\ref{rem:tier60}\textup{)}, $c$ is divisible by none of them, so $c=1$ or
$c\ge173$. The branch $c=1$ is $D'=2D+1$, equivalently $a'=a+1$ and $(a+1)'=a$; it is the closest a
solution could come to $\sigma(a)=1$, and it is exactly the configuration a mass argument cannot
reach. Below $3\times10^{6}$ the equation $a'=a+1$ has only the solutions $30$, $858$, $1722$ and
$66{,}198$, and none extends: for the first three $a+1$ is prime, so $(a+1)'=1$, and for the last
$a+1$ is not squarefree. The branch is not excluded, only unsearched, since the barrier puts it above
$10^{56}$. The level--$59$
computation that gives $|U|\ge60$ is, in these terms, an exhaustive verification that
$\sigma(a)>t_{59}$, and Proposition~\ref{prop:levelbarrier} says the same verification is unavailable
one level up.

\begin{remark}[The Lehmer reading of the tail is withdrawn]\label{rem:notlehmer}
Remark~\ref{rem:lehmer} below reads the surviving question as the primality of a term $q_n$ of the
\emph{exponential} Pell orbit, and Remark~\ref{rem:whatisleft} names that orbit as the only live
branch. Proposition~\ref{prop:tailbound} withdraws the reading. The orbit is genuinely infinite and
both square conditions do hold along all of it; what fails is that no term past $N$ can be prime.
The single identity responsible is $A-B=4D$, which the earlier treatment used only to form the Pell
system and never to factor $4Dq$.

Three consequences, in increasing order of what they cost.

\emph{The gate is finite, not Lehmer.} The tail primes of a base lie in the explicit finite set
$\{2(N\pm k)/m^2: m\mid 4D,\ k^2=AB+Dm^2\}$. So the second of the three gates is a search over the
divisors of $4D$, not a primality question along an exponential sequence, and the classification of
the residual difficulty as Lehmer--class rather than Schinzel--class does not survive. Checked sound
\emph{and complete} against exhaustive search on $1200$ bases with zero mismatches, and stepped out
along a real orbit, in \texttt{code/tailbound.gp}: on the base $D=1$, $N=7$ (bound $N=7$) the orbit
gives $q_1=7$ prime, then $q_2=741895$, $q_3=76921173511$, and every later term composite, with both
square conditions holding throughout.

\emph{The count reverses sign, and the search has now been run.} The expected number of admissible
$m$ is $\sum_{m\mid 4D}(AB+Dm^2)^{-1/2}\approx D^{-1/2}\sum_{m\mid4D}m^{-1}$, which on the first
immune base is $2.4\times10^{-56}$ and below $10^{-38}$ summed over all $34$ immune families, the
opposite of the ``weakly favours \textsc{yes}'' of Remark~\ref{rem:killcount} on exactly the families
that reading was about. That prediction is now checked. The full divisor space is
$4\cdot2^{58}\approx1.15\times10^{18}$ per family and is not enumerable, but because the hit
probability falls like $1/(m\sqrt D)$ the expected count is carried by small $m$: restricting to
$\omega(m)\le8$ leaves $99.9998\%$ of the expected count. Over all $34$ families that slice is
$3.0777\times10^{11}$ candidates, and
\[ \textup{it contains no }m\textup{ with }AB+Dm^2\textup{ a perfect square} \]
($4{,}933$ survive the modular cascade, exact \texttt{issquare} kills all of them; the earlier
$\omega(m)\le7$ pass, $4.7081\times10^{10}$ candidates and $771$ survivors, agrees).

The same search has been run over \emph{every} admissible base, which matters because
Proposition~\ref{prop:tailbound} requires no primality of $A_S$ or $B_S$, only the algebra. The $34$
immune families are the ones left after the Jacobi kill and the factoring attack, and their immunity
is conditional on a Baillie--PSW verdict; the other $3{,}444$ of the $3{,}478$ open families are open
precisely because $A_S$ or $B_S$ is a balanced semiprime nobody can split. The tail search does not
care. Over all $49{,}961$ admissible bases at $\omega(m)\le4$, which is $99.2\%$ of the expected count
per base, the run is $9.1296\times10^{10}$ candidates, $1{,}469$ survive the modular cascade, and
exact \texttt{issquare} kills all $1{,}469$. Every survivor was additionally checked to satisfy
$m\mid 4D$, the hypothesis the proposition needs. So
\[ \textup{no admissible base admits a Pythagorean tail prime with }\omega(m)\le4, \]
unconditionally and with no primality input. Where Proposition~\ref{prop:tailkill} leaves $3{,}478$
of the one--new--prime families at level $60$ open, of which $34$ only conditionally immune, this
leaves none open on the stated slice. It is a slice result and not a proof: the untested remainder is
$0.8\%$ of the expected count per base, and $\omega(m)\ge5$ is untested.
The run is \texttt{code/tailsearch.rs} (a cascade of quadratic--residue filters modulo primes coprime
to $2D$, $771$ survivors) with every survivor then tested exactly by \texttt{issquare} in
\texttt{code/tailsearch\_verify.gp} ($0$ squares). A positive control recovers the known solutions of
the toy base $D=1$, $N=7$, and an exhaustive exact run at $\omega(m)\le2$ agrees with the filter, so
the cascade is not silently over--rejecting. This is \textsc{verified} on the stated slice, not on the
full divisor set: the residual $0.0018\%$ of the expected count is untested.

The filter primes must be coprime to $2D$, which is itself a structural fact rather than a
convenience. Since $D$ contains every prime up to $167$, for any $\ell\mid D$
\[ AB+Dm^2-N^2=D\bigl(m^2-4D\bigr)\equiv0 \pmod{\ell}, \]
so the value is the square $N^2$ modulo $\ell$ for \emph{every} $m$, and as $\gcd(D,N)=1$ Hensel
lifts that root to every power of $\ell$. The tail search therefore has no local obstruction at any
prime dividing $D$, the analogue of Proposition~\ref{prop:localcomplete} for this formulation; the
classical $64/63/65/11$ pre--filter for squares rejects nothing here, measured at a $75\%$ pass rate
modulo $64$ and $100\%$ modulo $63$. Formalised as \texttt{tail\_value\_mod\_dvd} and
\texttt{hensel\_lift\_identity}; the induction from the step to all powers is routine and is not
formalised.

\emph{The paper's own height estimate now says the same thing.} Remark~\ref{rem:lehmer} puts the
least point of a nonempty fiber at height $\exp\bigl(10^{112}\bigr)$ on regulator grounds, so its
$q\approx x^2/A$ carries some $10^{112}$ digits, against a bound of $115$. Two independent estimates,
one from divisor counting and one from the class--number formula, agree that no immune family admits
a Pythagorean tail prime. Neither is a proof: the bound $q\le N$ is proved, the emptiness is
\textsc{heuristic}.

Formalised in \texttt{lean/Erdos307/TailBound.lean} (\texttt{tail\_factor}, \texttt{tail\_quadratic},
\texttt{tail\_discriminant}, \texttt{tail\_bound}, \texttt{tail\_finite}) on the three standard
axioms, with a non--vacuity witness at the base above.
\end{remark}

\begin{proposition}[The level--by--level programme terminates at $60$]\label{prop:levelbarrier}
Let $U$ be a finite set of primes with $\sum_{p\in U}1/p>2$ and let $B$ be arbitrary. Then there is a
prime $q>B$ with $q\notin U$ such that $U\cup\{q\}$ again has $\sum 1/p>2$. Consequently, for every
$n>59$ there are infinitely many $n$--element prime supports of mass exceeding $2$.
\end{proposition}

\begin{proof}
The primes are infinite, so choose a prime $q$ exceeding both $B$ and $\max U$; then $q\notin U$ and
$\sum_{p\in U\cup\{q\}}1/p=\sum_{p\in U}1/p+1/q>2$. Iterating from the first $59$ primes, whose
reciprocal sum is $N_{59}/P_{59}=2.00235\ldots>2$, produces admissible supports of every cardinality
above $59$, each containing a prime beyond any prescribed bound.
\end{proof}

\begin{remark}[What the level barrier costs, and what it does not]\label{rem:levelbarrier}
Proposition~\ref{prop:close59} closes level $59$ by enumerating its $49{,}961$ admissible supports,
and Proposition~\ref{prop:tailkill} together with the search of Remark~\ref{rem:notlehmer} closes
level $60$ by adjoining one prime to each. The obvious continuation is to enumerate the admissible
$60$--element supports and repeat. Proposition~\ref{prop:levelbarrier} says there is no such
enumeration.

The finiteness at level $59$ is an accident of one inequality being tight: $59$ primes reach mass $2$
only just, at $2.00235\ldots$, so demanding mass $>2$ from exactly $59$ primes pins them near the
smallest $59$. Allow a sixtieth and the first $59$ already carry the mass alone, leaving the sixtieth
entirely free. The mass condition, having been met, constrains nothing further.

Two consequences, and the second is the one that matters. The searches of
Remark~\ref{rem:notlehmer}, however far they are pushed in $\omega(m)$, bear on level $60$ alone;
widening them is not a route upward. And any further progress on \#307 must be
\emph{non--enumerative}: exhausting supports one level at a time is not merely expensive beyond
level $60$, it is impossible, so the method behind every unconditional exclusion in this note is now
spent. That is a no--go, not a resource limit.

Formalised in \texttt{lean/Erdos307/LevelBarrier.lean} (\texttt{admissible\_extends},
\texttt{admissible\_of\_card\_add}, \texttt{level\_enumeration\_terminates}) on the three standard
axioms, with a non--vacuity witness exhibiting the first $59$ primes as a set the hypotheses apply
to.
\end{remark}
\begin{remark}[The immune signal, and where its Lehmer reading broke]\label{rem:lehmer}
Superseded by Proposition~\ref{prop:tailbound} and Remark~\ref{rem:notlehmer}: the orbit described
here is real, but the inference that primality along it is the residual question is not, because the
tail prime is bounded by $N$. What follows is the orbit, the three gates, and the structural points
that survive.

For a reciprocity--immune base the Pell system of Remark~\ref{rem:tier60},
$B_Sx^2-A_Sy^2=-4D_S^2$, is locally soluble everywhere ($\kappa_S>0$), and a Pythagorean pair in the
tail family $S\cup\{q\}$ is \emph{exactly} a solution with $q=(x^2-D_S)/A_S=(y^2-D_S)/B_S$ a
positive prime. Its solutions grow like $\varepsilon^n$ in the fundamental unit of the order of
discriminant $4A_SB_S$ ($224$ digits, $A_SB_S$ a nonsquare, on an explicit immune base). Writing
$x_{n+1}=tx_n-x_{n-1}$ with $x_n=(\alpha^n+\beta^n)/2$ and $\alpha\beta=1$, the squares $x_n^2$ have
characteristic roots $\alpha^2,\beta^2,1$ and the constant shift is absorbed by the root $1$, so
\[ q_n=\frac{x_n^2-D_S}{A_S}\qquad\text{satisfies}\qquad q_{n+3}=(t^2-1)q_{n+2}-(t^2-1)q_{n+1}+q_n, \]
a non--degenerate ternary linear recurrence (verified on four Pell parameters). A
reciprocity--killed base is the complementary degeneration: there the same system is locally
\emph{insoluble} ($\kappa_S=0$), so no candidate exists at all, and the dichotomy of
Proposition~\ref{prop:tailkill} is Pell--solubility versus Pell--insolubility.

The return is milder than ``$k=0$'' suggests: on the Pythagorean locus the deviation $k=(a'-b)/a$ is
an \emph{integer} (the identity $v(D(u)-v)=u(u-D(v))$ with $\gcd(u,v)=1$), forced into $\{0,-1\}$ by
the near--balanced immune ratio, so $k=0$ is a sign choice rather than a measure--zero coincidence,
and the naive Wieferich shortcut that would pin $q$ is vacuous, since $x^2\equiv D_S\pmod q$ holds
identically from $x^2=A_Sq+D_S$.

Neither square condition carries the difficulty alone. Imposing only the minus--square,
$d^2=B_Sq+D_S$, the admissible $d$ satisfy $d^2\equiv D_S\pmod{B_S}$, and along such a class
$d=d_0+tB_S$ one has
\[ q(t)\;=\;\frac{d_0^2-D_S}{B_S}\;+\;2d_0\,t\;+\;B_S\,t^2, \]
a quadratic in $t$ (toy bases: $3+54t+173t^2$, $7+238t+1693t^2$, $2+446t+17357t^2$): a Bunyakovsky
question, of the same class as the $p=5$ family of Proposition~13.15\textup{(iii)} of the full version, and
conjecturally routine. It is the \emph{conjunction} of the two squares that replaces the quadratic
parameter by a fundamental unit and turns the orbit exponential.

\emph{Where the reading broke.} That exponential growth was taken to place the question in the
Lehmer class, on two grounds: no algebraic factorisation is available, $D_S$ being squarefree so
that $x^2-D_S$ is irreducible, and no covering congruence can kill the family, by
Proposition~\ref{prop:localcomplete}. Both observations are correct, and neither is the relevant
one. The factorisation that does exist is not of $x^2-D_S$ but of $x^2-y^2=4D_Sq$, available because
$A_S-B_S=4D_S$ \emph{exactly}: the identity was used here only to build the Pell system, never
multiplicatively. Proposition~\ref{prop:tailbound} draws it, and the ternary recurrence above then
supplies no prime term past $q\le N$.

Three nested gates separate an immune family from a solution: fiber nonemptiness, a tail prime, and
the $k=0$ return. The first is a class--group condition, since Hasse supplies only rational points,
and it is quantitative: immunity makes $A_SB_S$ a product of two distinct primes, hence squarefree,
so the conic lives in the near--maximal order of $\mathbb{Q}(\sqrt{A_SB_S})$ of discriminant
${\sim}A_SB_S\approx10^{224}$, and the class--number formula puts the expected regulator at
$R=\sqrt\Delta\,L(1,\chi)/(2h)\approx10^{112}$; barring an anomalously large class number, the
\emph{least} point of a nonempty fiber sits at height $\exp\bigl(10^{112}\bigr)$. That estimate is
used in Remark~\ref{rem:notlehmer} as independent corroboration that the second gate is empty. A
direct sweep to $q\le10^{12}$ was structurally incapable of reaching even the first gate.
\end{remark}

\begin{remark}[The affine sieve does not reach a rank--one orbit]\label{rem:affinesieve}
The modern instrument for primes in orbits is the affine sieve of Bourgain, Gamburd and Sarnak
\cite{BGS2010}, which produces $r$--almost--primes in orbits of thin groups. It is the one method
designed for exactly the shape of question Remark~\ref{rem:lehmer} isolates, so its failure here is
worth recording precisely rather than left implicit.

The affine sieve needs a spectral gap: the congruence quotients of the acting group must form an
expander family, and that is what supplies the level of distribution the sieve consumes. Here the
$q_n$ run along a Pell orbit, so the acting group is the cyclic group generated by the fundamental
unit of the order of discriminant $4A_SB_S$. Cyclic groups are abelian, hence amenable, and have no
spectral gap; expansion is precisely the property an abelian group lacks. The sieve therefore
returns nothing, and it returns nothing on $2^n-1$ for the same reason. The obstruction is rank, not
thinness as such: non--elementary thin orbits, Apollonian packings among them, do yield
almost--primes.

The natural repair fails for a reason worth stating too. One may embed the rank--one orbit in a
larger thin group whose orbit does expand, and sieve there. What the affine sieve then delivers is
infinitely many almost--primes in the \emph{larger} orbit, inside which our sequence has relative
density zero; the sieve does not localise to a prescribed sub--orbit. Expansion is bought at the
cost of losing the sequence.

This meets the unit--rank trichotomy of Remark~\ref{rem:dichotomy} from a second direction. There,
rank zero gave a classical gate and denied a continuous family, positive rank the reverse. Here,
rank one is exactly the cell in which the affine sieve has no traction. Stated as a single missing
property: what is wanted is a substitute for expansion at rank one, and that is the Mersenne problem
itself, not a lemma about \#307.
\end{remark}

\begin{remark}[The gap to the record, measured]\label{rem:sparserecord}
It is worth stating how far the required statement is from the best that is known, because the
distance is not one of degree.

Write a sequence's \emph{density exponent} $\theta$ for $\#\{n\le X\}\asymp X^{\theta+o(1)}$. The
sparsest sets currently proved to contain infinitely many primes are the values of binary forms:
$a^2+b^4$ at $\theta=3/4$, by Friedlander and Iwaniec \cite{FI1998}, and $a^3+2b^3$ at $\theta=2/3$,
by Heath-Brown \cite{HB2001}, the latter being the sparsest polynomial form known to capture its
primes. Our $q_n$ grow like $\varepsilon^{2n}$, so $\#\{n:q_n\le X\}\asymp\log X$ and
\[ \theta=0. \]

The distance from $2/3$ to $0$ is categorical rather than quantitative. Both records are bilinear
arguments: they run a sieve of dimension one and close it with a Type~II estimate obtained by
averaging over a set of size a positive power of $X$. Every gain in such an argument is a gain of a
power of $X$, and against $\log X$ terms a power of $X$ gains nothing; below $\theta=1/2$ the
standard sieve axioms already fail, and at $\theta=0$ there is no set left to average over. This is
the same wall as Remark~\ref{rem:affinesieve} approached from the analytic rather than the spectral
side, and it is why every sequence of exponent zero, Mersenne, Fibonacci, Lucas and ours alike, is
open together.

The structure our sequence does carry is the standard one and confers nothing. From
$A_Sq_n=x_n^2-D_S$ one gets $x_n^2\equiv D_S$ for every prime $p\mid q_n$ not dividing $D_S$, so all
prime factors of $q_n$ lie in the density--$\tfrac12$ set where $D_S$ is a quadratic residue. That
constrains the factorisations without excluding any, which is exactly the position the primitive
divisor theorems already left us in: divisors, never primality.
\end{remark}

\begin{remark}[The residue, atomised, and where its difficulty comes from]\label{rem:atomised}
Decomposing the question into independently known statements locates the obstruction on a single
atom and identifies its source.

Two atoms carry theorems. Every $q_n$ beyond the thirtieth has a primitive prime divisor
\cite{BHV2001}; and the greatest prime factor satisfies
$P^+(q_n)>n\exp\bigl(\log n/(104\log\log n)\bigr)$, by Stewart \cite{Stewart2013}, which settled
questions of Schinzel and Erd\H{o}s and was the first general advance on Bang and Carmichael. Every
other atom, that $q_n$ be squarefree infinitely often, that $\omega(q_n)$ be bounded infinitely
often, that $q_n$ be a $P_k$, and primality itself, is open.

The chain that would close the crux is $P^+$ large, hence $\omega$ small, hence prime. It breaks at
the first link, and by a measurable amount. Stewart's bound certifies
$\log P^+\ge(1+o(1))\log n$ against $\log q_n\approx2n\log\varepsilon$, so the fraction of the
term's size that is certified is of order $\log n/n$: $10.6\%$ at $n=10$, $0.32\%$ at $n=10^3$,
$4\times10^{-4}\%$ at $n=10^6$. Primality is the fraction $1$. No sharpening of the constant $104$
changes this, since $n^{1+o(1)}$ is the shape linear forms in logarithms produce. The Subspace
Theorem is the other available instrument and does not help: it is stronger under weaker hypotheses
but ineffective \cite{BCZ2003}, and an ineffective statement cannot exhibit the single prime term a
construction needs.

The source is visible in the comparison with the polynomial case. For $n^2+1$ the greatest prime
factor exceeds $n^{1.3}$ infinitely often \cite{Merikoski2019}, that is $q^{0.65}$, a \emph{positive
proportion} of the term's size, and on such sequences almost-primality is available too. For an
exponentially parametrised sequence only $(\log q)^{1+o(1)}$ is available, a vanishing proportion.
The gap between the two regimes is exactly the gap between polynomial and exponential
parametrisation, and Proposition~\ref{prop:nopoly} proves \#307 admits no polynomial one. So the
obstruction recorded in Remark~\ref{rem:affinesieve} as a question of rank, in
Remark~\ref{rem:sparserecord} as one of density exponent, and here as one of $P^+$, is a single
obstruction with three faces, and its source is Proposition~\ref{prop:nopoly}: forced exponentiality
is what places the residue in the class where nothing of the required strength is known.

Two things sharpen that attribution, and both make the obstruction less incidental rather than more.
First, the boundary is exactly where Proposition~\ref{prop:nopoly} sits. The fraction of a term's
size that Stewart's bound certifies is $\log P^+/\log q_n\approx\log n/\log q_n$, so it stays bounded
below precisely when $\log q_n\ll\log n$, that is when $q_n$ grows polynomially: degree $d$ gives the
constant $1/d$. Every super-polynomial rate sends it to $0$, and not only the exponential one:
$\exp(\sqrt n)$ and $n^{\log n}$ do as well. So polynomial growth is not one convenient case among
several, it is the whole of the favourable region, and it is exactly what
Proposition~\ref{prop:nopoly} excludes.

Second, the hypotheses of Proposition~\ref{prop:nopoly} are worth testing, since a gap there would
dissolve everything above. It assumes a \emph{fixed} number of polynomials, and its proof uses that
finiteness when it splits $\sum_i1/p_i$ into constants plus a remainder tending to $0$. A family
whose support grows with the parameter is therefore not covered. That gap is real and useless: a
growing support means a growing number of simultaneous primality conditions, which is a harder
demand than one, not an easier one, and it does not return the family to the polynomial regime. The
proposition forbids the case that would help and leaves uncovered only cases that would not.
\end{remark}

\begin{remark}[The squarefree atom, and what it reduces to]\label{rem:squarefree}
Among the four open atoms one is structurally easier, and it is worth saying why and how far it
gets. Squarefreeness of $q_n$ asks only for an \emph{upper} bound on a count of bad events, never
for the production of a prime, so it is the one atom not subject to the parity obstruction that
blocks bounded $\omega$, $P_k$, and primality.

The argument it wants is a union bound. For a prime $p\nmid2A_SD_S$, $p^2\mid q_n$ forces
$x_n^2\equiv D_S\pmod{p^2}$; the Pell orbit is purely periodic modulo $p^2$ with period
$\pi(p^2)$, there are at most two square roots of $D_S$, and each is met a bounded number of times
per period, so the density of bad $n$ attached to $p$ is at most $C/\pi(p^2)$. Where
$\pi(p^2)=p\,\pi(p)$, which is the generic behaviour, and $\pi(p)\asymp p$, this is $\asymp C/p^2$,
and the union bound closes as soon as $C\sum_pp^{-2}<1$. The constant is comfortable:
$\sum_pp^{-2}=0.4522474\ldots$, so any $C<2.21$ suffices, and a positive proportion of the $q_n$
would be squarefree, hence infinitely many.

What is missing is the hypothesis $\pi(p^2)=p\,\pi(p)$. Its failure is precisely a
Wall--Sun--Sun prime for the sequence, and there the density is $\asymp1/p$ rather than
$\asymp1/p^2$, whose sum diverges and destroys the bound. No such prime is known for any classical
sequence and their finiteness is unproved, so the atom does not close: it \emph{reduces} to a
Wall--Sun--Sun statement. That is a sharper position than ``open'', since it names the obstruction
and shows it is a single hypothesis rather than a programme, but it is a reduction and not a proof.

The unconditional pieces are formalised. \texttt{Erdos307.primeSquareSum\_certificate} gives
$\sum_pp^{-2}<\tfrac12$ in the finite form used above, the first six primes together with the bound
$\tfrac12\sum_{m\ge16}m^{-2}<\tfrac1{30}$ for the odd primes beyond; \texttt{good\_density\_pos} and
\texttt{squarefree\_positive\_density\_of\_bound} are the union bound and the conditional
implication. All carry $0$ \texttt{sorry} and the three standard axioms
\textup{(\texttt{lean/Erdos307/Squarefree.lean})}.

Two further points fix the standing of this atom, and both are negative. First, on the immune bases
the argument is blocked at its \emph{lowest} rung as well as its highest. The density attached to a
single small $p$ is $\#\{n\bmod\pi(p^2):p^2\mid q_n\}/\pi(p^2)$, read off the Pell orbit modulo
$p^2$, and that needs the fundamental solution, the object Proposition~\ref{prop:untestable} places
at $10^{24.7}$ operations. For a Lucas sequence with a known fundamental unit the small-prime rung is
a finite computation; here it is not. The ladder is therefore blocked at both ends, by
non-computability below and by Wall--Sun--Sun above. Second, and more to the point, squarefreeness is
not what the construction needs: Remark~\ref{rem:lehmer} requires $q$ to be \emph{prime}, and a
squarefree $q$ with several prime factors is useless there, so even the dream lemma of this atom
would leave \#307 exactly where it stands. Of the four atoms this is the easiest, and it is also the
only one whose resolution would not advance the problem.
\end{remark}

\begin{theorem}[Parity dichotomy]\label{thm:parity}
Let $a'=b,b'=a$ be a derivative two--cycle (so $a,b$ squarefree, $\gcd(a,b)=1$). Then exactly one of
the following holds.
\begin{enumerate}[label=\textup{(\roman*)},leftmargin=2.2em]
\item \emph{Mixed parity.} Exactly one of $a,b$ is even; say $a=2k$ with $k$ odd. Then $b=a'$ is odd
and $\omega(b)$ is even.
\item \emph{All odd.} Both $a,b$ are odd. Then $P\cup Q$ avoids $2$, so reaching reciprocal sum $>2$
needs $\ge 1412$ odd primes, and $\min(a,b)>10^{2519}$.
\end{enumerate}
\end{theorem}

\begin{proof}
They cannot both be even, as $\gcd(a,b)=\gcd(a,a')=1$. In case~(i), write $a=2k$, $k$ odd squarefree;
$b=a'=k+2k'$ is odd. For odd squarefree $b=\prod_{i}q_i$ one has
$b'=\sum_i b/q_i\equiv \omega(b)\pmod 2$ (each $b/q_i$ is odd). Since $b'=a$ is even,
$\omega(b)\equiv0\pmod2$. In case~(ii) the union of supports omits $2$; the reciprocal sum of the
first $1411$ odd primes is $<2$ while that of the first $1412$ exceeds $2$, so $|P\cup Q|\ge 1412$
and $\prod_{p\in P\cup Q}p\ge 10^{5040}$ (product of the $1412$ smallest odd primes, by direct
computation). Applying the same rigidity/ratio argument as in Theorem~\ref{thm:barrier}(b) (now
with $2\notin U$) gives $\min(a,b)\ge\sqrt{R/T(U)}>10^{2519}$. The all--odd case is therefore
vastly more constrained than the mixed case and, heuristically, empty.
\end{proof}

\begin{proposition}[Longer cycles are harder; two--cycles are extremal]\label{prop:kcycles}
Every member of a $k$--cycle of the arithmetic derivative is squarefree, and for $k\le3$ the members
are pairwise coprime (each pair is consecutive). Their masses $s(a_i)=a_{i+1}/a_i$ multiply to $1$,
so by AM--GM $\sum_i s(a_i)\ge k$; for $k\le3$ this is the union's reciprocal sum, whence the support
contains at least $m_k$ primes, where $m_k$ is least with $\sum_{i\le m_k}1/p_i\ge k$. Now $m_2=59$
and $m_3=361{,}139$ (last prime $5{,}195{,}977$), so a three--cycle has at least $361{,}139$ primes in
its support and product exceeding $10^{2.26\times10^{6}}$ (hence its largest member exceeds
$10^{752{,}194}$); by Mertens $m_k$ grows doubly exponentially. Two--cycles are thus the extremal,
and only practically relevant; case for that measure. For $k\ge4$ non--consecutive members may share
primes, so the displayed bound is one on the \emph{multiset} reciprocal sum; the union's own sum is
bounded in Proposition~\textup{\ref{prop:kuniform}}.
\end{proposition}

\begin{proposition}[The barrier is uniform in the period]\label{prop:kuniform}
Let $a_1\to a_2\to\dots\to a_k\to a_1$ be a cycle of the arithmetic derivative with $k\ge2$ and
$a_i\neq a_{i+1}$, and let $U=\bigcup_i\operatorname{supp}(a_i)$. Then
\[ \sum_{p\in U}\frac1p\ \ge\ \frac{k}{\lfloor k/2\rfloor}, \]
so $|U|\ge m(k)$, the least $n$ with $T_n\ge k/\lfloor k/2\rfloor$, and $\prod_{p\in U}p\ge\Pi_{m(k)}$.
Since $k/\lfloor k/2\rfloor\ge2$ always, with equality exactly for even $k$,
\[ |U|\ \ge\ 59\quad\text{and}\quad \prod_{p\in U}p\ >\ 8.77\times10^{112} \]
for a cycle of \emph{every} period.
\end{proposition}
\begin{proof}
Rigidity gives $\gcd(a_i,a_i')=1$, and $a_{i+1}=a_i'$, so $a_i$ and $a_{i+1}$ are coprime. Hence for
each prime $p$ the set $S_p=\{i:p\mid a_i\}$ contains no two cyclically consecutive indices, i.e.\ it
is an independent set in the cycle graph $C_k$, and therefore $|S_p|\le\lfloor k/2\rfloor$.
Consequently
\[ \sum_{i}\sigma(a_i)\;=\;\sum_{p\in U}\frac{|S_p|}{p}\;\le\;\Bigl\lfloor\tfrac k2\Bigr\rfloor
   \sum_{p\in U}\frac1p . \]
The masses satisfy $\prod_i\sigma(a_i)=\prod_i a_{i+1}/a_i=1$, so $\sum_i\sigma(a_i)\ge k$ by AM--GM.
Combining the two displays gives the claim, and $T$ is increasing.
\end{proof}

The bound is uniform rather than growing, and it is the union support that a search or a sieve sees.
Odd periods are genuinely harder, even ones are not:
\[
\begin{array}{l|cccccccc}
k & 2 & 3 & 4 & 5 & 6 & 7 & 8 & 9\\\hline
k/\lfloor k/2\rfloor & 2 & 3 & 2 & 5/2 & 2 & 7/3 & 2 & 9/4\\
m(k) & 59 & 361{,}139 & 59 & 1{,}413 & 59 & 400 & 59 & 231
\end{array}
\]
so $m(9)<m(7)<m(5)$: the sequence is not monotone in $k$, and $m(k)=59$ \emph{exactly} for every
even $k$, not merely in the limit. The consequence
worth naming is that Theorem~\ref{thm:barrier} is not a statement about period two at all: every
nontrivial cycle of the arithmetic derivative, of any length, carries at least $59$ primes and a
product past $8.77\times10^{112}$. The barrier therefore bounds the whole of
Ufnarovski--\AA hlander's Conjecture~3 and not merely its period--two case, which is \#307.

\begin{proposition}[The barrier as a function of period and smallest prime]\label{prop:twoparam}
Let $a_1\to\dots\to a_k\to a_1$ be a cycle with $k\ge2$, let $U$ be the union of the supports, and
suppose $\min U\ge P$. With $c(k)=k/\lfloor k/2\rfloor$, let $M(k,P)$ be the least $m$ for which the
$m$ smallest primes $\ge P$ have reciprocal sum at least $c(k)$. Then $|U|\ge M(k,P)$, and
$\prod_{p\in U}p$ is at least the product of those $m$ primes.
\end{proposition}
\begin{proof}
Proposition~\ref{prop:kuniform} gives $\sum_{p\in U}1/p\ge c(k)$, and among sets of $m$ primes all
$\ge P$ the reciprocal sum is maximised by the $m$ smallest such primes.
\end{proof}

\[
\begin{array}{l|ccccc}
M(k,P) & k=2 & k=3 & k=5 & k=7 & k=9\\\hline
P=2 & 59 & 361{,}139 & 1{,}413 & 400 & 231\\
P=3 & 1{,}412 & \text{---} & 361{,}138 & 40{,}260 & 15{,}473\\
P=5 & 40{,}259 & \text{---} & \text{---} & \text{---} & 1{,}266{,}094\\
P=7 & 588{,}500 & \text{---} & \text{---} & \text{---} & \text{---}
\end{array}
\]
(dashes exceed $\pi(4\times10^{7})$; \texttt{code/kcycle\_uniform.gp}). The family contains three
results proved elsewhere in this note as special cases, and reproduces each of them exactly:
$M(2,2)=59$ is Theorem~\ref{thm:barrier}; $M(2,3)=1412$, with product $10^{5040.1}$ and hence
$\min(\prod P,\prod Q)>10^{2519}$, is the all--odd branch of Theorem~\ref{thm:parity}; and
$M(3,2)=361{,}139$ is Proposition~\ref{prop:kcycles}. That three independent derivations land on the
same integers is the check the family deserves.

Two readings. The parity dichotomy is not special to period two: an all--odd cycle of \emph{any} even
period needs $1{,}412$ primes and a product past $10^{5040}$, and of period $9$ still needs
$15{,}473$. And the family supplies, in valid form, the conditional statements that
Proposition~18 of \cite{Kovic2012} asserts but does not establish: for $\min U\ge3$ that proposition
claims $r+s\ge57$ where the true bound is $M(2,3)=1{,}412$, and for $\min U\ge5$ it claims
$r+s\ge110$ where the true bound is $M(2,5)=40{,}259$ --- larger by factors of $24.8$ and $366$.

\section{Local anatomy and a heuristic model}\label{sec:anatomy}

\begin{lemma}[Local anatomy]\label{lem:anatomy}
Let $a$ be squarefree and $q$ a prime with $q\nmid a$. Then $q\mid a'$ if and only if
$\sum_{p\mid a}p^{-1}\equiv 0\pmod q$, where $p^{-1}$ is the inverse of $p$ modulo $q$. If $q\mid a$,
then $q\nmid a'$.

\textup{Formalised.} \texttt{Erdos307.anatomy\_iff}, with the factorisation
$a'\equiv a\sum_{p\mid a}p^{-1}$ proved as \texttt{csum\_eq\_mul\_recipSum} rather than assumed; the
second half is \texttt{anatomy\_on\_support} and is Rigidity
\textup{(\texttt{lean/Erdos307/Frame.lean})}.
\end{lemma}
\begin{proof}
$a'=a\sum_{p\mid a}1/p$. If $q\nmid a$ then $q\mid a'$ iff $q\mid a\sum_{p\mid a}p^{-1}$ iff
$\sum_{p\mid a}p^{-1}\equiv0\pmod q$. If $q\mid a$, then by Lemma~\ref{thm:rigidity}
applied to the prime set of $a$, $q\nmid a'$.
\end{proof}

Thus the cycle equation $a'=b$, $b'=a$ factors, prime by prime, into
$\omega(a)+\omega(b)$ local congruences of exactly the shape that governs amicable--pair heuristics:
for each $q\mid b$ one needs $\sum_{p\mid a}p^{-1}\equiv0\pmod q$, and symmetrically. These local
conditions, however, are individually \emph{empty of obstruction}:

\begin{proposition}[No single--modulus congruence obstruction]\label{prop:nolocal}
For every modulus $m$ with $2\le m\le210$, the reduced cross--match system $N_P\equiv D_Q$,
$N_Q\equiv D_P\pmod m$ is realisable by disjoint squarefree prime sets $P,Q$, in both structural
branches: $\gcd(m,D_PD_Q)=1$, and the branch in which a prime divisor of $m$ lies in $P\cup Q$. Hence
no individual modulus in this range can rule out a solution of \#307.
\end{proposition}
\begin{proof}[Construction]
Write $\Sigma_S\equiv\sum_{p\in S}p^{-1}\pmod m$, so $N_S\equiv D_S\Sigma_S$. The reduced system is
$\Sigma_P\equiv D_Q D_P^{-1}$, $\Sigma_Q\equiv D_P D_Q^{-1}\pmod m$; its only coupling is
$\Sigma_P\Sigma_Q\equiv1$, automatic, with $D_P,D_Q$ otherwise free. By Dirichlet each unit class mod
$m$ contains primes, fixing $D_P,D_Q$; appending a prime $r\equiv1\pmod m$ leaves $D_S$ fixed while
sending $N_S\mapsto rN_S+D_S\equiv N_S+D_S$, i.e.\ incrementing $\Sigma_S$ by $1$, so (when
$\gcd(D_S,m)=1$) every target class is reached with disjoint sets. In the branch $\ell\mid m$,
$\ell\in P$, reduction mod $\ell$ collapses the system to $N_Q\equiv D_P\equiv0\pmod\ell$, met by
inverse--paired residues in $Q$; composite $m$ needs no separate gluing, the increment acting directly
mod $m$. We verified both branches for all $2\le m\le210$ (for prime and prime--power moduli the
inside branch requires the directed construction above; seeding a prime divisor of $m$ into $P$,
a blind search missing these as a search--power artefact, not an arithmetic obstruction).
\end{proof}
\begin{remark}[Precise scope]
This excludes only a \emph{single fixed--modulus} congruence disproof. The witnesses are
$m$--dependent (the modulus--$6$ witness already fails modulo $210$), and a single $(P,Q)$ solving the
system for \emph{all} $m$ simultaneously is, for bounded integers, equivalent to the integer identities
$N_P=D_Q$, $N_Q=D_P$ (that is, to \#307 itself) so per--modulus solvability is \emph{no} evidence
for existence. Nor does it exclude other finite or local negative arguments (size/Mertens bounds,
density, magnitude--congruence hybrids, or a Hasse--type failure across infinitely many moduli);
indeed our own $|U|\ge59$ and $2.09\times10^{56}$ barriers are finite--data results fully consistent
with it. Together with Theorem~\ref{thm:ff} (every non--archimedean shadow trivially negative for
degree reasons) it brackets the problem: no single--modulus obstruction over $\mathbb{Z}$, fatal
non--archimedean analogues over $\mathbb{F}_q[t]$, the difficulty living at the archimedean place.
\end{remark}

\paragraph{A first--order count.}
Model a candidate squarefree $a$ as ``$P$--abundant'' ($\sum_{p\mid a}1/p>1$) and ask for
$b=a'$ to itself be $Q$--abundant with $\sum_{q\mid b}1/q=1/s$. Two empirical inputs (Section
\ref{sec:computations}) calibrate this:
\begin{itemize}[leftmargin=1.4em]
\item the density of $P$--abundant $n$ is $\delta=0.0420$ (and $\approx0.0176$ among squarefree
$n$); a theorem, not merely an observation (Proposition~\ref{prop:density});
\item for a qualifying squarefree $a$, $b=a'$ is squarefree about $95\%$ of the time, but its mass
$\sum_{q\mid b}1/q$ has mean $\approx0.047$ against the required $\approx0.934$; the empirical tail
probability $\Pr[\sum_{q\mid b}1/q\ge 1/s]$ is $0$ in $40{,}920$ samples at these scales
(Section~\ref{sec:computations}).
\end{itemize}
The gap is structural: at small scale a $P$--abundant $a$ hogs the cheap mass--carrying small primes,
and $b=a'$, being coprime to $a$, is both too small and too coprime to assemble reciprocal mass
$1/s$. The extremal case is stark: for the primorial $a=\prod_{i\le k}p_i$, though $\sum_{p\mid a}1/p$
rises through $2$ at $k=59$, the cofactor sum $a'$ is essentially prime, with $\sum_{q\mid a'}1/q<0.2$
for every $k<100$ (and only $0.0014$ at $k=59$): the family that \emph{maximises} the forward mass is
the most starved on the return. Equivalently, for squarefree $n$ with $n'$ squarefree,
$n''=n\,\sigma(n)\,\sigma(n')$, so the two--step return ratio $n''/n=\sigma(n)\sigma(n')$ is the \#307
product itself; over \emph{all} such $n<10^6$ it stays strictly below $1$, never exceeding $0.54$
(maximal at $n=969969=3\cdot7\cdot11\cdot13\cdot17\cdot19$), so $D$ is two--step contracting throughout
the tested range and must reverse to unit ratio for a solution to exist. The first scale at which $b$
could plausibly carry the required mass is of the order of the
barrier ($\sim10^{56}$); this dynamical mass--gap is \emph{consistent with} the proved barrier rather
than an independent derivation of it (the empirical tail probability $0$ at small scale is equally
consistent with \textsc{no} and with \textsc{yes}--beyond--the--barrier). A naive ``expected number
of solutions'' built from the local congruences (one factor $\sim1/q$ per condition) diverges like a
constant times $\log X$; this is a divergent naive count of exactly the type known to be unreliable,
so we treat it only as very weak, non--rigorous suggestive evidence that solutions could exist at or
beyond the barrier.

\begin{remark}[Calibration]
We present the count of the previous paragraph as a heuristic, not a theorem. Heuristics of
amicable--pair type are known to be delicate; we report it as weak evidence, consistent with the
divergence of the refined expected count, that the answer to \#307 is plausibly \textsc{yes}
(and hence that the period--one form of the Ufnarovski--\AA hlander conjecture is plausibly false).
\end{remark}

\begin{remark}[A cautionary refinement]\label{rem:refined}
One can try to sharpen this into a quantitative rate using a Kubilius/Erd\H{o}s--Wintner model,
stratified over which small primes lie in the support of $a$ versus remain available to $b$. Two
lessons emerge, both methodological. First, a plain Monte--Carlo estimate \emph{undercounts} the rate
by some $10^{61}$: the expectation is dominated by rare, balanced splits of the small primes that
sampling almost never sees, so a naive estimate of this kind is misleadingly small and only an
explicit stratified sum corrects it. Second, and this is why we record no numerical prediction,
the stratified estimate is itself unstable: refining the stratification from the first $8$ to the
first $14$ primes drops the rate by some $17$ orders of magnitude and pushes the implied first
occurrence from $\sim10^{10^{38}}$ out past $\sim10^{10^{55}}$, with no sign of convergence. The only
robust conclusion is negative: every version places a hypothetical first solution \emph{doubly
exponentially} beyond any conceivable search, consistent with the constructive deficit of
Section~\ref{sec:breeder}; the model does not reliably decide existence in either direction.
\end{remark}

\begin{proposition}[The abundance density is a theorem]\label{prop:density}
The additive function $f(n)=\sum_{p\mid n}1/p$ satisfies the Erd\H{o}s--Wintner three--series
criterion \cite{Tenenbaum} (here $\sum 1/p^2$ and $\sum 1/p^3$ converge), so $f$ has a limiting
distribution and the
density $\delta=\lim_{x\to\infty}\tfrac1x\#\{n\le x:f(n)>1\}$ exists, equal to $\Pr\bigl[\sum_p
X_p/p>1\bigr]$ for independent $X_p\sim\mathrm{Bernoulli}(1/p)$. Moreover the constant carries a
\emph{certified} enclosure:
\[ 0.04202\;\le\;\delta\;\le\;0.04223. \]
Writing $S_1=\sum_{p\le P_1}X_p/p$, $P_1=3\times10^4$, and $T\ge0$ for the tail, monotone upper and
lower cdf envelopes of $S_1$ are propagated through the Bernoulli convolution on a grid of step
$10^{-7}$ with the shift $1/p$ rounded down for the upper and up for the lower envelope (which
preserves them, cdfs being nondecreasing) so $\delta\ge1-U(1)$ costs nothing (the tail is
nonnegative), while $\delta\le1-L(1-a)+\Pr[T\ge a]$ with the explicit Chernoff bound
$\Pr[T\ge a]\le\exp(1.36-P_1a)$, $a=10^{-3}$ (from
$\log\mathbb{E}\,e^{\lambda T}\le\lambda s_2+0.72\lambda^2s_3$ at $\lambda=P_1$, using
$s_2<1/P_1$, $s_3<1/(2P_1^2)$ and $e^x-1\le x+0.72x^2$ on $[0,1]$); float slack $10^{-9}$ exceeds
the accumulated rounding by nine orders, and the machinery reproduces an exact four--prime
reference to $2\times10^{-16}$ (\texttt{code/delta\_cert.py}). The point value $0.0420$, at the
bottom of the enclosure, agrees to four places with the sieve to $2\times10^7$; the upper edge is
conservative, carried almost entirely by the $a$--window. As for the density of abundant integers,
no closed form is expected.
\end{proposition}

\begin{remark}[Two dynamical probes: no second obstruction]\label{rem:dynamical}
Two measurements sharpen the picture beyond the first--order mass count, and both cut towards rather
than against a solution. \emph{(i) Image density.} A two--cycle needs \emph{both} members in the image
$D(\mathbb{N})$ (each is the other's derivative), so one might fear a second obstruction; that the
high--mass numbers a cycle requires are seldom derivatives. The opposite holds: of the integers below
$10^7$, $68.3\%$ lie in $D(\mathbb{N})$, and the fraction \emph{rises} with mass, from $61.7\%$ at
$\sum_{p\mid n}1/p<0.2$ to $78.2\%$ at $\sum_{p\mid n}1/p>1$ and $80\%$ beyond. The high--mass numbers
are thus over--represented among derivatives; the second cycle condition is comfortably met.
\emph{(ii) The barrier is where the balanced mass--product reaches $1$.} Set
$r(n)=\bigl(\sum_{p\mid n}1/p\bigr)\bigl(\sum_{q\mid n'}1/q\bigr)=D^2(n)/n$, so $r=1$ is exactly the
cycle condition. Exhaustively over $n\le X$ the maximum of $r$ climbs $0.50\to0.62$ as $X$ runs
$10^3\to10^9$ (with no exact cycle, as the barrier demands), with no ceiling; the climb follows the
balanced split $r\approx(S(k)/2)^2$, $S(k)=\sum_{i\le k}1/p_i$, which reaches $1$ precisely when
$S(k)=2$, i.e.\ $k\approx59$. The dynamical trajectory thus reproduces the $|U|\ge59$ barrier from the
opposite direction. Neither probe decides existence; both \emph{remove} hypothetical obstructions,
leaving the exact archimedean closure as the sole gate.
\end{remark}

\begin{remark}[Why the analytic method is starved: $N(r)\le1$]\label{rem:nostructure}
The rank--one rigidity (Section~\ref{sec:status}) has a sharp consequence for the circle--method
machinery that resolves the \emph{unrestricted} Egyptian--fraction problem. For a rational $r$ let
$N(r)=\#\{\,P:\sum_{p\in P}1/p=r\,\}$ count the finite prime sets with reciprocal sum $r$. Then
\[N(r)\le1:\]
by the Rigidity Lemma (Theorem~\ref{thm:rigidity}) the sum $\sum_{p\in P}1/p=N_P/D_P$ is already in
lowest terms, so the reduced denominator of $r$ equals $D_P=\prod P$ and $P$ is exactly its set of
prime divisors; uniquely determined ($N(r)=1$ precisely when that denominator is squarefree and the
numerator of $r$ is its cofactor sum). The circle method needs the opposite: a target hit by a positive
proportion of subsets, so a main term dominates. This is exactly what Bloom's shift to $\{p-1\}$
manufactures, $p-1$ being composite and divisor--rich, a rational is the reciprocal sum of
\emph{many} subsets of $\{1/(p-1)\}$ (already among the first fourteen shifted primes the multiplicity
reaches $3$ and grows), whence \cite{BloomShifted} represents every positive rational there. Over the
primes themselves $N(r)\le1$ leaves nothing to average: the analytic engine that reaches $\{p-h\}$ is
provably starved on $\{p\}$, and this is the method--level reason \#307 (the rank--one, exact--prime
case) lies beyond it.
\end{remark}

\section{A function--field lens: the obstruction is archimedean}\label{sec:ff}

Define an arithmetic derivative on $\mathbb{F}_q[t]$ by $P'=1$ for monic irreducibles and the Leibniz
rule, so $f'=\sum_i e_i\,f/P_i$ for $f=\prod_i P_i^{e_i}$ (the exponents $e_i$ read modulo the
characteristic).

\begin{theorem}\label{thm:ff}
Over $\mathbb{F}_q[t]$ this derivative has no two--cycles and no nonzero fixed points. Indeed
$\deg f'\le\deg f-1$ for every nonconstant $f$. The analogues of the integer fixed points $p^{p}$ also
vanish: $(P^{p})'=p\,P^{p-1}P'=0$ in characteristic $p$. The same holds verbatim for every twisted
variant ($P'=u_P\in\mathbb{F}_q^\times$, Leibniz).
\end{theorem}
\begin{proof}
Each term $e_i\,f/P_i$ has degree $\deg f-\deg P_i\le\deg f-1$, so $\deg f'\le\deg f-1$. A two--cycle
$f'=g$, $g'=f$ would force $\deg g\le\deg f-1$ and $\deg f\le\deg g-1$, hence $\deg f\le\deg f-2$;
a fixed point $f'=f$ is excluded likewise. (Verified exhaustively over $\mathbb{F}_2$ to degree $7$,
$\mathbb{F}_3$ to degree $5$, and $\mathbb{F}_5$ to degree $4$; twisted variants over $\mathbb{F}_2$ and
$\mathbb{F}_5$ on small supports through $2.3\times10^{6}$ twist configurations, none closing.)
\end{proof}

\noindent\textup{Formalised.} \texttt{Erdos307.ff\_no\_cycle} is the theorem, on the single
hypothesis $\deg f'+1\le\deg f$; \texttt{deg\_drop\_of\_terms} is that hypothesis from the term
bound, and \texttt{no\_fixed\_point\_of\_drop} excludes fixed points. The argument is stated on the
degrees alone, so it holds verbatim for every twisted variant, the twist multiplying each term by a
unit \textup{(\texttt{lean/Erdos307/FunctionField.lean})}.

\begin{corollary}[The function--field Erd\H{o}s--Barbeau equation is insoluble]\label{cor:ff307}
For every finite field $\mathbb{F}_q$, all nonempty finite sets $P,Q$ of monic irreducibles, and all
unit weights $u_\pi,v_\rho\in\mathbb{F}_q^\times$,
\[
\Bigl(\sum_{\pi\in P}\frac{u_\pi}{\pi}\Bigr)\Bigl(\sum_{\rho\in Q}\frac{v_\rho}{\rho}\Bigr)\ne1.
\]
\end{corollary}
\begin{proof}
Write $f=\prod P$, $g=\prod Q$, $D_u(f)=\sum_\pi u_\pi f/\pi$; the equation reads $D_u(f)\,D_v(g)=fg$.
If some $\pi\in P\cap Q$, then $\pi^2\mid fg$ while $D_u(f)\equiv u_\pi f/\pi\not\equiv0$ and
$D_v(g)\equiv v_\pi g/\pi\not\equiv0\pmod\pi$, so $\pi\nmid D_u(f)D_v(g)$; impossible; hence
$\gcd(f,g)=1$. The same congruence gives $\gcd(f,D_u(f))=\gcd(g,D_v(g))=1$, so $g\mid D_u(f)$ and
$f\mid D_v(g)$, forcing $D_u(f)=c\,g$, $D_v(g)=c^{-1}f$ with $c\in\mathbb{F}_q^\times$: a twisted
two--cycle, excluded by the degree drop. The weighted product equation is thus insoluble over every
$\mathbb{F}_q(t)$; it \emph{is} soluble over the unit--rich number rings of
Theorem~\ref{thm:existence}, since any twisted cycle gives $(D_u(a)/a)(D_v(b)/b)=(b/a)(a/b)=1$, and over
$\mathbb{Q}$, where the admissible weights are $\pm1$, it is \#307 itself with the signed barrier of
Remark~\ref{rem:ordered}, open beyond $2.09\times10^{56}$. The three regimes (non--archimedean
monotonicity (closed, negative), unit--broken monotonicity (closed, positive), ordered--archimedean
(open, barred)) now have theorems on two of three doors.
\end{proof}

Over $\mathbb{Z}$, by contrast, $n'=n\sum_i e_i/p_i$ can exceed $n$ precisely because the archimedean
absolute value lets a sum of strictly smaller terms outgrow the original; the ``mass $>1$''
phenomenon $\sum 1/p>1$, which has no degree--theoretic analogue. Theorem~\ref{thm:ff} thus localises
the entire existence question to the archimedean place: every non--archimedean shadow of \#307 is
trivially negative. This is consistent with (and partly explains) both the census finding of
Section~\ref{sec:computations} and Proposition~\ref{prop:nolocal} (no single modulus obstructs the
cross--match system in the tested range); the difficulty is not local. Any eventual proof must therefore distinguish $\mathbb{Z}$ from $\mathbb{F}_q[t]$; over the
latter the decisive tools are the degree and the Mason--Stothers/Wronskian inequality
\cite{Stothers}, neither of which has a $\mathbb{Z}$--analogue. The conjectural analogue, $abc$,
does not fill the gap, and quantifiably so. Every cycle inherits, from $N'=a^2+b^2$ with $N=ab$, the
relation $(a-b)^2+4ab=(a+b)^2$ on coprime squarefree members, but this triple is $abc$--unexceptional,
of quality $q=\log(a+b)^2/\log\operatorname{rad}\!\bigl((a-b)\cdot ab\cdot(a+b)\bigr)\approx\tfrac12$
(median $0.543$, never above $0.74$ across $6\times10^4$ coprime squarefree pairs; the only $q>1$ cases
are degenerate smooth--gap pairs such as $(5,3)$, vanishingly far below the $10^{56}$ scale a solution
must occupy). Granting $abc$ in its strong explicit form would therefore yield no bound here: the
function--field collapse rests on an inequality genuinely \emph{absent} over $\mathbb{Z}$, not merely
unproven, and the size mismatch is by a factor of two in the exponent. \#307 lives in the
$abc$--comfortable zone.

\begin{remark}[The free grading, and what evaluation destroys]\label{rem:grading}
Theorem~\ref{thm:ff} has a structural reading. In the free Leibniz world (the polynomial ring
$\mathbb{Z}[x_p]$ with $x_p'=1$) the derivative of a squarefree monomial of degree $n$ is the
elementary symmetric form $e_{n-1}$, homogeneous of degree $n-1$: the formal derivative strictly lowers
the grading, cycles are impossible, and over $\mathbb{F}_q[t]$ the degree realises this grading
faithfully. Over $\mathbb{Z}$ the evaluation $x_p\mapsto p$ destroys homogeneity, and the data show it
destroys it completely: for random squarefree $a$ of fixed size, $\omega(a')$ is governed by the size of
$a'$ alone (Erd\H{o}s--Kac), flat in $\omega(a)$ within each parity class (measured: mean
$\omega(a')=2.7$--$3.0$ as $\omega(a)$ runs over $3,5,7$, and $3.6$--$4.0$ over $4,6,8$, the offset being
the parity anatomy of Theorem~\ref{thm:parity}). A two--cycle is therefore a doubly anti--typical event:
an additively assembled integer must carry multiplicative rank far in the Erd\H{o}s--Kac tail \emph{and}
on an exactly prescribed support; the two independent prices whose product is the heuristic rarity of
Section~\ref{sec:anatomy}. The deeper additive--multiplicative content collapses, on inspection, into
the $S$--unit form of Proposition~\ref{prop:sunit}: dividing the cycle equation by $D_Q$ exhibits it as
an $n$--term unit equation; one more coordinate system in which the same exactness reappears
unchanged.
\end{remark}

\begin{remark}[The obstruction is \emph{ordered}--archimedean]\label{rem:ordered}
The barrier is finer than ``archimedean'': its engine is AM--GM on real, \emph{positive} masses, i.e.\
the ordering of $\mathbb{R}$. Over $\mathbb{Z}[i]$ (a PID, so the Rigidity Lemma and cross--matching
survive verbatim, with cycles defined up to units), the masses $\sum1/\pi$ are \emph{complex}: phases can
align so that $\bigl|\sum1/\pi\bigr|>1$ with as few as three Gaussian primes (e.g.\
$\tfrac1{1+i}+\tfrac1{2+i}+\tfrac1{2-i}=1.3-0.5i$), and the $59$--prime bound has no analogue. The
ordered--archimedean thesis thus predicts that small derivative two--cycles \emph{may} exist over
$\mathbb{Z}[i]$ where they are barred over $\mathbb{Z}$, and the prediction is confirmed. Two facts
make the contrast precise.
\emph{First}, phases do not help over $\mathbb{Z}$: for any \emph{signed} derivative ($p'=\pm1$, Leibniz)
a two--cycle still satisfies $|s||t|=1$ with $|s|\le\sigma(a)$, $|t|\le\sigma(b)$ and disjoint supports
(the rigidity proof is unchanged), so $\sigma(a)+\sigma(b)\ge2$ and the full $59$--prime,
$2.09\times10^{56}$ barrier applies to every choice of signs: real phases cannot cancel.
\emph{Second}, fourth roots of unity suffice instantly: Leibniz derivatives on $\mathbb{Z}[i]$ with unit
prime--values ($\pi'\in\{\pm1,\pm i\}$) possess genuine two--cycles at norm ${\sim}10^4$. The smallest:
\[ a=(1+i)(2+7i)(6+11i)=-129-i,\qquad b=3\,(2+i)(1+2i)(6+5i)=-75+90i, \]
with $a'=b$ under the \emph{standard} values $\pi'=1$ on the primes of $a$, and $b'=a$ exactly under
$\pi'=i$ on $3$, $2+i$, $1+2i$ and $\pi'=1$ on $6+5i$; seven prime classes
($N(a)=16{,}642$, $N(b)=13{,}725$) against the fifty--nine that $\mathbb{Z}$ demands. Exhaustive search
finds exactly sixteen such cycles up to conjugation and direction having a member of norm
$\le2\times10^7$ with $\omega\le5$ on the enumerated side ($\omega$--profiles $(3,4)$, $(4,4)$,
$(5,4)$, and even $(4,6)$) and the population grows steadily ($4$ below $3\times10^5$, $16$ below
$2\times10^7$), consistent with an infinite, slowly growing family; in six of the sixteen one side
requires no twist at all. All sixteen were verified exactly, by independent recomputation; no twisted
\emph{fixed} point ($a'=ua$) exists in the same range. The strict two--sided normalisation $\pi'=1$
(first--quadrant representatives) has no cycle up to units (nor any fixed point, zero--derivative, or
unit--derivative element) with $N(a)\le2\times10^9$ ($9.4\times10^8$ candidates, partners by full
factorisation); this is the expected order rather than a separate rigidity, a single fixed phase system
being heuristically $4^{\omega-1}$--fold rarer than the union of all systems. The contrast with
$\mathbb{Z}$ survives this accounting: over $\mathbb{Z}$ \emph{every} one of the $2^{\omega-1}$ sign
systems is barred by the theorem above, while over $\mathbb{Z}[i]$ the union of phase systems is
demonstrably populated from norm $1.7\times10^4$ on. The barrier of \#307 is thus localised precisely: it is the impossibility of phase
cancellation in the ordered field $\mathbb{R}$, where the only available phases $\pm1$ provably gain
nothing; while the moment fourth roots of unity are admitted, the analogous existence problem closes
at seven primes. The cycles found come in conjugate pairs (supports not conjugation--stable); a
conjugation--\emph{stable} Gaussian cycle would descend to rational data through the norm form
$a^2+b^2$, i.e.\ precisely the Pythagorean layer of Proposition~\ref{prop:pairform}; now a concrete
target rather than a speculation.
\end{remark}

\begin{remark}[The descended operator; the symmetric sector]\label{rem:descent}
Conjugation--stable supports admit an exact descent to $\mathbb{Z}$. Such a squarefree Gaussian integer
is an associate of an odd squarefree rational $m$ (split primes entering as whole conjugate pairs; the
$(1+i)$--sector is empty by a short parity argument), and a conjugation--equivariant unit assignment
makes the twisted derivative rational:
\[ D(m)\;=\;\sum_{\substack{p\mid m\\ p\equiv1\,(4)}} t_p\,\frac mp\;+\;\sum_{\substack{q\mid m\\
q\equiv3\,(4)}}\varepsilon_q\,\frac mq, \qquad t_p\in\{\pm2x_p,\pm2y_p\},\quad \varepsilon_q=\pm1, \]
where $p=x_p^2+y_p^2$; a Leibniz derivative on odd squarefree integers whose prime values have size
$2\sqrt p$ rather than $1$, and whose masses ${\sim}2/\sqrt p$ reach $2$ with four primes, so no
$59$--prime barrier exists for it. A two--cycle of $D$ is exactly a conjugation--stable Gaussian cycle
and would plant the phenomenon on rational soil. Parity constrains the hunt: $D(m)$ is odd iff $m$
carries an odd number of primes $\equiv3\pmod4$, so every member must contain such a prime.
One level deeper, since $q^{-1}\equiv q\pmod 8$ for odd $q$, the descended operator obeys
$D(m)\equiv m\,W(m)\pmod 8$ with $W(m)=\sum_{p\mid m}t_p\,p$, so a two--cycle forces
$W(m)\equiv W(m')\equiv mm'\pmod 8$; a mod--$8$ constraint on the sign assignment that prunes
the $4^{\#\mathrm{split}}2^{\#\mathrm{inert}}$ phase enumeration severalfold in future searches
(verified on $2\times10^5$ random assignments; \texttt{code/peeling\_groupoid.py}). An
exhaustive search over members $m\le10^9$ (Gaussian norm $10^{18}$, $\omega\le8$, $1.4\times10^9$
closure tests) found no cycle, no fixed point $|D(m)|=m$, and no vanishing $D(m)=0$. The search
has since been extended two decades further in $m$, on a box narrower in the other parameters:
$\omega\in[4,6]$ and odd primes below $600$, with $m$ running over $[10^9,10^{10}]$ and then
$[10^{10},10^{11}]$. Both tiers are exhaustive over their $5{,}778$ prime pairs and return nothing:
$10{,}004{,}024$ members and $6.94\times10^8$ closure tests on the first, $22{,}218{,}430$ members
and $1.82\times10^9$ closure tests on the second, so $3.22\times10^7$ members and
$2.52\times10^9$ closure tests in all, with no two--cycle at any point
\textup{(\texttt{hunt/descended\_hunt2.rs}, \texttt{hunt/run\_tier2\_par.sh})}. The tier boxes are
not nested in the one above, being narrower in $\omega$ and in the prime range while wider in $m$,
so the two counts complement rather than supersede one another; and the tiers test closure only, not
the fixed--point and vanishing conditions recorded above. We claim no
obstruction from this: equivariance prunes the available phases fourfold per split prime, and a phase
count makes emptiness at this depth the expected order. The descent target remains open over
$\mathbb{Z}[i]$, but not in the real--quadratic world, where the corresponding sector contains a fully
rational pair (Remark~\ref{rem:orderable}). Finally, the
descended coefficients are classical objects: at a split prime $p$ the four values
$\{\pm2x_p,\pm2y_p\}$ are exactly the Frobenius traces of the congruent--number curve $y^2=x^3-x$
(CM by $\mathbb{Z}[i]$) and of its twists ($a_p\in\{\pm2x_p,\pm2y_p\}$ verified by point counting
for all $p<120$) so the conjugation--stable sector is the arithmetic derivative twisted by the
Hecke eigensystem of $\mathbb{Q}(i)$, and its mass statistics are governed unconditionally by Hecke's
equidistribution theorem: the first point of contact between \#307 and automorphic forms. Concretely:
over random squarefree $m$ the three--series conditions converge (the variance term is $\sum 2/p^2$), so
the descended mass $D(m)/m$ has an Erd\H{o}s--Wintner \emph{limiting distribution} whose prime components
are arcsine laws scaled by $1/\sqrt p$; the required input, equidistribution of the Hecke angles, is
visible exactly (Kolmogorov--Smirnov distance $0.0041$ from uniform over the $8{,}977$ split primes below
$2\times10^5$, well inside the $95\%$ band $0.0144$).
\end{remark}

\begin{remark}[The wall is the orderability of $\mathbb{Q}$]\label{rem:orderable}
The cycles over $\mathbb{Z}[i]$ permit a dissection of the barrier that no amount of theory over
$\mathbb{Z}$ could provide. The entire structural chain of this note (rigidity, cross--matching, the
Pythagorean identity, the split discriminant) transfers verbatim to twisted derivatives over
$\mathbb{Z}[i]$: Leibniz holds for disjoint supports, so a cycle $(a,b)$ has $D(N)=a^2+b^2$ for $N=ab$,
with $D(N)-2N=(a-b)^2$ and $2N-D(N)=\bigl(i(a-b)\bigr)^2$ (all verified exactly on the cycles found).
The last equation exposes the wall: over $\mathbb{Z}[i]$, where $-1=i^2$ is a square, \emph{both}
discriminant factors are squares automatically, and the plus/minus distinction of
Proposition~14.20 of the full version (which over $\mathbb{Z}$ carries the entire obstruction, via
$(a-b)^2\ge0\Rightarrow\sigma(N)\ge2$) has no content. The unique non--transferring link is the
sign--definiteness of squares: the barrier is, in the end, the \emph{formal reality} of $\mathbb{Q}$ in
the sense of Artin--Schreier. (The two Artin--Schreier theorems thus bracket the problem: their
characteristic--$2$ map builds the tower of Proposition~13.8 of the full version, their theory of orderable
fields names the wall.) This classifies where the mechanism can operate at all: it needs an ordering
making squares nonnegative (formal reality, hence a real embedding) \emph{and} bounded twisting phases
(finite unit group); by Dirichlet's unit theorem, unit rank $r_1+r_2-1=0$ with $r_1\ge1$ forces degree
one. \emph{$\mathbb{Q}$ is the only number field over which the mass barrier can operate}: over
imaginary quadratic fields squares lose sign--definiteness (over $\mathbb{Z}[i]$ even $-1$ is a square),
and over every other field the unit rank is positive and twisted masses are unbounded. This is a
classification of the wall's habitat, not an existence claim elsewhere, but over $\mathbb{Z}[i]$
existence is no longer a claim: it is a list of sixteen. The ladder of imaginary quadratic rings
confirms the mechanism quantitatively. Over $\mathbb{Z}[\sqrt{-2}]$, whose only units are $\pm1$,
precisely the twisting freedom that provably gains nothing over $\mathbb{Z}$; sign--twisted cycles
nevertheless exist from norm $11{,}462$ on ($a=-90-41\sqrt{-2}$, $b=75-45\sqrt{-2}$, seven prime
classes): the complex embedding alone carries the phenomenon. Over the Eisenstein ring
$\mathbb{Z}[\omega]$, with six units, cycles appear at norm $2{,}964$ with \emph{six} prime classes
($a=-40+22\omega$, $b=56-21\omega$); the smallest derivative two--cycle exhibited in any ring,
and the population at fixed height grows with the unit group: at norm $\le2\times10^7$ the censuses
read $14$, $56$, and $278$ twisted--cycle instances for unit groups of order $2$, $4$, $6$ (a factor
${\sim}4$--$5$ per two extra units, as a phase count predicts) while the strict canonical systems of
all three rings remain empty to norm $2\times10^9$ ($3\times10^9$ candidates in total, partners by full
factorisation). Two further objects appeared over $\mathbb{Z}[\sqrt{-2}]$: a composite with \emph{unit}
derivative ($a'\sim1$; impossible over $\mathbb{Z}$, where $n'=1$ characterises the primes), and an
\emph{anti--holomorphic} cycle: $a=-3163-1474\sqrt{-2}$ satisfies $a'=-\bar a$ exactly, under a single
conjugation--equivariant sign assignment closing both directions ($N(a)=N(b)=14{,}349{,}921$, five prime
classes each side). The equation $a'=u\bar a$ is vacuous over $\mathbb{Z}$ only in its
\emph{one--dimensional} reading, where it says $n'=\pm n$. In the reading that matches the problem it is
not vacuous but central: by Proposition~13.2 of the full version the Pythagorean layer over $\mathbb{Z}$
\emph{is} the anti--holomorphic equation $\mathcal{D}z=\mu\bar z$ for the coordinatewise derivative on
$z=a+bi$, and \#307 is exactly its unit--multiplier case. The ring ladder and the main problem are
therefore instances of one equation rather than analogues; what singles out $\mathbb{Z}$ is the extra
constraint $\operatorname{Im}\mu=1$, which confines the multiplier to a line meeting
$\mathbb{Z}[i]^{\times}$ in the single point $\mu=i$. That the anti--holomorphic sector is thin is
uniform across the ladder: over $\mathbb{Z}[\sqrt{-2}]$ the cycle above is still the \emph{only} one
after widening the census simultaneously in all three directions; prime--norm $\le2\times10^4$,
$\omega\le8$, product--norm $\le10^{10}$ (a factor $700$ past the cycle itself), $8.09\times10^8$
supports, \texttt{hunt/antiholo\_hunt.rs}, and over
$\mathbb{Z}$ the sector is empty below $10^{112}$ by Corollary~\ref{cor:emptytest}. We record the
plateau as consistent with genuine finiteness, not as evidence for it: a random model for
$a'=-\bar a$ predicts a count growing like $\log X$, which the tested range cannot separate from a
constant. All cycles were verified exactly by independent recomputation. Finally, the second pillar was tested in isolation: over
$\mathbb{Z}[\sqrt2]$ (\emph{totally real}, so squares are sign--definite in both embeddings, but with
infinite unit group $\pm(1+\sqrt2)^{\mathbb{Z}}$) twisted cycles appear already at norm $882$: with
unit twists of height $\le2$, $a=21\sqrt2$ (classes $\sqrt2$, $-1+2\sqrt2$, $1+2\sqrt2$, $3$) cycles
exactly with $b=51+13\sqrt2$, of norm $2263$ and only \emph{two} prime classes (against the
fifty--nine that $\mathbb{Z}$ demands) and the composite $369-15\sqrt2$ has derivative exactly
$\varepsilon^6=99+70\sqrt2$, a unit. Both pillars of the classification are thus load--bearing, and each
is now confirmed by explicit counterexample on the soil where it alone fails: drop formal reality and
cycles appear (imaginary quadratic rings, norms $10^3$--$10^4$); keep formal reality but drop unit
finiteness and cycles appear (real quadratic, norm $882$, the smallest in any ring); over $\mathbb{Q}$,
where both hold, the barrier stands at $10^{112}$ by theorem. The bracketing of $\mathbb{Q}$ is complete
from both axes. The deep census sharpens it: at norm $\le2\times10^7$ the real--quadratic ring carries
$30{,}204$ twisted--cycle instances against $14$, $56$, $278$ for the imaginary rings with $2,4,6$ units
infinite unit rank dominates by two orders of magnitude, and contains both the \emph{minimal}
architecture, $a=\sqrt2\,(5-\sqrt2)$ cycling with $b=7$ at norms $(46,49)$, two prime classes per member
(the rational prime $7$ being composite in the ring), the smallest derivative cycle found in any ring;
and a \emph{fully rational} pair:
\[ D(3707)=1547,\qquad D(1547)=3707\qquad(3707=11\cdot337,\ \ 1547=7\cdot13\cdot17), \]
exactly, with unit heights $\le2$ on both sides, verified independently: two ordinary integers in a
derivative two--cycle. The descent phenomenon that is empty over $\mathbb{Z}[i]$ to norm $10^{18}$ is
realised over $\mathbb{Z}[\sqrt2]$ at four digits. The canonical system of $\mathbb{Z}[\sqrt2]$ remains
cycle--free to norm $2\times10^9$ ($7.7\times10^8$ candidates), now with seven unit--derivative
composites.
\end{remark}

The cycle phenomenon over rings with units is governed by the theory of unit equations, which supplies
its rigidity theorem.

\begin{proposition}[$S$--unit rigidity of twisted cycles]\label{prop:sunit}
Let $\mathcal{O}$ be the ring of integers of a number field and $S$ a finite set of primes of
$\mathcal{O}$. If $(a,b)$ is a twisted two--cycle with $\operatorname{supp}(ab)\subseteq S$, then
dividing $D(a)=b$ by $b$ exhibits a solution of
\[ x_1+\dots+x_{\omega(a)}=1,\qquad x_i=u_i\,\frac{a/\pi_i}{b}, \]
in the $S$--unit group of the field, and likewise for $b$. By the theorem of
Evertse--Schlickewei--Schmidt \cite{ESS2002} (cf.\ \cite{Evertse1984}), such an equation has only
finitely many nondegenerate solutions; hence \emph{every fixed support carries at most finitely many
twisted cycles free of vanishing subsums}, and the species can be infinite only through growing
supports; precisely the behaviour of every census above. (Vanishing subsums correspond to partial
zero--derivative patterns; none occurred in any search. The relation was verified exactly on the
$\mathbb{Z}[\sqrt2]$ specimen: the four $S$--unit terms of $a=21\sqrt2$ sum to $1$.)
\end{proposition}

\begin{remark}[Galois folding, and the thinness of symmetric sectors]\label{rem:folding}
Over a quadratic field with automorphism $\sigma$, a $\sigma$--equivariant choice of unit values
($u_{\pi^\sigma}=u_\pi^{\,\sigma}$) makes $D$ commute with conjugation, so the \emph{single} equation
$D(a)=a^\sigma$ already implies that $(a,a^\sigma)$ is a twisted two--cycle (checked symbolically and on
$1{,}988$ random instances). In the two--prime case the partner is forced:
$\rho=-\bigl(\bar u v\,\pi+\bar v\,(\pi^\sigma)^2\bigr)/\bigl(N(u)-N(\pi)\bigr)$, with integrality a
congruence and the entire residue of the problem the primality of $|N(\rho)|$ along Pell--exponentially
growing sequences; Lehmer territory rather than Bunyakovsky. Whether a folded sector is starved or
rich is decided by its degrees of freedom. When the fold leaves no slack the sector is thin: the folded
\emph{two}--prime family over $\mathbb{Z}[\sqrt2]$ is empty in the tested box ($36{,}000$ parameter
triples, none even integral), the conjugation--stable sector over $\mathbb{Z}[i]$ is empty to Gaussian
norm $10^{18}$, and the anti--holomorphic cycle over $\mathbb{Z}[\sqrt{-2}]$ is unique in its census.
But with \emph{three} primes the folded equation
$u_1\pi_2\pi_3+u_2\pi_1\pi_3+u_3\pi_1\pi_2=(\pi_1\pi_2\pi_3)^\sigma$ carries three unit unknowns against
two real equations, and this underdetermined system is \emph{populated}: an exhaustive box ($p_i\le137$,
unit heights $\le4$) contains twenty--four solutions over eleven prime triples, the smallest being
\[ a=(3+\sqrt2)(5-2\sqrt2)(5-\sqrt2)=57-16\sqrt2,\qquad D(a)=a^\sigma=57+16\sqrt2, \]
with unit values $\varepsilon^2$, $\varepsilon$, $-\varepsilon^{-1}$: \emph{the derivative conjugates the
element}, at norm $2737$ on both sides, the equivariant return being automatic (verified exactly).
Existence of twisted cycles over the four rings tested is a theorem by exhibition; for \emph{every} ring
with infinite unit group it is now reduced to the abundance of lattice points on these underdetermined
unit systems (measured plentiful already at $k=3$) a geometry--of--numbers question rather than a
primality one.
\end{remark}

\begin{remark}[{The forced third prime: a Schinzel--type reduction over $\mathbb{Z}[\sqrt2]$, unblocked}]
\label{rem:forced}
Within that abundance the arithmetic is sharper than free lattice--counting. Fixing $\pi_1,\pi_2$ and the
units $u_1,u_2,u_3$, the folded equation rearranges to
$\pi_3\,(u_1\pi_2+u_2\pi_1)-(\pi_1\pi_2)^\sigma\pi_3^\sigma=-u_3\pi_1\pi_2$, a $2\times2$ integer linear
system for $\pi_3=(x,y)$ whose determinant is $N(u_1\pi_2+u_2\pi_1)-N(\pi_1\pi_2)$ (an identity, verified
on all $168$ solutions of the box). So $\pi_3$ is \emph{forced} by $(\pi_1,\pi_2,u_1,u_2,u_3)$, and
infinitude of twisted cycles is exactly the primality of this forced $\pi_3$ for infinitely many such data
a Schinzel/Bunyakovsky--type condition over $\mathbb{Z}[\sqrt2]$. Crucially it is \emph{not} the regime
of Proposition~\ref{prop:nopoly}: the exponential unit freedom $u_i=\pm\varepsilon^{k_i}$ keeps the family
non--polynomial yet populated, so the polynomial--rigidity that forbids reducing the \emph{rational} \#307
to such a hypothesis does not apply here. A Schinzel--type theorem over $\mathbb{Z}[\sqrt2]$ would render
twisted infinitude unconditional, whereas none touches \#307 itself; the twisted problem sits strictly
closer to the reach of current sieve theory than its rational shadow. Concretely the twisted solution set
is \emph{abundant} at accessible scale; thousands of cycles at modest norm (Proposition~\ref{prop:geometry}),
and the count of distinct prime--norm signatures grows roughly in proportion to the norm bound in our census
(e.g.\ $6,12,20$ at bounds $40,80,160$, unit heights saturating); whereas the rational Pythagorean layer
is \emph{empty} below $10^{112}$ (Corollary~\ref{cor:emptytest}). The twisted and rational problems thus differ
not in the difficulty of the search but in whether solutions exist to be found at all: for the twist they are
plentiful and the only question is their primality, exactly the Schinzel--shaped residue above. That residue is a
\emph{pair}, not a triple or a Lehmer sequence: fixing $\pi_1$ and bounding the units, the cycle count still
grows with $N(\pi_2)$ (verified: $33,48,60$ for $\pi_1=3+\sqrt2$ at $N(\pi_2)\le60,120,240$), so $\pi_1$'s
primality is free and only the \emph{simultaneous} primality of $\pi_2$ and the forced $\pi_3$ remains. Setting up
the sieve pins the difficulty precisely. Writing $C=u_1\pi_2+u_2\pi_1$, $E'=(\pi_1\pi_2)^\sigma$,
$F=-u_3\pi_1\pi_2$, the forced third prime is $\pi_3=(\bar CF+E'\bar F)/\det$ with $\det=N(C)-N(E')$ a binary
quadratic form in $(c,d)=\pi_2$. The first condition is the clean form $N(\pi_2)=c^2-2d^2$; the second, however,
is \emph{not} a second form but the rational condition $N(\pi_3)=N(G)/\det^2$ ($G=\bar CF+E'\bar F$ quadratic,
$N(G)$ quartic, $\gcd(N(G),\det)=1$), evaluated on the integral locus $\det\mid G$, and $N(\pi_3)$ is unbounded,
spiking as $\det\to0$ (values to ${\sim}10^{11}$ already at $N(\pi_2)\le2\times10^4$). On this \emph{frozen}
$(\pi_1,\text{units})$ slice the question is primality along the folded conic bundle, and the slice is thin: its
unimodular fibres $\det=\pm1$ are \emph{empty} by a genus obstruction (there $\det=\pm1$ would need
$x^2-2y^2\in\{-55,-43\}$, both non--residues). But the thinness is an artefact of freezing exactly the parameters
that carry the abundance. The \emph{full} twisted variety $D(a)=a^\sigma$ ($\pi_1,\pi_2,\pi_3$ and their units all
free) is a cubic threefold whose solution count grows with positive exponent; a positive--density attack must run
there, not on a fixed fibre, so the fixed--fibre conic sieve is a closed dead end. The honest placement is then: the
reachable object is the full threefold at the Heath--Brown $x^3+2y^3$ frontier, compounded by a \emph{coupled} second
primality ($N(\pi_2)$ and $N(\pi_3)$ simultaneously prime); the two--sparse--conditions/parity barrier where the
circle method loses minor--arc cancellation. This is strictly below the exponential/Lehmer tier of the immune
families (Remark~\ref{rem:lehmer}, where no sieve reaches at all), but a generation beyond present technology; it is
where genuinely new analytic input for the twist would first have to bite.
\end{remark}

The natural attempt to \emph{prove} that infinitude (forcing one prime to be the conjugate of another
so that a single free prime remains, whence Dirichlet would suffice) is blocked, and the block is a
clean structural law.

\begin{proposition}[No conjugate prime pair in a twisted cycle]\label{prop:noconj}
Let $D$ be a Leibniz derivative with unit prime--values over a quadratic ring with automorphism $\sigma$,
and $a$ a squarefree twisted cycle $D(a)=a^\sigma$. Then no two prime factors of $a$ are
$\sigma$--conjugate. Equivalently a twisted cycle's support is disjoint from its conjugate; it is never
conjugation--stable, which is exactly why the exhibited $\mathbb{Z}[i]$ cycles are asymmetric, coming in
conjugate pairs of \emph{distinct} supports rather than on one symmetric support. This \emph{separates}
two problems easily conflated: the twisted equality $D(a)=a^\sigma$ treated here, and the symmetric
target of Remark~\ref{rem:descent}, which is a genuine two--cycle $D(m)=m'$ ($m\ne m'$, both
conjugation--fixed) of the \emph{descended} rational operator; not a twisted equality. The present
proposition bars the first from the symmetric locus; it does not bear on the second, which remains
genuinely open (the exhaustive emptiness there, Remark~\ref{rem:descent}, is unexplained by it).
\end{proposition}
\begin{proof}
Suppose a split prime $\pi$ and its distinct conjugate $\pi^\sigma$ both divide $a$. Reducing
$D(a)=\sum_i u_i\,(a/\pi_i)$ modulo $\pi$ kills every term but the one at $\pi_i=\pi$, so
$D(a)\equiv u_\pi\,(a/\pi)\not\equiv0\pmod\pi$ ($u_\pi$ a unit and $\pi\nmid a/\pi$ by squarefreeness).
Yet $a^\sigma=\prod_i\pi_i^\sigma$ carries the factor $(\pi^\sigma)^\sigma=\pi$, so
$a^\sigma\equiv0\pmod\pi$; hence $D(a)\ne a^\sigma$. An exhaustive $\mathbb{Z}[\sqrt2]$ search confirms
no cycle carries a conjugate pair.
\end{proof}

We codify what is theorem and what remains open in this circle.

\begin{theorem}[Existence of twisted derivative two--cycles]\label{thm:existence}
Leibniz derivatives with unit prime--values admit genuine two--cycles over $\mathbb{Z}[i]$,
$\mathbb{Z}[\sqrt{-2}]$, $\mathbb{Z}[\omega]$, and $\mathbb{Z}[\sqrt2]$. The proof is by exhibited
certificate, each verified exactly by independent recomputation: the smallest instances are, in order of
norm, $\bigl(\sqrt2(5-\sqrt2),\,7\bigr)$ over $\mathbb{Z}[\sqrt2]$ at norms $(46,49)$; the Eisenstein
cycle at norm $2{,}964$; the sign--twisted cycle over $\mathbb{Z}[\sqrt{-2}]$ at norm $11{,}462$, and the
cycle over $\mathbb{Z}[i]$ at norm $16{,}642$; together with the rational pair
$D(3707)=1547$, $D(1547)=3707$ over $\mathbb{Z}[\sqrt2]$.
\end{theorem}

\begin{proposition}[Geometry of the solution set]\label{prop:geometry}
Over $\mathbb{Z}[\sqrt2]$, consider the Galois--folded system
$u_1\pi_2\pi_3+u_2\pi_1\pi_3+u_3\pi_1\pi_2=(\pi_1\pi_2\pi_3)^\sigma$ of
Remark~\textup{\ref{rem:folding}}. \textup{(i)} For \emph{every} triple of split primes the system is
solvable with the $u_i$ on the real unit hyperbola $\{|st|=1\}$: solving the second embedding for $u_2$
and letting the third coordinate tend to zero, the first--embedding residual sweeps $\pm\infty$ while the
remaining terms stay bounded, so a root exists by continuity; the real solution set is a nonempty curve
\textup(verified numerically on $200$ random triples, all solvable\textup). The volume of the solution
variety never vanishes. \textup{(ii)} By Proposition~\textup{\ref{prop:sunit}} each fixed support carries
finitely many integral solutions, so infinitude can only come from growing supports. \textup{(iii)} The
measured integral hit rates are $11$ of $455$ prime triples in the smallest box, and $30{,}204$ cycles of
norm $\le2\times10^7$ at twist height $\le2$. \textup{(iv)} What remains open is precisely whether the
discrete unit lattice meets the real curve in integer points for infinitely many supports,
inhomogeneous Diophantine approximation on exponentially parametrised curves. Two shortcuts are closed:
bounded--unit polynomial families are obstructed (the rational component of the return forces unit
coefficients to track growing prime coordinates, which bounded units cannot), and even the cleanest
forward family ($D\bigl(\sqrt2(x-\sqrt2)\bigr)=x+2$ \emph{identically in} $x$, supplying a candidate
cycle whenever $x^2-2$ and $x+2$ are prime) has sporadic returns: the closing two--term unit equation,
effectively finite for each $x$, is solvable at $x=5$ and provably insolvable at $x=15$ and $x=21$.
Those failures are, however, exactly the rigid case: a \emph{prime} return $x\pm2$ leaves the closing
system no slack. The identity has four unit branches $D\bigl(\sqrt2(x\pm\sqrt2)\bigr)=x\pm2$ (independent
signs, twists $1$ and $\pm\varepsilon^{\pm1}$), and when $x\pm2$ is composite squarefree the closing
system gains one unit unknown per prime class against the same two equations, and does close:
$x=107,121,387,1693$ yield cycles with rational partners $b=105,119,385,1695$, each verified exactly.
Infinitude along this family is thereby reduced to Bunyakovsky for $x^2-2$ together with recurring
composite--return closure: the softest form the gap takes anywhere in this paper. Even the primality
half of that reduction can be removed. By the half--dimensional sieve
\cite{Iwaniec1978,LemkeOliver2012} (applied to $(2y+1)^2-2=4y^2+4y-1$ to keep $x$ odd), $x^2-2$ is a
$P_2$ infinitely often; every prime factor of $x^2-2$ has $2$ as a quadratic residue and therefore
splits in $\mathbb{Z}[\sqrt2]$, so $\sqrt2(x-\sqrt2)$ is infinitely often squarefree and all--split
with $\omega\in\{2,3\}$. The candidate supports thus exist \emph{unconditionally}; all that remains
open on this route is whether the unit gates open along that supply infinitely often.
Exactness migrates; it does not disappear.
\end{proposition}

\begin{proposition}[No vanishing derivatives, no transport, and the price of exactness]\label{prop:novanish}
\textup{(i)} No squarefree $z$ with $\omega(z)\ge1$ has $D(z)=0$, for any twisted derivative over any of
the rings considered (indeed over any UFD): modulo a prime $\pi\mid z$ every cofactor term but one
vanishes, leaving $D(z)\equiv u_\pi\,z/\pi\not\equiv0\pmod\pi$. The zero--derivative counters of every
census above are therefore theorems, not observations. \textup{(ii)} Consequently cycles cannot be
transported multiplicatively: for squarefree $c$ coprime to $ab$, Leibniz gives $D(ca)=D(c)\,a+c\,D(a)$,
so $(ca,cb)$ is a cycle exactly when $D(c)=0$; impossible. Cycles do not breed; each one requires
fresh exactness. \textup{(iii)} The two cycle equations can be bought by construction, at a conserved
price. Buying \emph{neither} leaves the generic stratum: two exact unit conditions ($30{,}204$
specimens). Buying \emph{one} is free in two dual ways; forward, $\pi\coloneqq q-\varepsilon^m\sqrt2$
makes $D(\sqrt2\,\pi)=q$ an identity (the four--branch family above is the case $m=\pm1$); return,
$\pi\coloneqq(\rho+\varepsilon^{-1}\rho^\sigma)/\sqrt2=(c+2d)-d\sqrt2$ for $\rho=c+d\sqrt2$ closes
$D(b)=\sqrt2\,\pi$ identically for every associate $b$ of $\rho\rho^\sigma$, and the residue is one
primality, $|N(\pi)|$, plus one exact unit gate, measured at $2$ of $50$ pairs with both primalities in
the smallest box. The two hits re--derive, with no search, precisely the two census minimal--shape cycles
with prime rational partner: norms $(46,49)$ and $(206,49)$. Buying \emph{both} collapses the parameters
to unit--indexed quadratics in $q$ with a square--discriminant condition; the resulting $q=7$ line
carries the known cycles at $m=1$ and $m=-2$ (discriminants $25$ and $81$) together with one Galois
image, and every other prime--norm rung out to $|m|\le14$ fails its return, verified exactly. Each
purchased identity is repaid as a sparser arithmetic condition elsewhere: the hardness of exactness is
conserved, never destroyed.
\end{proposition}

\begin{remark}[The two faces of the gap: unit rank dichotomises the obstruction]\label{rem:dichotomy}
The invariant that governs abundance in Remark~\ref{rem:orderable} (the unit rank
$r=r_1+r_2-1$ of Dirichlet) also governs the \emph{type} of the residual open problem, and the two
types are mutually exclusive. When $r=0$ (the rings $\mathbb{Z}$, $\mathbb{Z}[i]$,
$\mathbb{Z}[\sqrt{-2}]$, $\mathbb{Z}[\omega]$) the unit group is finite, so a given pair of supports
admits only $|\mathcal{O}^\times|$ twist assignments and hence finitely many candidate cycles: the
solution locus is a set of \emph{isolated points}, zero--dimensional, and infinitude can arise only from
infinitely many \emph{distinct} supports independently passing a closing test; a
Bateman--Horn/Schinzel--type recurrence over prime tuples, with no continuous family to deform along.
The constituent primalities are individually classical (each support prime is a value of a binary
norm--form, so by Fermat--Dirichlet (and, for $\mathbb{Z}[i]$, by the Green--Sawhney theorem with
\emph{both} coordinates prime) there is no shortage of supports); it is their \emph{simultaneous}
closure that is not classical. This was tested directly: the minimal shape $a=(1+i)\pi$ over
$\mathbb{Z}[i]$ yields \emph{no} cycle for any Gaussian prime $\pi$ of norm $\le4000$ under any of the
sixteen twists, and the sixteen genuine cycles are irregular mixtures of inert and split classes
($\omega=3,4,5$) with no one--parameter family among them; the signature of isolated points, not a
curve. When $r\ge1$ (every real quadratic ring, e.g.\ $\mathbb{Z}[\sqrt2]$),
Proposition~\ref{prop:geometry} makes the real solution locus a positive--dimensional curve in each
support, but the unit lattice is an infinite rank--$r$ discrete set, and infinitude requires it to meet
that curve in \emph{integer} points infinitely often; inhomogeneous Diophantine approximation on an
exponentially parametrised curve. No unit rank delivers both a classical gate and a continuous family:
rank zero gives the former and denies the latter, positive rank the reverse. This is why no single ring
lets the infinitude be \emph{finished} by present methods; a structural account of the obstruction's
universality, not a barrier theorem. It also locates $\mathbb{Z}$ exactly: the unique ring carrying
\emph{both} the rank--zero sporadicity \emph{and}, by formal reality, the mass barrier; the worst
cell of the trichotomy of Corollary~\ref{cor:ff307}, which is precisely why \#307 is the hard
representative of its family.
\end{remark}

\section{A Th\=abit--type rule and its honest limits}\label{sec:frame}

The bridge recasts the search for a solution as the search for a derivative two--cycle, and the
derivative's Leibniz structure makes the cycle equations \emph{linear} in the ``closing'' primes.
This yields a constructive calculus analogous to Th\=abit ibn Qurra's rule for amicable pairs and to
Euler's divisor method.

\begin{theorem}[Frame Rule]\label{thm:frame}
Let $M,N$ be coprime squarefree integers with derivatives $M',N'$, and set $\Delta\coloneqq MN-M'N'$.
Suppose
\[
p=\frac{MN'+N^2}{\Delta},\qquad q=\frac{M'N+M^2}{\Delta}
\]
are integers, both prime, each coprime to $MN$, with $p\ne q$. Then $a\coloneqq Mp$, $b\coloneqq Nq$
form a derivative two--cycle $a'=b$, $b'=a$; equivalently $(P,Q)$ with $P$ the prime set of $Mp$ and
$Q$ that of $Nq$ solves Erd\H{o}s \#307.

\textup{Formalised.} \texttt{Erdos307.frame\_cycle} verifies both cycle equations from the two
divisibility hypotheses, and \texttt{frame\_solve} runs the derivation backwards, so the two
quotients are \emph{forced} rather than guessed and the rule is an equivalence. The Leibniz step
$a'=M'p+M$ is \texttt{csum\_insert\_prime}. What no algebra reaches is the primality of the two
quotients \textup{(\texttt{lean/Erdos307/Frame.lean})}.
\end{theorem}

\begin{proof}
With $p$ prime and $p\nmid M$, $a=Mp$ is squarefree and $a'=M'p+M$; similarly $b'=N'q+N$. The cycle
conditions $a'=b$, $b'=a$ read
\[
M'p+M=Nq,\qquad N'q+N=Mp,
\]
a linear system in $(p,q)$ with determinant $-(MN-M'N')=-\Delta$. Solving gives exactly
$p=(MN'+N^2)/\Delta$ and $q=(M'N+M^2)/\Delta$; substituting back verifies both equations
identically. (Both identities were checked symbolically.)
\end{proof}

A two--primes--per--side variant (``Rule~2'') is obtained identically: with $a=Mp_1p_2$, $b=Nq_1q_2$
and elementary symmetric functions $e_1,e_2,f_1,f_2$, the cycle equations are linear in the $f_i$,
\[
f_2=\frac{Me_1+M'e_2}{N},\qquad f_1=\frac{MNe_2-MN'e_1-M'N'e_2}{N^2},
\]
and a solution is realised whenever $f_1^2-4f_2$ is a perfect square with prime roots.

\paragraph{An Euler-style divisor trick.}
Fix $N$ and let $M=M_0r$ with $r$ a new prime, so $M'=M_0'r+M_0$. Put
$\Delta_0\coloneqq M_0N-M_0'N'$.

\begin{proposition}[Divisor trick]\label{prop:divisor}
With the above notation and $p_{\mathrm{num}}\coloneqq MN'+N^2$:
\begin{enumerate}[label=\textup{(\alph*)},leftmargin=2.2em]
\item $\Delta(r)\coloneqq MN-M'N'=r\,\Delta_0-M_0N'$ \emph{(linear in $r$)};
\item $\Delta_0\,p_{\mathrm{num}}-M_0N'\,\Delta(r)=K$ where $K\coloneqq \Delta_0N^2+(M_0N')^2$ is
\emph{constant in $r$};
\item consequently $\Delta(r)\mid p_{\mathrm{num}}\Rightarrow \Delta(r)\mid K$, and the converse
holds provided $\gcd(\Delta(r),\Delta_0)=1$ \textup{(}equivalently $\gcd(M_0N',\Delta_0)=1$, since
$\Delta(r)\equiv-M_0N'\bmod\Delta_0$\textup{)};
\item each positive divisor $d\mid K$ with $d\equiv-M_0N'\ (\mathrm{mod}\ \Delta_0)$ yields a
candidate $r=(d+M_0N')/\Delta_0$ and recovers $\Delta(r)=d$; this enumeration is complete (it
captures every genuine solution). The resulting $p=p_{\mathrm{num}}/d$ is integral when
$\gcd(M_0N',\Delta_0)=1$, and otherwise must be checked.
\end{enumerate}
\end{proposition}

\begin{proof}
(a) and (b) are direct expansions (verified symbolically; $\partial K/\partial r=0$). (c): from
(b), $\Delta(r)\mid p_{\mathrm{num}}\Rightarrow\Delta(r)\mid\Delta_0p_{\mathrm{num}}
=M_0N'\Delta(r)+K\Rightarrow\Delta(r)\mid K$. For the converse, $\Delta(r)\mid K$ gives
$\Delta(r)\mid\Delta_0p_{\mathrm{num}}$, whence $\Delta(r)\mid p_{\mathrm{num}}$ iff
$\gcd(\Delta(r),\Delta_0)=1$. (d) inverts (a). Excluded degeneracies: $\Delta_0=0$ ($r$ undefined)
and $K\le0$.
\end{proof}

Thus, exactly as in Euler's amicable--pair machinery, one factors the single number $K$ once and
reads off candidates from its divisors lying in a fixed residue class modulo $\Delta_0$; only the
primality of $r,p,q$ then remains. Two caveats deserve emphasis:
the converse in (c) and the automatic integrality of $p$ require $\gcd(M_0N',\Delta_0)=1$ (it fails,
e.g., for $M_0=2$, $N=19\cdot47$, $r=5$, where $\Delta(r)\mid K$ but $\Delta(r)\nmid
p_{\mathrm{num}}$), and the count of admissible classes is governed by $\Delta_0$, not merely by
$\tau(K)$.

When the closing side carries a \emph{single} new prime, the rule collapses to a normal form that
separates the two layers of the problem.

\begin{proposition}[One--prime normal form; the exactness stratum]\label{prop:oneprime}
\textup{(a)} Erd\H{o}s \#307 has a solution if and only if there exist squarefree $M,N$ with
\[ MN-M'N'=M^2, \]
such that $p\coloneqq N'/M$ is a prime not dividing $M$. Indeed the equation forces $M\mid M'N'$, hence
$M\mid N'$ by $\gcd(M,M')=1$, and $N=M+M'p$; then $a=Mp$, $b=N$ satisfy $a'=M'p+M=N$ and $b'=N'=Mp=a$
(coprimality and $a\ne b$ are automatic). Conversely every solution $(a,b)$ and every prime $p\mid a$
yield such a pair with $M=a/p$, so each solution carries $\omega(a)$ witnesses. The problem thus splits
into an \emph{exactness} layer (the Diophantine equation) and a single \emph{primality} layer (one prime
value of the determined quantity $N'/M$).
\textup{(b)} \emph{The exactness layer alone carries the full barrier.} Call $(M,N)$ a \emph{relaxed
witness} if $MN-M'N'=M^2$ with $p=N'/M$ composite. The identity
$\bigl(\sigma(M)+\tfrac1p\bigr)\sigma(N)=\frac{N}{Mp}\cdot\frac{Mp}{N}=1$ holds for \emph{every} witness,
so by strict AM--GM ($\sigma(N)=1$ being impossible) $\sigma(M)+\sigma(N)>2-\tfrac1p$. Rigidity forces
three structural facts. The supports of $M$ and $N$ are \emph{automatically} disjoint: a common prime
$q$ would divide $N'=Mp$ through $q\mid M$, against $\gcd(N,N')=1$. Likewise $\gcd(p,MN)=1$: a prime
shared by $p$ and $M$ propagates to $N=M+M'p$ and lands in the previous case, while one shared by $p$
and $N$ divides $N'$. And if $p$ is even, both members are odd: $2\mid N'=Mp$ forces $N$ odd by the same
lemma, while $M$ even would make $N=M+M'p$ even. The union of supports therefore avoids every prime
divisor of $p$, and for even $p$ avoids $2$; exact greedy computation then gives, for \emph{every}
composite $p$,
\[ \omega(MN)\ \ge\ 59, \qquad MN\ >\ 8.77\times10^{112} \]
the cycle barrier itself; with strictly stronger floors at every finite $p$: $\omega\ge194$,
$MN>10^{497}$ at $p=9$; $\omega\ge230$, $MN>10^{609}$ at $p=4$; the even--$p$ floors rising towards the
all--odd bound $10^{5040}$, and the infimum $8.77\times10^{112}$ approached only as $p\to\infty$ through
composites free of prime factors $\le277$. Deleting primality thus lowers the barrier by
\emph{nothing}: the $59$--prime, $10^{112}$ wall is carried entirely by the exactness equation, and the
primality of $N'/M$ is a separate demand on top of it. (No witness, prime or composite, exists with
$M\le120$, $p\le2000$, exhaustively.)
\end{proposition}

\paragraph{The decisive obstruction: a circular objective.}
The heuristic yield of the divisor trick from one frame is
\[
\mathbb{E}[\text{cycles}] \ \approx\ A\cdot\frac{1}{\ln r\,\ln p\,\ln q},
\qquad A=\#\{d\mid K:\ d\equiv -M_0N'\ (\mathrm{mod}\ \Delta_0)\}\ \approx\ \frac{\tau(K)}{\Delta_0},
\]
the last step being an equidistribution heuristic. To make $\mathbb{E}\ge1$ at the barrier scale
($1/\ln^3\approx 5\times10^{-7}$) one needs $A\gtrsim 2\times10^6$ admissible divisors per frame,
hence $\Delta_0=O(1)$ with $\tau(K)\sim10^6$. The natural objective is therefore ``engineer frames
with $\Delta_0$ small.'' But here is the obstruction, an exact identity:
\[
\boxed{\ \Delta_0 \;=\; M_0N\bigl(1-s_{M_0}\,s_N\bigr),\qquad
s_{M_0}=\!\!\sum_{p\mid M_0}\!\frac1p,\ \ s_N=\sum_{p\mid N}\frac1p.\ }
\]
So $\Delta_0$ is small \emph{relative to} $M_0N$ exactly when $s_{M_0}s_N\approx1$ (which is the
Erd\H{o}s~\#307 condition for the frames themselves) and $\Delta_0=O(1)$ in absolute terms at
scale requires $s_{M_0}s_N=1-O(1/M_0N)$, i.e.\ a reciprocal--product within $O(1/M_0N)$ of $1$,
\emph{strictly harder} than \#307. The frame--engineering objective is thus circular: the Frame Rule
is a faithful \emph{parametrisation} of solutions, not a reduction in difficulty.

\begin{remark}[Honest status of the constructive program]
The Frame Rule (Theorem~\ref{thm:frame}) and divisor trick (Proposition~\ref{prop:divisor}) are
exact theorems; they make the existence of a solution \emph{falsifiable in principle by finite
computation per frame}, and they place \#307 in the same constructive lineage as Th\=abit's and
Euler's rules. They do not, as they stand, lower the difficulty below \#307 itself: by the boxed
identity the bottleneck (small $\Delta_0$) is the original problem in disguise, and a complete
small--frame sweep ($2.2\times10^5$ frames) produces zero cycles with total expected count $<1$, as
the barrier demands. We regard the calculus as an exact parametrisation of solutions whose central
obstruction we have now identified precisely; not a pathway that lowers the difficulty, and not
evidence that a solution is within reach.
\end{remark}

\section{A breeder form and the feasibility of certificate search}\label{sec:breeder}

Allowing \emph{two} free primes on one side turns the divisor trick into a cleaner bilinear form,
Euler's amicable--pair method, in the shape used by te~Riele's breeding algorithms
\cite{teRiele1984}, and lets us state precisely the strongest honest conclusion: a
\emph{conditional} reduction, which nonetheless does not become a feasible finite search.

\begin{proposition}[Bilinear breeder form]\label{prop:bde}
Fix coprime squarefree $M_0,N$ and seek a two--cycle $a=M_0rp$, $b=Nq$ with $r,p,q$ distinct primes
coprime to $M_0N$. Put $G:=NM_0-N'M_0'\ (=\Delta_0)$, $c:=M_0N'$, and $K:=GN^2+c^2$. Then the cycle
equations $a'=b$, $b'=a$ are equivalent to
\[
(Gr-c)(Gp-c)=K, \qquad q=\frac{M_0rp-N}{N'}.
\]
\end{proposition}
\begin{proof}
Expanding, $(Gr-c)(Gp-c)-K=-N'G\,(a'-b)$ (verified symbolically), so the bilinear equation is
equivalent to $a'=b$; the second cycle equation $b'=a$ then determines $q$ as stated.
\end{proof}

\begin{proposition}[The breeder is integral]\label{prop:bdeint}
Keep the notation of Proposition~\textup{\ref{prop:bde}} and let $G\ne0$. Then
\[ (Gr-c)(Gp-c)=K \iff G\,rp-c(r+p)=N^{2}, \]
and every solution of that equation has
\[ N\ \big|\ a', \qquad q=\frac{M_0rp-N}{N'}=\frac{a'}{N}\ \in\ \mathbb{Z}. \]
\end{proposition}
\begin{proof}
Expanding, $(Gr-c)(Gp-c)-c^{2}=G\bigl(G\,rp-c(r+p)\bigr)$ while $K-c^{2}=GN^{2}$, so the two
equations differ by the factor $G\ne0$. Substituting $G=NM_0-N'M_0'$, $c=M_0N'$, $a=M_0rp$ and
$a'=M_0'rp+M_0(r+p)$ into $G\,rp-c(r+p)=N^{2}$ gives
$NM_0rp-N'M_0'rp-M_0N'(r+p)=N^{2}$, that is
\[ N\,(a-N)\;=\;N'\bigl(M_0'rp+M_0(r+p)\bigr)\;=\;N'a'. \]
Hence $N\mid N'a'$, and $N$ squarefree gives $\gcd(N,N')=1$, so $N\mid a'$ and $q=a'/N$.
\end{proof}

This is worth isolating because the displayed $q=(M_0rp-N)/N'$ is a quotient by a quantity of the
size of the barrier, and one might expect a candidate to be usable only when $N'$ happens to divide
$M_0rp-N$ --- a further demand of probability about $1/N'$, which would leave the construction
hopeless for reasons having nothing to do with primality. There is no such demand. The expectation
$\mathbb{E}\approx\bigl(\tau(K)/G\bigr)(\ln r\ln p\ln q)^{-1}$ is therefore the correct count: the
breeder loses candidates only to primality, never to integrality. A scan of $51{,}132$ frames
(\texttt{code/breeder\_integral.gp}) produced $9$ admissible candidates, every one with integral $q$.
Formalised in \texttt{lean/Erdos307/Breeder.lean}.

Each factorisation $K=d\cdot(K/d)$ with $d\equiv-c\ (\mathrm{mod}\ G)$ proposes a candidate
$(r,p)=\bigl((d+c)/G,(K/d+c)/G\bigr)$; a solution is any such candidate whose $r,p,q$ are all prime.
A dual, partition--side form is occasionally useful:

\begin{proposition}[Subset form]\label{prop:subset}
For a support $U$ with $D=\prod_{p\in U}p$ and $\Sigma=\sum_{p\in U}D/p$, a partition $U=P\sqcup Q$
solves \#307 iff $S=P$ satisfies $A(\Sigma-A)=D^2$, where $A:=\sum_{p\in S}D/p$. Moreover
$A(\Sigma-A)-D^2=-(1-s_Ps_Q)\,D^2$.
\end{proposition}

(Both identities verified symbolically and on $2\times10^4$ random supports.) The last identity says
the ``analytic'' residual $1-s_Ps_Q$ and the ``arithmetic'' residual $A(\Sigma-A)-D^2$ are one
equation up to the factor $D^2$: the two obvious steering handles carry no independent information.

\paragraph{A conditional reduction.} Several ingredients of a construction are unconditional. Greedy
Egyptian--fraction choice gives disjoint $P,Q$ with $0<1-s_Ps_Q=O(1/T)$ (solvability to any
precision); a prime $r$ nearest $M_0N'/\Delta_0$ gives $G\le\Delta_0^{0.475}(M_0N')^{0.525}$
unconditionally (Baker--Harman--Pintz \cite{BHP2001}), and each small prime $\ell$ can be forced to
divide $K$ by one congruence on a free prime slot of $N$ (Dirichlet then supplies the prime).
Granting these, \#307 reduces to a single prime--density statement,
\[
(\mathrm{H})\qquad\text{among the admissible candidates so produced, some triple $(r,p,q)$ is prime,}
\]
and $(\mathrm{H})\Rightarrow$ \#307 is \textsc{yes}. This $(\mathrm{H})$ is a
Hardy--Littlewood/Bateman--Horn--type assertion for an explicit, computable, \emph{non--polynomial}
family; we know of no named conjecture implying it. That is not an accident of our construction:

\begin{proposition}[No polynomial families]\label{prop:nopoly}
Let $p_1(t),\dots,p_k(t)$, $q_1(t),\dots,q_l(t)$ be integer polynomials, each taking \emph{positive}
prime values on a common infinite branch $t\to+\infty$ (a fixed--shape Schinzel--type family), such
that
$\bigl(\sum_i1/p_i(t)\bigr)\bigl(\sum_j1/q_j(t)\bigr)=1$ for infinitely many $t$. Then every $p_i$ and
$q_j$ is constant: the family is a single fixed solution. Consequently no polynomial parametrisation can
reduce \#307 to Bunyakovsky, Schinzel's Hypothesis~H, or Bateman--Horn.

\textup{Formalised.} \texttt{Erdos307.const\_prod\_eq\_one\_of\_tendsto} is the limit step,
\texttt{eps\_eq\_zero\_of\_identity} the vanishing step, and \texttt{nonconstant\_empty} the
conclusion that both index sets of nonconstant members are empty; \texttt{nopoly} chains them. The
polynomial input enters only as the two hypotheses it supplies, that each nonconstant member
contributes a strictly positive reciprocal and that the remainders tend to $0$, so what is imported
is visible. $0$ \texttt{sorry}, three standard axioms
\textup{(\texttt{lean/Erdos307/NoPoly.lean})}.
\end{proposition}
\begin{proof}
Holding at infinitely many $t$, the relation holds identically as rational functions. A nonconstant
$p_i$ taking infinitely many prime (hence positive) values has positive leading coefficient on an
infinite branch, so $p_i(t)\to+\infty$ there; split each side into constants and nonconstant members,
$\sum_i1/p_i=A_P+\varepsilon_P(t)$ with $\varepsilon_P(t)\to0^+$. Letting $t\to\infty$ in the identity
gives $A_PA_Q=1$ (in particular both sides contain constants), and subtracting,
$A_P\varepsilon_Q(t)+A_Q\varepsilon_P(t)+\varepsilon_P(t)\varepsilon_Q(t)=0$ for all large $t$. Every
term on the left is nonnegative, and strictly positive as soon as one nonconstant member exists; hence
there are none.
\end{proof}

\begin{proposition}[No fixed--shape family, of any growth rate]\label{prop:nofamily}
Let $P(t)=\{p_1(t),\dots,p_k(t)\}$ and $Q(t)=\{q_1(t),\dots,q_l(t)\}$ be families of integers of
\emph{any} shape --- no polynomial, algebraic or analytic hypothesis is imposed --- indexed by
$t\in\mathbb{N}$, with $k,l$ fixed. Suppose each member takes positive prime values, each nonconstant
member tends to $+\infty$, and
$\bigl(\sum_i1/p_i(t)\bigr)\bigl(\sum_j1/q_j(t)\bigr)=1$ for \emph{infinitely many} $t$. Then every
member is constant.
\end{proposition}
\begin{proof}
Let $\varphi$ enumerate, in increasing order, the $t$ at which the relation holds; $\varphi$ is
strictly monotone, so $\varphi(n)\to\infty$. Split each side into its constant members, contributing
$A_P$ and $A_Q$, and its nonconstant ones, contributing $\varepsilon_P(t),\varepsilon_Q(t)>0$ which
tend to $0$ because each nonconstant member tends to $+\infty$ and there are finitely many. Along
$\varphi$ the relation holds at every index and the remainders still tend to $0$, so
Proposition~\ref{prop:nopoly}'s two steps apply verbatim: the limit gives $A_PA_Q=1$, and subtracting
leaves $A_P\varepsilon_Q+A_Q\varepsilon_P+\varepsilon_P\varepsilon_Q=0$ with every term nonnegative
and some term strictly positive if a nonconstant member exists.
\end{proof}

The proof of Proposition~\ref{prop:nopoly} uses polynomiality in exactly one place, to pass from
``holding at infinitely many $t$'' to ``holding identically'', and restricting to the subsequence
removes that need. What the argument actually consumes is three hypotheses --- finitely many members,
positive values, nonconstant members growing --- and none of them is about shape.

The strengthening is not cosmetic. Excluding polynomial families rules out a reduction of \#307 to
Bunyakovsky, Schinzel or Bateman--Horn, but it says nothing about \emph{exponential} families, and
those are what the classical technique for problems of this kind actually uses: Th\=abit's rule for
amicable pairs runs on primes $3\cdot2^{n}-1$, $3\cdot2^{n-1}-1$, $9\cdot2^{2n-1}-1$, and Euler's
generalisation is of the same shape. A member $3\cdot2^{n}-1$ is positive and tends to infinity, so
it satisfies the hypotheses above and the family is excluded. Proposition~\ref{prop:nofamily}
therefore closes the whole Th\=abit/Euler programme for \#307, not merely its polynomial shadow: no
fixed--shape rule of any growth rate can generate solutions, and the constructive calculus of
Section~\ref{sec:frame} is not failing for want of a cleverer ansatz.

Two limits on the statement should be read with it. The hypothesis that the families have
\emph{fixed} cardinality is used, in the splitting into constants plus a remainder tending to $0$, and
a family whose support grows with $t$ is not covered; that gap is real, and harmless, since a growing
support is a growing number of simultaneous primality demands rather than a weaker one. And
positivity is a genuine hypothesis, not a consequence, for the reason recorded in
Remark~\ref{rem:nopolyerratum}. Formalised as \texttt{Erdos307.no\_family\_of\_any\_shape} on three
standard axioms, in the same shape as Proposition~\ref{prop:nopoly}: the abstract core is proved, and
the two facts the family supplies --- that each nonconstant member contributes a strictly positive
reciprocal and that the remainders tend to $0$ --- enter as hypotheses rather than being derived, so
what is imported stays visible. The theorem's statement quantifies over arbitrary index types and
arbitrary $f_P:\mathbb{N}\to\iota\to\mathbb{R}$; no polynomial occurs anywhere in it.

\begin{remark}[Erratum: positivity is a hypothesis, not a consequence]\label{rem:nopolyerratum}
Earlier versions asked only that each polynomial take prime values for infinitely many integers $t$,
without requiring the values to be positive on a common branch. As printed that statement is
\emph{false}: take $P=(2,2,t,-t)$ and $Q=(2,2)$. Each member takes prime values for infinitely many
integers, and
\[ \Bigl(\tfrac12+\tfrac12+\tfrac1t-\tfrac1t\Bigr)\Bigl(\tfrac12+\tfrac12\Bigr)=1 \]
identically, with two nonconstant members present. The proof's step ``every term on the left is
nonnegative'' silently assumes positivity, which is exactly the missing hypothesis; with it, as now
stated, the proof is correct and the counterexample is excluded.

The formalisation carries the repaired statement, not the printed one:
\texttt{Erdos307.nonconstant\_empty} takes $0<f_P(i)$ for every nonconstant member as an explicit
hypothesis, so the Lean development was already stating the corrected proposition. This is the one
place in the paper where formalisation caught a missing hypothesis before review did.
\end{remark}

\noindent This upgrades the earlier computational observation (that the breeder family is not
polynomially parametrised) to a theorem for \emph{all} fixed--shape families: the prime--density
hypothesis behind \#307 is irreducibly non--polynomial, which is why no standard conjecture applies.

\begin{proposition}[Near--miss frames are extremal for $G$]\label{prop:nearframe}
Let $a$ be squarefree with $a'$ squarefree, and take the frame $(M_0,N)=(a,a')$, so that
$M_0'=a'$ and $N'=a''$. Then
\[ G\;=\;NM_0-N'M_0'\;=\;a'\,(a-a''), \]
and since $a''\ne a$ unless \#307 is solved, $|a-a''|\ge1$ and therefore $|G|\ge a'$.
\end{proposition}
\begin{proof}
Direct substitution; the inequality is $|a-a''|\ge1$ for the nonzero integer $a-a''$.
\end{proof}

The point is not that the bound is large but that it runs the wrong way. Writing $r=a''/a$ for the
near--miss ratio, $G=aa'(1-r)$, so a frame built from a good near--miss has $1-r$ small and $aa'$
correspondingly large, and the two cancel \emph{exactly}, because $a(1-r)=a-a''$ is an integer and a
nonzero one. Improving the near--miss therefore does not shrink $G$; it drives $a'$ up and $G$ with
it. The witness of Proposition~\ref{prop:nearmiss}, with $|1-r|=1.65\times10^{-13}$, gives
$|G|\ge a'>10^{112}$, against the $\log_{10}G\approx13$ that the Baker--Harman--Pintz route reaches
at the barrier. So the natural reading of ``steer $\Delta_0$ small by steering $s_{M_0}s_N$ to $1$''
is not merely circular in the sense of Section~\ref{sec:frame}: on near--miss frames it is inverted.

Small $|G|$ does occur, but only on small frames: over $1{,}143{,}888$ coprime squarefree frames with
$M_0,N\le2000$ the minimum is $|G|=5$, at $(M_0,N)=(2,3)$. That is the supply--demand conflict in raw
form, and it can be measured. Binning by $|G|$ and taking the best frame in each bin
(\texttt{code/breeder\_supply.gp}) gives
\[
\begin{array}{l|cccccc}
\log_{10}|G| & 0.5 & 1.5 & 2.5 & 3.5 & 4.5 & 5.0\\\hline
\max \tau(K)/|G| & 0.600 & 0.313 & 0.214 & 0.053 & 0.0065 & 0.0030
\end{array}
\]
The ratio is largest on the smallest frames, where it is already below $1$, and decays monotonically
from there. Since a productive frame needs $\tau(K)/|G|$ well above $1$, no frame in the range
supplies one, and the trend is against scale rather than with it.

\paragraph{The descent, and why it cannot be iterated.} Adjoining a prime $w$ to $M_0$ acts on
$G$ \emph{linearly}: with $M_0\mapsto M_0w$ one has $M_0'\mapsto M_0'w+M_0$ and hence
\[ G\ \longmapsto\ wG-M_0N', \qquad\text{so}\qquad |G_{\mathrm{new}}|=|G|\cdot|w-w^{*}|,
   \qquad w^{*}:=\frac{M_0N'}{G}=\frac{\sigma(N)}{\varepsilon}, \]
writing $\varepsilon=1-\sigma(M_0)\sigma(N)=G/(M_0N)$. Since $G=0$ is \#307, this is a descent
towards a solution, and on the $\varepsilon$ scale it is quadratic: $\varepsilon_{\mathrm{new}}
=\varepsilon^{2}\theta/\sigma(N)$ with $\theta=w-w^{*}$. It is a Newton iteration for \#307.

Three facts close it. First, it cannot start below the barrier: $w^{*}=\sigma(N)/\varepsilon\ge2$
needs $\varepsilon\le\sigma(N)/2$, i.e.\ $\sigma(M_0)\sigma(N)$ near $1$, which by AM--GM needs
$\sigma(M_0N)\ge2$ and so $\omega(M_0N)\ge59$. Over $4{,}284{,}896$ coprime squarefree frames with
$M_0\le5000$, $N\le3000$ the smallest $|\varepsilon|$ is $0.470$, and $w^{*}<2$ throughout.

Second, only $w=\lfloor w^{*}\rfloor$ or $\lceil w^{*}\rceil$ can shrink $|G|$, since
$|G_{\mathrm{new}}|<|G|$ demands $|w-w^{*}|<1$. So a step descends only when one of those two
integers is prime, with probability $\approx2/\log w^{*}$.

Third, and decisively, a successful step squares the difficulty. From $w^{*}=M_0N'/G$,
\[ w^{*}_{\mathrm{new}}=\frac{M_0wN'}{G\theta}=w^{*}\cdot\frac{w}{\theta}\ \approx\
   \frac{(w^{*})^{2}}{\theta}, \]
so $\log w^{*}$ more than doubles at every step while $\log|G|$ falls by only $|\log\theta|\approx
0.69$. The gain is linear in the number of steps $k$ and the cost is quadratic in the exponent:
$|G|$ falls like $2^{-k}$, whereas the probability of surviving $k$ steps is
$\prod_{j<k}2/\log w^{*}_j\asymp2^{-k^{2}/2}$. Starting from the greedy frame $N=30$,
$M_0=\prod_{7\le p\le271}p$, where $\log|G|=249$ and $w^{*}=428.2$
(\texttt{code/breeder\_descent.gp}):
\[
\begin{array}{l|cccccc}
\text{paths available} & 10^{2} & 10^{4} & 10^{6} & 10^{8} & 10^{10} & 10^{12}\\\hline
\text{reachable depth }k & 3 & 5 & 6 & 7 & 8 & 8\\
|G|\text{ shrinks by} & 8 & 32 & 64 & 128 & 256 & 256
\end{array}
\]
against the factor $10^{108}$ required to reach $G=0$. Even $10^{12}$ independent descent paths buy
eight steps and a factor $256$. This is the quantitative form of the circularity of
Section~\ref{sec:frame}: the Newton iteration for \#307 exists, is exact, and converges quadratically,
and each step costs a prime whose size is the square of the last one's.

\paragraph{Why $(\mathrm{H})$ is not a finite search.} It is tempting to read $(\mathrm{H})$ as
``run a certificate hunt that is expected to succeed.'' It is not, for a quantitative reason. The
expected number of certificates from one frame is $\mathbb{E}\approx
\bigl(\tau(K)/G\bigr)\cdot(\ln r\,\ln p\,\ln q)^{-1}$, and \emph{both} factors are pinned by the
geometry.
\begin{itemize}[leftmargin=1.4em]
\item \emph{Demand.} $G=\Delta_0=M_0N(1-s_{M_0}s_N)$ exactly, so $G$ is small only in the circular
regime $s_{M_0}s_N\approx1$; over $2\times10^4$ random coprime squarefree frames the ratio
$G/(M_0N)$ has median $1$ and best value $\approx10^{-0.23}$, and the honest Baker--Harman--Pintz
reduction still leaves $\log_{10}G\approx13$ at the barrier.
\item \emph{Supply.} $\tau(K)$ is bounded by the \emph{size} of $K=GN^2+c^2\ (\ge c^2)$, which is an
\emph{output} of the frame, not a free dial. The maximal order of the divisor function
\cite[Thm.~317]{HardyWright} caps $\log_2\tau(K)$ at $\approx36$ for $K\le10^{54}$ (the barrier
geometry) and at $\approx112$ even for $K\le10^{240}$; reaching $\tau(K)=2^{120}$ needs
$K\gtrsim10^{263}$, i.e.\ frame scale far beyond the barrier.
\end{itemize}
Enlarging a frame grows $K\propto N^2$ but gains divisors only at rate $\log2/\log\log K<1$ per
log--unit of $K$, while it grows $G$ at rate $1$ per log--unit of $N$; so $\tau(K)/G$ cannot be
driven up. Optimising $\mathbb{E}$ over all $(M_0,N)$ splits with geometry--consistent $K$ and the
maximal--order $\tau$, and granting the most favourable demand and singular--series constants, yields
$\log_{10}\mathbb{E}\le-4$ at the barrier and \emph{decreasing} with scale; empirically $385$
factorable frames realise a median of $0$ and a maximum of $1$ admissible divisor, against the
$\sim10^{5}$ a single productive frame would need. A self--consistent accounting; capping the
number $w$ of forced prime factors of $K$ by the frame's own residue freedom (counted generously),
sharpens this to $\log_{10}\mathbb{E}\le-8$ at $B=56$, falling to $\approx-260$ at $B=3000$: the
deficit is \emph{scale--invariant}, so no choice of scale, slot count, or budget split reaches
expectation $1$. (The construction is even self--similar: engineering $K$'s algebraic; e.g.\
Gaussian--integer; factorisation to supply divisors reduces to prescribing $N$ and $N'$
simultaneously, an arithmetic--antiderivative inverse problem of the same difficulty class as \#307.)
The breeder thus makes \#307 conditional on $(\mathrm{H})$ but does \emph{not} render it a feasible
finite computation: the supply--demand conflict through $K=GN^2+c^2$ is exactly the circularity of
Section~\ref{sec:frame} in quantitative form. There is also a structural reason the analogy with
amicable pairs was bound to fall short: te~Riele's breeding amplifies a \emph{known seed} pair into
new ones, whereas the barrier (Theorem~\ref{thm:barrier}) forbids any seed below $2.09\times10^{56}$
the method imports a precondition the problem cannot supply.

\section{Computations}\label{sec:computations}

All verdicts use exact integer/rational arithmetic; floating point is used only for pre--screening.
\begin{itemize}[leftmargin=1.4em]
\item \emph{Exhaustive non--existence below $10^8$.} Enumerating all prime sets $P$ with
$\sum1/p>1$ and $\prod P\le10^8$ (forcing $Q$ as the prime support of $N_P$ and testing
$N_Q=D_P$): $1{,}075{,}419$ candidate sets, no solution. Subsumed by Theorem~\ref{thm:barrier} but an
independent sanity check.
\item \emph{Derivative sieve.} No solution of $n''=n$ with $n'\ne n$ and both members $\le
2\times10^7$ (smallest--prime--factor sieve of the derivative).
\item \emph{Function field.} The arithmetic derivative on $\mathbb{F}_q[t]$ has no two--cycles and no
nonzero fixed points, verified exhaustively over $\mathbb{F}_2$ (degree $\le7$), $\mathbb{F}_3$
($\le5$), $\mathbb{F}_5$ ($\le4$), confirming Theorem~\ref{thm:ff}.
\item \emph{Density.} The sieve gives density $0.042029$ to $2\times10^7$ ($0.017623$ among
squarefree); the independent Erd\H{o}s--Wintner convolution over primes $\le10^5$ gives
$\delta=0.04202$, agreeing to four places (Proposition~\ref{prop:density}).
\item \emph{Conditional mass.} Of $43{,}087$ qualifying squarefree $a$ sampled, $b=a'$ is squarefree
$95.0\%$ of the time; on the resulting $40{,}920$ squarefree $b$ the mean reciprocal mass is $0.047$
vs.\ the required $1/s\approx0.934$, with empirical tail probability $\Pr[\sum_{q\mid b}1/q\ge1/s]=
0/40{,}920$.
\item \emph{Frame/ledger sweep.} Over $2.2\times10^5$ factorable small frames: the divisor--trick
identities hold on every instance; the smallest $\Delta_0$ observed is $22$; the equidistribution
$A\approx\tau(K)/\Delta_0$ holds only in aggregate ($\sum A=127$ vs.\ $\sum\tau(K)/\Delta_0=39.8$)
and fails per frame ($A=0$ for $99.95\%$ of frames); total expected cycles $<1$; zero cycles, as the
barrier requires.
\item \emph{$\mu$--Sondow line census.} Enumerating the lines $n'=n\pm d$ (squarefree members coprime to
$d$) finds no adjacent opposite--sign pair at additive distance $d$: exhaustively to $2\times10^7$ for
$|d|\le200$, and along the small--$|\delta|$ genealogy tree out to members of more than a thousand digits.
Every $\delta=\pm1$ member encountered is a known Giuga or primary pseudoperfect number; the closest
opposite--sign approach recorded is $|y-(x+d)|=11$ (at $d=1$, $30$ vs.\ $42$, and $d=37$, $210$ vs.\ $258$),
attained only among small members; the two trees do not come close at height.
\end{itemize}
Code and logs accompany this note.

\section{Formalisation in Lean~4}\label{sec:lean}

The rigidity and barrier results are formalised against mathlib (Lean~4, toolchain
\texttt{v4.30.0}). The following are machine--checked and \texttt{sorry}--free, depending only on the
standard axioms \texttt{propext}, \texttt{Classical.choice}, \texttt{Quot.sound}:
\begin{itemize}[leftmargin=1.4em]
\item Lemma~\ref{thm:rigidity} (\texttt{rigidity\_coprime}) and Theorem~\ref{thm:structure}
(\texttt{solution\_structure});
\item Lemma~\ref{lem:amgm} (\texttt{reciprocal\_sum\_gt\_two}); the exact thresholds
$\sum_{\text{first }58}1/p<2<\sum_{\text{first }59}1/p$ grounding $|U|\ge59$
(\texttt{recipSum58\_lt\_two}, \texttt{recipSum59\_gt\_two}); the numeric core
$\Pi_{59}^2\ge(4\times10^{112})N_{59}$ (\texttt{barrier\_numeric}), and the algebraic core of the
barrier (\texttt{barrier}), giving $D_P^2\ge4\times10^{112}$;
\item Theorem~\ref{thm:halflyap} (\texttt{log\_lt\_half\_sub}, \texttt{delta\_decreasing},
\texttt{two\_le\_add\_of\_mul\_eq\_one}) and Theorem~\ref{thm:lyapfalse}
(\texttt{criterion\_fails\_of\_pos}, \texttt{criterion\_fails\_of\_nonpos},
\texttt{no\_lambda\_works}, \texttt{lyapunov\_criterion\_false}), neither of which uses
\texttt{native\_decide};
\item Proposition~\ref{prop:masssumthreshold} (\texttt{half\_works}, \texttt{no\_f\_above\_two},
\texttt{mass\_sum\_threshold}) and Proposition~\ref{prop:lyapuniversal}
(\texttt{exists\_f\_of\_injective}, \texttt{increment\_iff},
\texttt{no\_two\_cycle\_of\_decreasing}), which together machine-check the closure of the sign
branch in Theorem~\ref{thm:noinvariant};
\item Proposition~\ref{prop:infzero} (\texttt{exists\_block\_sum\_near}, \texttt{infimum\_zero});
\item Theorem~\ref{thm:noinvariant}(1) in its mechanism (\texttt{Erdos307.Congruential}) and
Proposition~15.28 of the full version (\texttt{live\_slot\_threshold}, \texttt{threshold\_bracket}).
\end{itemize}
The single non--elementary ingredient, that the $k$ smallest primes minimise $\prod$ and maximise
$\sum1/p$ among $k$--element prime sets, is isolated as an explicit hypothesis of \texttt{barrier}
(the ratio bound $R/T(U)\ge\Pi_{59}/T_{59}$), so the dependency is transparent rather than hidden.
It is no longer an open dependency: \texttt{Erdos307.Extremal} proves it
(\texttt{prod\_first\_primes\_le}, \texttt{recipSum\_le\_first\_primes},
\texttt{hRatio\_of\_extremal}), which is why Theorem~\ref{thm:barrier} appears above in the
hypothesis-free form \texttt{erdos307\_barrier\_closed}.

Two remarks on the trust boundary, since it moved at v1.4.6. The numeral bridge
$\texttt{np0}\dots\texttt{np59}$ is checked by the Lean kernel through \texttt{decide}, not by
\texttt{native\_decide}, so \texttt{card\_ge\_59} and \texttt{erdos307\_barrier\_closed} carry only
the three standard axioms. The whole development contains exactly one \texttt{native\_decide} site,
\texttt{dfs\_run}, the pruned-search execution behind Proposition~\ref{prop:close59}; it is the only
non-logical input anywhere, and \texttt{erdos307\_sixty} is the only theorem that carries it.
Proposition~\ref{prop:close59}'s product bounds are proved in this paper and not in Lean:
\texttt{erdos307\_sixty} establishes $|P\cup Q|\ge60$ alone.

\section{Status and questions}\label{sec:status}

Erd\H{o}s~\#307 remains \textbf{open}. It is Conjecture~4 of Ufnarovski--\AA hlander
\cite{UfnarovskiAhlander2003}, so it is equivalent to the existence of a non--trivial two--cycle of
the arithmetic derivative and has been open in that literature since $2003$ under another name. We
establish rigidity and a barrier, $|P\cup Q|\ge60$ with $\min(\prod P,\prod Q)>3.50\times10^{57}$,
which voids all direct search; it appears to be the first valid unconditional bound of its kind,
the only prior candidate resting on an invalid step (Section~\ref{sec:prior}). The rigidity result
and the barrier are Lean--verified, extremality proved rather than assumed.

The barrier is one instance of a single inequality. Writing $n(c)=\min\{n:T_n\ge c\}$,
\[ \sum_{p\in U}1/p\ \ge\ c\ \Longrightarrow\ |U|\ge n(c),\qquad \prod_{p\in U}p\ \ge\ \Pi_{n(c)}, \]
and the arithmetic of each result enters only in producing its $c$ (Remark~\ref{rem:dictionary}).
Read in the size it gives the double--logarithmic frontier law
(Proposition~\ref{prop:frontier}); in the near--miss ratio, the cost family
(Corollary~\ref{cor:cost}), whose endpoint at $r=1$ is the barrier and which shows $a''\ge2a$ needs
$37{,}963$ primes; in the period, that the barrier is uniform in $k$, so it bounds the whole of
Ufnarovski--\AA hlander's Conjecture~3 and not only \#307 (Proposition~\ref{prop:kuniform}); in the
orbit length, that the derivative cannot outgrow its own support
(Proposition~\ref{prop:growth}); in the split, that $\min(\omega(a),\omega(b))\ge3$ whenever
$|U|\le68$ (Proposition~14.20 of the full version); and in the level, the forced core and the ceiling that
pin level $59$ from both sides and yield its admissible count $49{,}961$ in closed characterisation
(Propositions~\ref{prop:window},~\ref{prop:band}, Corollary~\ref{cor:count59}). The algebraic cores
of these are Lean--verified (\texttt{Dictionary.lean}).

On the constructive side we give a Th\=abit/Euler--type calculus and a bilinear breeder whose
candidate $q$ is always integral (Proposition~\ref{prop:bdeint}), a conditional reduction
$(\mathrm{H})\Rightarrow$\textsc{yes}, and a quantitative ledger showing $(\mathrm{H})$ is not a
feasible finite search. No fixed--shape family of any growth rate can parametrise \#307
(Proposition~\ref{prop:nofamily}), which closes the Th\=abit/Euler programme rather than only its
polynomial shadow. We also give a parity dichotomy, an Erd\H{o}s--Wintner density theorem, and a
function--field theorem localising the obstruction to the archimedean place. The heuristics favour
\textsc{yes}, and the model behind them survives a second--moment test
(Proposition~\ref{prop:anticorr}); we claim no proof in either direction.

Open problems: (i) close the residual $3{,}478$ single--tail families at level $60$, and close the
pair sector of Proposition~\ref{prop:sectors}. What this buys should be read with
Proposition~\ref{prop:nottree}: the admissible supports do not form a tree under adjunction, so
completing every one--new--prime family at level $60$ would still not close level $60$; for which Proposition~\ref{prop:pairlocal} rules
out every congruence kill, so only global, per--family input (factorisation or direct search) can
serve (the canonical tail family
is closed, Proposition~\ref{prop:tailkill}; the finite verification of
Proposition~\ref{prop:close59} is Lean--checked, \texttt{erdos307\_sixty}; the conditional $3/4$
law is explained; it is the reciprocity switch of Remark~\ref{rem:fiveeighths}, made exact by
the class--vector functional there); (ii) decide
the parity dichotomy's all--odd case (we expect it empty); (iii) sharpen the certified enclosure
$0.04202\le\delta\le0.04223$ of Proposition~\ref{prop:density}, whose width is carried almost
entirely by the tail window $a$; three digits are certified, and a longer envelope convolution
with a smaller window certifies more; (iv) settle the density hypothesis $(\mathrm{H})$ of
Section~\ref{sec:breeder}, or exhibit a non--circular frame family escaping the $\tau(K)/G$ ceiling;
(v) find a $\mathbb{Z}$--side structure detecting the archimedean obstruction of
Section~\ref{sec:ff}; one that must be non--local, since by Proposition~\ref{prop:nolocal} no single
modulus obstructs; (vi) settle \#307.

\begin{remark}[The place of the problem: a local--global dictionary]\label{rem:place}
The pieces above assemble into a dictionary. \textup{(i)} For a fixed support $U$, the subset form of
Proposition~\ref{prop:subset}, $A(\Sigma-A)=D^2$, satisfies the \emph{Hasse principle as an equation}:
an integer solution $A$ exists iff one exists over $\mathbb{R}$ and over every $\mathbb{Z}_\ell$ (the
discriminant $\Sigma^2-4D^2$ must be a nonnegative square, and squares obey local--global). Its
real--place condition is $\Sigma\ge2D$, i.e.\ $\sigma(U)\ge2$: \emph{the barrier is precisely the
real--place solvability condition of the subset form}. \textup{(ii)} The barrier never needed
primality: pairwise coprime integers have distinct smallest prime factors, so
$\sum 1/a_i\le\sum1/p_i$, and \emph{any} pairwise--coprime support of mass $\ge2$ has at least $59$
elements and product at least $\Pi_{59}$. Coprimality and the ordering of $\mathbb{R}$ carry the wall;
primality enters only through rigidity. \textup{(iii)} What remains unsolved is the \emph{realization}
of the root $A^*=\bigl(\Sigma+\sqrt{\Sigma^2-4D^2}\bigr)/2$ as a sub--multiset sum of the cofactors
$D/p$. Its local conditions are exactly the local anatomy of Lemma~\ref{lem:anatomy}: modulo $\ell^2$
for $\ell\in U$ the subset sums occupy precisely the two cosets
$\ell\mathbb{Z}\cup(D/\ell+\ell\mathbb{Z})$ ($2\ell$ residues, observed exactly ($6$ of $9$, $10$ of
$25$, $14$ of $49$, $22$ of $121$)) and these are the same constraints the equation itself forces:
no completion separates the equation from its realization. In the language of integral points, then:
\#307 is everywhere locally soluble, with the barrier as its archimedean height condition, and its
resolution is a \emph{strong approximation} question for the thin, prime--structured set of subset sums
a failure mode invisible to Brauer--Manin--type obstructions, consistent with every no--go above
and with the existence of cycles over the neighbouring rings. We record this as a dictionary, not a
method.
\end{remark}

\paragraph{Attribution.}
We are careful to separate prior, concurrent, and present contributions (details in
Section~\ref{sec:prior}). \emph{Prior/concurrent:} the two--cycle reformulation, and its
equivalence with the arithmetic derivative, are due to Ufnarovski--\AA hlander
\cite{UfnarovskiAhlander2003} (2003, Conjecture~4); the naming of it as \#307 to Seva
\cite{Seva323194} (2019); an independent rederivation to Bado \cite{Bado2026}; the bound $|P\cup Q|\ge59$ to
J.~Rosen and the disjointness valuation $v_p(\sum1/p_i)=-1$ to R.~Israel \cite{MO320838}; the source
identification to ``Woett'' \cite{MO320838}. Bado \cite{Bado2026} independently obtained the rigidity
cross--matching, the local anatomy, the parity restriction, a one--sided exact test, approximate
solvability, and the union criterion of Proposition~\ref{prop:pyth}. \emph{Present note:} the identification of \#307 with
Ufnarovski--\AA hlander's Conjecture~4 and the resulting two--way bridge to the $n''=n$ literature; the product barrier $2.09\times10^{56}$; the
function--field theorem and archimedean localisation; the constructive deficit analysis and seed
obstruction; the Erd\H{o}s--Wintner density theorem, and the $k$--cycle barriers. Both external
sources \cite{Seva323194,Bado2026} have been examined directly and the overlaps above reflect their
actual contents.

\paragraph{Acknowledgments.}
This work was carried out by the author with substantial assistance from Anthropic's Claude. The
AI assistant contributed to the derivations, the exact and heuristic computations, the Lean~4
formalisation, and the drafting of this note; it was also used to stress-test and attempt to
refute each claim. All statements, proofs, attributions, and the final text were reviewed and are
the responsibility of the author. The author thanks the contributors to the MathOverflow thread
\cite{MO320838}; in particular R.~Israel and J.~Rosen, whose observations are credited above.

\section{Prior and concurrent work}\label{sec:prior}

We record the overlaps with four items, each examined directly.

\paragraph{Ufnarovski--\AA hlander (2003) \cite{UfnarovskiAhlander2003}.}
Their Theorem~5 gives $p^{k}\|n$, $0<k<p$ $\Rightarrow$ $p^{k-1}\|n'$, and their Theorem~4 gives
$n^{(j)}\to\infty$ when $p^{p}\mid n$ properly; together with Corollary~2 ($n$ squarefree
$\Leftrightarrow$ $(n,n')=1$) these force a period--two point to have squarefree members with
disjoint supports. The argument is sound: for $2\le k<p$ two applications of Theorem~5 give
$v_p(n'')=k-2\ne k$, while $k\ge p$ falls to Theorem~4 or to the fixed point $n=p^p$, and the edge
cases $p=2$, $k\in\{2,3\}$ are covered by the latter because $4=p^p$. They then state the
period--two case as \textbf{Conjecture~4}: $\bigl(\sum_{i\le k}1/p_i\bigr)\bigl(\sum_{j\le l}
1/q_j\bigr)=1$ has no solution in distinct primes. This is \#307, and priority for the two--cycle
reformulation belongs here rather than to 2019: what \cite{Seva323194} adds is the identification of
that equation as \#307, which \cite{UfnarovskiAhlander2003} does not make --- they cite
\cite{Barbeau1961} for the derivative, not \cite{Barbeau1976} for the problem. The consequence is
that \#307 has been an open conjecture in the arithmetic--derivative literature since 2003 under
another name, and that the bounds of Section~\ref{sec:barrier} are bounds on Conjecture~4.

\paragraph{Seva (MathOverflow 323194, 2019) \cite{Seva323194}.}
Defines $s(Q)=Q\sum_{p\mid Q}1/p$, derives the valuation property $v_p(s(Q))=v_p(Q)-1$, asks for
$s(P)=Q$, $s(Q)=P$, and states this is Problem~\#307 ``in disguise.'' Priority for \emph{naming}
the two--cycle question as \#307 belongs here. He further asks whether every $s$--trajectory stabilises at $0$; the
$s$--analogue of the Ufnarovski--\AA hlander conjecture, posed independently of it. (Note $s\ne n'$
globally, e.g.\ $s(p^\alpha)=p^{\alpha-1}$, but $s=n'$ on squarefree integers, where both two--cycle
questions live.)

\paragraph{Kovi\v{c} (2012), Propositions~17 and~18 \cite{Kovic2012}.}
Proposition~18 claims $r+s\ge34$ unconditionally ($\ge57$ if $\min\{p_1,q_1\}\ge3$, $\ge110$ if
$\ge5$), together with $p_1<r$, $q_1<s$ and $\max\{m,n\}\ge\prod_{i\le17}p_i\approx1.92\times
10^{21}$. Its proof does not establish them. Substituting $n'=m$, $m'=n$ into its two correct pairs
of bounds gives
\[ \text{(A)}\ \ \frac{rn}{p_r}<m<\frac{rn}{p_1}, \qquad
   \text{(B)}\ \ \frac{sm}{q_s}<n<\frac{sm}{q_1}, \]
after which it asserts ``Hence: $p_1q_s<rs$ and $q_1p_r<rs$'', and deduces everything from these via
$2\max\{p_r,q_s\}<N^2/4$ against $\max\{p_r,q_s\}\ge P_N$. But products of inequalities are valid
only in matching directions, and (A) with (B) admits exactly two: upper$\,\times\,$upper gives
$p_1q_1<rs$, lower$\,\times\,$lower gives $rs<p_rq_s$. The asserted $p_1q_s<rs$ pairs an upper bound
on $m$, carrying $p_1$, with a lower bound on $n$, carrying $q_s$. Nor is it recoverable: the cycle
says precisely $\sigma(n)\in(r/p_r,\,r/p_1)$ and $\sigma(m)=1/\sigma(n)\in(s/q_s,\,s/q_1)$, so
$\sigma(n)$ lies in $(r/p_r,\,r/p_1)\cap(q_1/s,\,q_s/s)$; nonemptiness of that intersection is
equivalent to those two products and to nothing further, whereas $p_1q_s<rs$ reads $q_s/s<r/p_1$, a
comparison of the two \emph{upper} endpoints. In
$n=2\cdot5\cdot13\cdot37\cdot53\cdot73\cdot79\cdot89\cdot97\cdot101$,
$m=3\cdot7\cdot19\cdot31\cdot43\cdot59\cdot67\cdot71\cdot83\cdot103$ --- squarefree, disjoint
supports, $r=s=10$ --- all four inequalities hold while $p_1q_s=206>100=rs$ and $q_1p_r=303>100$
(\texttt{code/kovic\_prop18\_audit.gp}; the valid products and the invalidity are machine--checked in
\texttt{lean/Erdos307/KovicProp18.lean}). Two arithmetic slips sit alongside: Table~2 lists
$P_{11}=27$, and the ``product of the first $17$ primes'' is printed with $27$ for $29$.

What fails is the \emph{inference}, not the \emph{statement}: no two--cycle is known, so no
counterexample to $r+s\ge34$ exists to be given, and none is claimed. The conclusion is unproven
rather than false, and is in fact true, since Theorem~\ref{thm:barrier} gives $|U|\ge60$ by a route
touching none of this. Propositions~16 and~19 are sound, and so is~17, but~17 assumes the
\emph{smaller} member is odd and \cite{Kovic2012} notes that no estimate follows when it is even. The
strongest valid unconditional bound in the prior literature is therefore that paper's computer search
$x>10^4$; Theorem~\ref{thm:barrier} improves it by $52$ orders of magnitude, and $|U|\ge60$ is the
first valid cardinality bound for arithmetic--derivative two--cycles rather than an improvement on an
earlier one. None of this diminishes \cite{Kovic2012}, which remains a standard reference and whose
Propositions~16, 17 and~19 we use.

\paragraph{Bado (preprint, Aix--Marseille, 2026) \cite{Bado2026}.}
Working with $A(S)=\sum_{p\in S}M(S)/p$ and $M(S)=\prod_{p\in S}p$ (our $N_S,D_S$), Bado independently
obtains: the lowest--terms lemma (his Lemma~3.1, our Lemma~\ref{thm:rigidity}); the forcing identities
$A(P)=M(Q)$, $A(Q)=M(P)$ with $P\cap Q=\varnothing$ (his Thm.~A, our Theorem~\ref{thm:structure}); the
two--cycle reformulation $T(T(P))=P$ (his Thm.~B); the valuation signature (his Thm.~6.1, = Israel's);
the local zero--sum congruences (his Prop.~8.1, our Lemma~\ref{lem:anatomy}); the parity restrictions
(his Prop.~8.2, the parity part of Theorem~\ref{thm:parity}); an exact one--sided test (his Thm.~9.1); the deviation parameter, i.e.\ the integer $t$ with
$A(P)=n+tm$, $A(Q)=m-tn$ and solubility exactly at $t=0$ (his Thm.~7.6, our rung $k$ of
Proposition~13.1 of the full version); the intrinsic union criterion (his Thm.~7.3, our
Proposition~\ref{prop:pyth} in \emph{iff} form); uniqueness of the splitting (his Cor.~7.5);
the approximation theorem (his Thm.~10.2, sharper than our greedy version in that all primes exceed a
given bound), and the union discriminant / Pythagorean structure (his Thms.~7.1, 7.4). He cites only
Barbeau, Bloom, and Erd\H{o}s--Graham: crucially, he does \emph{not} identify the map with the
arithmetic derivative, and his note contains none of the present note's $n''=n$ bridge, explicit
barrier constant, function--field theorem, density theorem, $k$--cycle bounds, or constructive
deficit analysis. We regard the overlap on the shared structural layer as strong independent
validation, and import his union criterion with credit. Bado has confirmed this account in
correspondence (August 2026): he agrees that the results listed above were present in his
preprint and that the overlap is concurrent and independent, and he has confirmed the
identification $A(S)=M(S)'$, which he had not made, and is revising his historical discussion
accordingly. The concurrency is therefore recorded on both sides rather than asserted on one:

\begin{proposition}[Single--integer test; derivative form of Bado's Thms.~7.1, 7.4]\label{prop:pyth}
If $(a,b)$ is a two--cycle of the arithmetic derivative, then $N=ab$ is squarefree and satisfies
$N'=a^2+b^2$, whence $N'^2-4N^2=(a^2-b^2)^2$ is a perfect square. Conversely, a squarefree $N$ with
$N'^2-4N^2=\square$ determines a unique unordered factorisation $N=ab$ with $N'=a^2+b^2$, on which the
cycle condition reduces to a single equation (Bado's deviation parameter $t=0$, his Thm.~7.6). This
gives a one--integer necessary test;
computationally, no squarefree $N<2\times10^5$ with $\omega(N)\ge2$ passes it ($0$ of $103{,}596$,
consistent with Theorem~\ref{thm:barrier}, which forces $N>4\times10^{112}$).
\end{proposition}
\begin{proof}
$N=ab$ with $a,b$ coprime squarefree gives $N$ squarefree and, by Leibniz with $a'=b,\,b'=a$,
$N'=a'b+ab'=b^2+a^2$; then $N'^2-4N^2=(a^2+b^2)^2-4(ab)^2=(a^2-b^2)^2$. The converse is Bado's
\cite{Bado2026}; the emptiness of the test below $10^{112}$ is upgraded from a sieve to a theorem in
Corollary~\ref{cor:emptytest}.
\end{proof}

The Pythagorean equation $N'=a^2+b^2$ deserves to be isolated from the (harder) cycle condition.

\begin{proposition}[Pair form; the one--equation relaxation]\label{prop:pairform}
For coprime squarefree $a,b$, the single equation $a^2+b^2=(ab)'$ holds if and only if there is an
integer $k$ with $a'=b+ak$ and $b'=a-bk$; a two--cycle is exactly the case $k=0$. We call a solution
of $a^2+b^2=(ab)'$ (any $k$) a \emph{Pythagorean pair}.
\end{proposition}
\begin{proof}
The Leibniz identity $a^2+b^2-(ab)'=a(a-b')-b(a'-b)$ (a polynomial identity, verified symbolically
and on $2000$ random pairs) makes the equation equivalent to $a(a-b')=b(a'-b)$; since $\gcd(a,b)=1$
this forces $a\mid a'-b$ and $b\mid a-b'$ with equal quotients $k$, i.e.\ $a'=b+ak$, $b'=a-bk$.
Setting $k=0$ recovers $a'=b$, $b'=a$. (Equivalently, in $\mathbb{Z}[i]$: \#307 asks for $z=a+bi$ with
coprime coordinates and $ab$ squarefree satisfying $|z|^2=(\operatorname{Im}z^2/2)'$ and $k=0$.)
\end{proof}

\begin{corollary}[The single--integer test is provably empty below $10^{112}$]\label{cor:emptytest}
Every Pythagorean pair satisfies $\sum_{p\mid a}1/p+\sum_{q\mid b}1/q=a/b+b/a\ge2$ (divide
$a^2+b^2=(ab)'$ by $ab$ and apply AM--GM), so its support carries at least $59$ primes and
$N=ab\ge\prod(\text{first }59\text{ primes})>7.9\times10^{112}$. Hence no squarefree $N<10^{112}$ with
$\omega(N)\ge2$ satisfies $N'^2-4N^2=\square$: the empirical ``$0$ of $103{,}596$'' of
Proposition~\ref{prop:pyth} is a theorem, and the barrier of Theorem~\ref{thm:barrier} extends from
two--cycles to the strictly weaker one--equation (Pythagorean) problem.

\textup{Formalised.} \texttt{Erdos307.pyth\_sum\_of\_squares} is $N'=a^2+b^2$ from Leibniz and the
cycle, \texttt{pyth\_discriminant} the resulting square, \texttt{split\_identity} and
\texttt{split\_product} the two--layer form of Proposition~14.20 of the full version, and
\texttt{mass\_ge\_two\_of\_pythagorean} with \texttt{card\_ge\_59\_of\_pythagorean} this corollary,
glued by \texttt{mass\_ge\_two\_of\_pyth\_support} into \texttt{emptytest}, which takes the
Pythagorean equation itself rather than an assumed mass bound. Earlier editions proved the two ends
and not the link, so this tag overstated what was present.
The AM--GM step of Proposition~\ref{prop:kcycles} is \texttt{sum\_ge\_card\_of\_prod\_eq\_one}, valid
for every $k$ at once and proved from $\log x\le x-1$ rather than a $k$--term AM--GM. Bado's converse
is cited, not reproved \textup{(\texttt{lean/Erdos307/Pythagorean.lean})}.
\end{corollary}

\begin{remark}
(a) As $N'=a^2+b^2$ is a sum of two coprime squares, every odd prime dividing $N'$ is
$\equiv1\pmod4$; a cheap union--side pre--filter complementing Bado's congruence filters. (When
exactly one of $a,b$ is even this is the hypotenuse of the primitive triple $(|a^2-b^2|,2ab,a^2+b^2)$.)
(b) Since Seva's $s$ agrees with $n'$ on squarefree integers, and his valuation property forces
two--cycle members squarefree, his question is \emph{exactly} \#307, not merely an analogue.
(c) Primitive Pythagorean triples form the Barning--Hall ternary tree (the orbit structure of a
reflection group in $O(2,1;\mathbb{Z})$), and a \#307 solution's triple $(|a^2-b^2|,2ab,a^2+b^2)$
occupies one node with generator $(a,b)=(\prod P,\prod Q)$. The tree nonetheless yields no descent on
the problem: its moves $(a,b)\mapsto(2a\mp b,\,a),\,(a+2b,\,b)$ are linear, preserving coprimality but
not squarefreeness (over random squarefree coprime generators only $37\%$ of children keep both
members squarefree) so the locus of \#307--eligible nodes (squarefree generators of disjoint prime
support meeting $a^2+b^2=(ab)'$) is not tree--invariant. The additive tree is transverse to the
multiplicative constraint; the arithmetically natural recursion is instead the prime--append
$\delta$--tree of Proposition~15.33 of the full version, $\delta(Mp)=p\,\delta(M)+M$, which preserves
squarefreeness and tracks the line parameter, which is why the classical members organise on
\emph{that} tree, not Barning--Hall.
\end{remark}

The deviation parameter of Proposition~\ref{prop:pairform} organises the one--equation problem into a
short graded family.

\section{Two closures on the existence route}\label{sec:exclosure}

Atom $\mathrm{A}10$ asks for a prime term of the ternary recurrence of
Remark~\textup{\ref{rem:lehmer}}. Before attacking it, two families of approach can be removed:
the relaxations that would give the construction more freedom, and the algebraic methods that
would avoid the archimedean place. Both are closed outright.

\begin{proposition}[No free slot]\label{prop:nofreeslot}
Let $P,Q$ be finite sets of primes with $\sigma(P)\sigma(Q)=1$, and put $a=\prod P$, $b=\prod Q$.
Then $b=a'$ and $a=b'$ exactly. In particular $Q$ is \emph{determined} by $P$: it is the set of
prime factors of $a'$, and the existence route carries exactly one parameter, the set $P$.

\textup{Formalised.} \texttt{Erdos307.Q\_determined}: given $a'=b$, the set $Q$ is exactly the set of
prime factors of $a'$. The one--prime normal form of Proposition~\ref{prop:oneprime} is
\texttt{oneprime\_forced}, \texttt{oneprime\_cycle}, \texttt{oneprime\_converse}
\textup{(}which closes the equivalence; earlier editions formalised one direction only\textup{)},
\texttt{oneprime\_mass\_identity} and \texttt{oneprime\_mass\_bound}, the last two giving
part~\textup{(b)}: the mass identity holds for
\emph{every} witness, prime or composite, so the exactness layer alone carries the barrier. The slot
calculus of Proposition~15.17 of the full version is \texttt{slot\_layers} and \texttt{slot\_recovery}, the
filter of Corollary~15.16 of the full version is \texttt{qr\_filter}, resting on \texttt{cofactor\_survival},
the congruence $N'+2N\equiv N/\ell\pmod\ell$, which is proved rather than assumed (an earlier edition
stated \texttt{qr\_filter} with that arithmetic as a hypothesis on abstract residues, which left it
contentless); Lemma~15.1 of the full version is \texttt{Frame.csum\_eq\_mul\_recipSum} for its first half
and \texttt{symbolfact} for the Jacobi factorisation; and the reduction mechanism behind
Proposition~\ref{prop:noconj}, Rigidity and Lemma~\ref{lem:anatomy} alike is isolated once as
\texttt{dvd\_sum\_erase\_iff} \textup{(\texttt{lean/Erdos307/NormalForm.lean})}.
\end{proposition}
\begin{proof}
$\sigma(a)=a'/a$ and $\sigma(b)=b'/b$, both in lowest terms since $\gcd(m,m')=1$ for squarefree
$m$. Then $\sigma(a)\sigma(b)=1$ reads $a'b'=ab$. From $\gcd(a',a)=1$ and $a'\mid ab$ we get
$a'\mid b$, say $b=a'k$; symmetrically $b'\mid a$, say $a=b'j$. Substituting,
$a'b'=ab=b'j\cdot a'k$, so $jk=1$ and $j=k=1$.
\end{proof}

\begin{corollary}[The multi-slot relaxations are vacuous]\label{cor:noslack}
There is no version of the construction in which two or more primes of $Q$ may be chosen freely.
In particular one cannot fix a base $B\subset Q$ and solve $1/q_1+1/q_2=r/s$ for the remaining two
primes by ranging over the $\tau(s^2)$ factorisations $(rq_1-s)(rq_2-s)=s^2$: by
Proposition~\ref{prop:nofreeslot} the pair $(q_1,q_2)$ is already forced to be the corresponding
part of the factorisation of $a'$. The exponential parametrisation of
Remark~\textup{\ref{rem:lehmer}} is therefore intrinsic, not an artefact of the one-free-prime
ansatz, and no amount of extra slots converts it into a Bunyakovsky one. What does dispose of the
exponential growth is not extra slots but Proposition~\ref{prop:tailbound}, which bounds the tail
prime outright and so replaces the orbit by a finite divisor search.
\end{corollary}

The second closure is not new here, but it belongs in this list because it is what forces the whole
existence route to be analytic. Theorem~\ref{thm:ff} of Section~\ref{sec:ff} shows that over
$\mathbb{F}_q[t]$ the arithmetic derivative satisfies $\deg f'\le\deg f-1$, hence has no two--cycles
and no nonzero fixed points, and Corollary~\ref{cor:ff307} concludes that the function--field
Erd\H{o}s--Barbeau equation is insoluble outright. \textup{(An independent brute-force recheck,
$8192$ squarefree monics over $\mathbb{F}_2$ to degree $13$ and $6561$ over $\mathbb{F}_3$ to
degree $8$, $0$ violations and $0$ two--cycles: \texttt{code/ff\_nocycle.rs}.)} What is added here
is not that statement but its consequence for \emph{method}, recorded next.

\begin{remark}[Why the existence route cannot be algebraic]\label{rem:archimedean}
The mechanism of Theorem~\ref{thm:ff} is the ultrametric inequality. The cycle equation is
$\sigma(a)\sigma(b)=1$, so one of the two masses must reach $1$; over $K[t]$ the mass satisfies
$|\sigma(f)|=\max_i|1/\pi_i|<1$ by ultrametricity and can never do so, whereas over $\mathbb{Z}$ the
archimedean absolute value permits $\sum_{p\mid n}1/p>1$. So \#307 is possible at all only because
$|\cdot|_\infty$ is not ultrametric.

This has a consequence for method. Any argument for \emph{existence} that is purely algebraic, in
the sense that it applies verbatim with $\mathbb{Z}$ replaced by $K[t]$, would prove a statement
that Theorem~\ref{thm:ff} refutes; hence no such argument exists. Together with
Proposition~\ref{prop:nopoly} (no polynomial parametrisation) and Proposition~13.11 of the full version
(no place of $\mathbb{Q}$ carries the descent) this closes the algebraic and geometric toolkit on
\emph{both} routes: the negative route has no place to descend at, and the positive route has no
place to construct at. What remains for $\mathrm{A}10$ is archimedean and analytic, which is
precisely the regime in which lower bounds for primes in sparse sequences are unavailable. That is
a sharper statement of the obstruction than Remark~\ref{rem:lehmer} gives, and it is the reason the
present work does not resolve the existence direction.
\end{remark}

There is a third closure, and it removes the one concrete attack that survives the other two.
The $34$ immune families of Remark~\textup{\ref{rem:campaign}} are precisely the tail families the
congruence campaign cannot kill, so each is a live candidate and the obvious move is to test them
one at a time. That move is not available, and the reason is quantitative rather than a want of
effort.

\begin{proposition}[The immune families cannot be tested]\label{prop:untestable}
Let $S$ be one of the $34$ immune bases. Deciding whether the tail family $S\cup\{q\}$ contains a
Pythagorean pair is, by Remark~\textup{\ref{rem:lehmer}}, deciding whether
$B_Sx^2-A_Sy^2=-4D_S^2$ has an integral solution with $q=(x^2-D_S)/A_S$ prime. Neither of the two
routes to that decision is within reach.
\begin{enumerate}[label=\textup{(\roman*)},leftmargin=2.2em]
\item \emph{Global route.} Producing a solution requires the fundamental unit, equivalently the
class group, of the real quadratic order of discriminant $4A_SB_S$. On the first immune base that
discriminant has $225$ digits and is a nonsquare, so the subexponential index-calculus cost
$L_\Delta(1/2)=\exp\bigl(\sqrt{\log\Delta\,\log\log\Delta}\bigr)$ evaluates to $10^{24.7}$
operations.
\item \emph{Parametric route.} Since $A_S$ is prime and $(D_S\mid A_S)=+1$, a square root
$x_0^2\equiv D_S\pmod{A_S}$ exists and the family is genuinely parametrised by $x=x_0+tA_S$, with
$q(t)=\bigl((x_0+tA_S)^2-D_S\bigr)/A_S$. But a solution additionally needs $B_Sq(t)+D_S$ to be a
square, and that quantity has $223$ digits already at $t=0$, so the heuristic hit rate is
$10^{-112}$ per trial. Checked at $t=0,1,2,3$: not a square
\textup{(\texttt{code/a10\_pell\_feasibility.gp})}.
\end{enumerate}
There is a third route, which is to decline the computation rather than perform it, and it is worth
costing because it is the natural thing to try next and because nothing structural forbids it. The
size of the first candidate is governed by the regulator: $q_1\approx\varepsilon^2/A_S$ with
$R=\log\varepsilon$, so a base whose regulator happens to be near its minimum
$R_{\min}\approx\tfrac12\log\Delta$ would give $q_1$ of about $10^{113}$, which an Atkin--Morain
test settles in minutes. The obstruction is not that such bases are useless but that they are rare.
Near-minimal regulator means the fundamental solution has $v=1$, that is $\Delta=u^2-4$, and the
$\Delta\le X$ of that shape number about $\sqrt X$; so the fraction of discriminants with a testable
regulator is about $\Delta^{-1/2}=10^{-112.5}$, against a typical $R\approx\sqrt\Delta\approx
10^{113}$ which puts $q_1$ at $\exp(10^{112})$ and out of reach of notation, let alone computation.
Since $\Delta=4A_SB_S$ is determined by the base and not chosen, one may only sample: the $34$
immune bases give an expected $10^{-111}$ hits and all $49{,}961$ give $10^{-108}$, against the
$10^{112.5}$ bases one expected hit would need. The shortfall against the immune family is
$10^{111}$. So base selection does not circumvent the two routes above; the construction is one
fortunate discriminant away from being decidable, and the fortune required is $10^{112.5}$. The
estimate is \textup{HEURISTIC}, being a density count for the regulator.
\end{proposition}

One hope survives all of this on its face: a lower bound for \emph{some term of some sequence in a
family} is formally weaker than one for a single sequence, and Proposition~\ref{prop:close59}
supplies $49{,}961$ of them. It fails, and it fails for a reason no enlargement of the family can
repair.

\begin{proposition}[Multiplicity cannot help]\label{prop:multiplicityfail}
Every known lower-bound method for primes in a set requires that set to have counting function at
least $X^{\theta}$ for some fixed $\theta>0$. The union over all $49{,}961$ bases of the terms of
their sequences below $X$ has counting function $C\log X$ with $C=49{,}961/\log\varepsilon$, since
each sequence grows like $\varepsilon^n$ and so contributes only $O(\log X)$ terms. Hence
\[ \frac{\log(\text{union below }X)}{\log X}=\frac{\log\bigl(C\log X\bigr)}{\log X}\;\longrightarrow\;0 , \]
and the family size $C$ is powerless because it sits inside a logarithm.
\end{proposition}
\begin{proof}[Measurement]
On an immune base $\varepsilon$ has $224$ digits, so at $X=10^{500}$ each sequence contributes
$2.2$ terms and the whole family contributes $10^{5.05}$, a density exponent of $0.0101$. Reaching
$\theta=\tfrac12$ at that $X$ would need $10^{249.7}$ bases against the $10^{4.7}$ available, a
shortfall of $10^{245}$. Even the absolute ceiling on the family at \emph{every} level combined,
all $2^{139}=10^{41.8}$ subsets of the primes below $800$, falls short by $10^{207.8}$
\textup{(\texttt{code/a10\_multiplicity.gp})}.
\end{proof}

Each of the four blocks was then attacked on its own terms. None fell; each is now located more
precisely, and two further routes are closed.

\begin{proposition}[The number--field barrier transfer is \textsc{false}]\label{prop:fieldoptimal}
Theorem~\ref{thm:ff} kills the algebraic route by ultrametricity, which suggests replacing $K[t]$ by
a number ring $\mathcal{O}_K$, where the absolute value is archimedean and the obstruction
disappears. It does disappear. Earlier versions of this paper asserted that nothing is gained, namely
that for every number field $K$ the least $\prod N(\pi)$ over prime ideals with $\sum1/N(\pi)\ge2$ is
at least that of $\mathbb{Q}$. \emph{That assertion is false.} It is withdrawn, and what replaces it
is the following.
\begin{enumerate}[label=\textup{(\roman*)},leftmargin=2.2em]
\item \emph{The true infimum is $16$.} In the biquadratic field
$K=\mathbb{Q}(\sqrt{17},\sqrt{33})$, of discriminant $314{,}721$, the prime $2$ splits completely:
there are four prime ideals of norm $2$. They carry mass $4\cdot\tfrac12=2$ exactly, and their
product of norms is $2^4=16$. Conversely $2$ is the least possible ideal norm, so mass $2$ needs at
least four ideals of norm $2$ and any admissible family has product at least $16$; a field can carry
at most $[K:\mathbb{Q}]$ ideals of norm $2$, so degree $4$ is also minimal. Hence $16$ is the exact
infimum over all number fields, against the claimed $\prod_{i\le59}p_i>8.77\times10^{112}$: the
withdrawn statement was wrong by a factor $10^{112}$.
\item \emph{Where the argument broke.} The Chebotarev computation is correct as a statement about
\emph{density}: in a degree-$d$ Galois field the completely split primes have density $1/d$, so they
carry mass $d\cdot(1/d)\log\log Y=\log\log Y$, exactly as over $\mathbb{Q}$, while each costs
$d\log p$, and inert primes contribute $1/p^{d}$, that is, nothing. The error was to read a
statement about the tail as a statement about the minimum. The barrier consumes only \emph{finitely
many} primes, and for any finite set of small primes one may choose $K$ in which all of them split;
Chebotarev constrains how often splitting happens, not which particular primes split.
\item \emph{What survives, and what closes the route instead.} For a \emph{fixed} field the
measurement stands: $\mathbb{Q}$ needs $59$ primes and $10^{112.9}$, whereas $\mathbb{Q}(i)$ needs
$226$ prime ideals and $10^{735.1}$ \textup{(\texttt{code/a10\_fieldbarrier.gp})}. What does not
survive is the use of this proposition as a \emph{block}. The route is nevertheless closed, and by a
sounder argument, which is Proposition~\ref{prop:distinct} below: the ideal--theoretic \#307 is
\emph{trivially soluble}, so it carries no information about the rational problem.
\end{enumerate}
\end{proposition}

\begin{proposition}[The difficulty of \#307 is distinctness, not mass]\label{prop:distinct}
Allow $P$ and $Q$ to be finite \emph{multisets} of primes rather than sets. Then the \#307 equation
is soluble with four primes in total:
$\bigl(\tfrac12+\tfrac12\bigr)\bigl(\tfrac12+\tfrac12\bigr)=1$. With $P,Q$ genuine sets the same
equation forces $|P\cup Q|\ge59$ and $\prod>10^{112.9}$ by Theorem~\ref{thm:barrier}. The barrier is
therefore a statement about the injectivity of $p\mapsto1/p$ on the primes, not about the arithmetic
of the mass condition.

Two consequences, and the second is the operative one.

\textup{(i)} \emph{The ideal--theoretic problem is trivial.} In $K=\mathbb{Q}(\sqrt{17},\sqrt{33})$
the prime $2$ splits completely into four pairwise coprime prime ideals of norm $2$; taking
$P=\{\pi_1,\pi_2\}$ and $Q=\{\pi_3,\pi_4\}$ gives $\sigma(P)=\sigma(Q)=1$ and $|P\cup Q|=4$
\textup{(\texttt{code/ideal307.gp})}. Nor does the derivative transfer: $\pi_1+\pi_2$ is the unit
ideal, so $a'$ collapses. This is the correct reason the number--field route is closed, and it
explains why refuting the previous proposition reopened nothing: the transferred problem was never
the same problem.

\textup{(ii)} \emph{A no--go for transfers generally.} Any setting admitting two distinct objects of
the same norm $n\ge2$ on each side solves the mass equation at size four, whatever the setting is.
So a dictionary permitting repeated denominators \emph{dissolves} \#307 rather than relocating its
difficulty, and can carry no information about it. This prunes a class of attacks in one line, and
it is the test any future transfer should face first. Corollary~\ref{cor:coprime60} passes it only
because pairwise--coprime integers still have distinct least prime factors, which is distinctness
re-entering through the back door.

\textup{Formalised.} \texttt{Erdos307.multiset\_solution} is the four--element witness,
\texttt{distinctness\_carries\_the\_barrier} the contrast against
\texttt{card\_ge\_59\_of\_recipSum\_ge\_two}, stated as a conjunction so the two halves cannot drift
apart, and \texttt{transfer\_dissolves} the general no--go
\textup{(\texttt{lean/Erdos307/Distinctness.lean})}.
\end{proposition}

\begin{remark}[The three other blocks, attacked]\label{rem:blockattacks}
\textup{(i) Against Proposition~\ref{prop:nofreeslot}.} The rigidity $\prod Q=a'$ would dissolve if
disjointness were an assumption one could relax. It is not: $P\cap Q=\varnothing$ is
\emph{derived} in Theorem~\ref{thm:structure} by the Israel mechanism, $r\in P\cap Q$ forcing
$v_r(st)=-2\neq0$. There is nothing to relax.

\textup{(ii) Against Proposition~\ref{prop:untestable}.} Homogenised, $B_Sx^2-A_Sy^2=-4D_S^2$ is the
ternary form $B_Sx^2-A_Sy^2+4D_S^2z^2$, a \emph{conic}, and conics obey Hasse--Minkowski. Since
Remark~\ref{rem:tier60} already gives local solubility everywhere, \emph{rational} solubility
follows unconditionally, with no class group anywhere in the argument. This sharpens the block
rather than removing it: what is missing is not solubility but \emph{integrality}, and extracting
integral points from the rational parametrisation returns the fundamental unit. Worse, the
constructive step is itself blocked, since Simon's algorithm must factor the determinant
$-4A_SB_SD_S^2$, whose cofactor after removing the known primes of $D_S$ is the $224$-digit balanced
semiprime $A_SB_S$: the same cryptographic wall as Remark~\ref{rem:campaign}, reached from a new
direction.

\textup{(iii) Against Proposition~\ref{prop:multiplicityfail}.} The template that exploits a family
without density is Heath-Brown's: Artin's conjecture holds for all but at most two primes, with no
means of saying which. It does not transfer. That argument needs each failure to impose a
\emph{detectable} structure, which several failures then cannot share. Here a failing family would
be one all of whose $q_n$ are composite, and the only known mechanism for that is a covering system,
which Proposition~\ref{prop:localcomplete} proves does not exist. A failure would therefore be
structureless, and structureless failures cannot be played against one another.
\end{remark}

All of the above bounds the negative direction. The positive direction admits a quantitative
shadow that had not been measured: a solution is exactly $r(a):=\sigma(a)\sigma(a')=1$, so one may
ask how close $r$ actually comes to $1$. The answer is instructive, and it disposes of the usual
way such heuristics are calibrated.

\begin{proposition}[The near-miss spectrum, and why the heuristic cannot be calibrated]\label{prop:nearmiss}
Over all $4{,}385{,}727$ integers $a\le10^7$ with $a$ and $a'$ both squarefree, the maximum of
$r(a)$ is $0.5535$, attained at $a=3{,}972{,}254$; the frontier advances by decade as
$0.333,\,0.345,\,0.502,\,0.502,\,0.536,\,0.554$. In particular $r\ge1$ never occurs, the defect
$|a''-a|$ is of size $a^{0.9}$ or larger in all but three cases, and the smallest relative defect is
$0.4465$. Consequently the constant in the heuristic count of two-cycles, $f(1)\log X$ with $f$ the
density of $r$ at $1$, is \emph{not measurable}, and the obstruction is two--sided.

By Theorem~\textup{\ref{thm:barrier}} the region $r\ge1$ lies above $10^{112.9}$. The window
immediately to the \emph{left} of $1$ is barred as well, by AM--GM rather than by the barrier
hypothesis: $\sigma(a)+\sigma(a')\ge2\sqrt{\sigma(a)\sigma(a')}=2\sqrt r$, and
$2\sqrt r>T_{58}=1.99874\ldots$ as soon as $r>(T_{58}/2)^2=0.998740\ldots$, whence $|P\cup Q|\ge59$
and the product again exceeds $\Pi_{59}>8.77\times10^{112}$. So the whole neighbourhood
$r\in(0.998740,\infty)$ of the value $1$ lies above the barrier, not merely the half--line $r\ge1$,
and no sweep reaches any of it \textup{(\texttt{code/a10\_nearmiss.rs},
\texttt{code/a10\_density1.rs})}. Earlier editions barred only $r\ge1$, which left the left window
formally open and the unmeasurability argument one--sided.
\end{proposition}

\begin{proposition}[An effective approximation theorem, with a certified witness]\label{prop:effapprox}
Bado's Theorem~10.2 \cite{Bado2026} asserts that for every $Y$ and every $\varepsilon>0$ there are
disjoint finite sets of primes $P,Q$, all greater than $Y$, with $0\le\sigma(P)\sigma(Q)-1<
\varepsilon$. The proof is a greedy accumulation and is ineffective. Made effective it is expensive:
to land in $[T,T+\eta)$ the accumulation must start above $\approx1/\eta$, and Mertens gives
$\sum_{Y_1<p\le W}1/p\approx\log\log W-\log\log Y_1=T$, so $W=Y_1^{e^{T}}$; two stages with
$T\approx1$ put the largest prime and the cardinality both at $\approx\varepsilon^{-e^{2}}=
\varepsilon^{-7.389\ldots}$.

A Newton ladder does the same work with a handful of primes. Take $Q=\{2,3,5\}$, so
$\sigma(Q)=31/30$ and $P$ must approach $30/31$; from a pool, repeat $p\leftarrow$ the least prime
$\ge1/d$ and $d\leftarrow d-1/p$, where $d$ is the running deficit. Each rung squares it,
$d_{\mathrm{new}}=d-1/p\approx t\,d^{2}$ with $t$ the local prime gap, so the largest prime is
$\approx\varepsilon^{-1/2}$ and the number of rungs is $O(\log\log(1/\varepsilon))$. Taking the pool
to be the primes $7\le p\le271$ and eleven rungs gives $|P|=66$, $|Q|=3$, largest prime of $2{,}010$
digits, and, in exact rational arithmetic,
\[ \bigl|\sigma(P)\sigma(Q)-1\bigr|\;=\;4.98523\ldots\times10^{-4016}, \]
a rational whose numerator has $129$ digits and whose denominator has $4{,}145$. Every rung prime is
\emph{proved} prime by \textup{APR--CL}, not merely a pseudoprime
\textup{(\texttt{code/nearmiss\_effective.gp})}. The certificate is tiered, and a reader may stop
where they like:
\[
\begin{array}{l|ccc}
\text{rungs} & 9 & 10 & 11\\\hline
\text{largest prime (digits)} & 505 & 1{,}007 & 2{,}010\\
\bigl|\sigma(P)\sigma(Q)-1\bigr| & 9.57\times10^{-1007} & 2\times10^{-2010} &
  4.99\times10^{-4016}\\
\text{cumulative certification} & <10\text{ s} & \approx110\text{ s} & \approx27\text{ min}
\end{array}
\]
\end{proposition}

The gap is not marginal. At $\varepsilon=10^{-13}$ the greedy needs a largest prime near $10^{94}$
and some $10^{92}$ primes; the ladder needs a largest prime near $10^{7}$ and three rungs. The
ladder is exponentially better in the size of the largest prime and doubly exponentially better in
the count, and the reason is visible: greedy convergence is linear in the number of primes, Newton
convergence is quadratic. Truncating gives certificates of any size --- five rungs and a $34$--digit
largest prime already give $6.07\times10^{-66}$, and seven rungs give $3.22\times10^{-254}$. The witness above approaches $1$ from below; taking
the last rung's prime just under $1/d$ instead overshoots, so either side is available and the
statement of Theorem~10.2 as printed is met.

What this is not: progress towards a solution. A free pair of disjoint prime sets carries none of
the rigidity an exact hit requires --- $\sigma(P)\sigma(Q)=1$ forces $\prod P=(\prod Q)'$ and
$\prod Q=(\prod P)'$, and nothing in the ladder produces that --- and by
Proposition~\ref{prop:nearmiss} the size of $|\sigma(P)\sigma(Q)-1|$ does not measure progress
towards the integer identity at all. The value of the construction is that it makes an existence
theorem effective and exhibits a witness anyone can check.

\begin{proposition}[The near--miss frontier obeys a double--logarithmic law]\label{prop:frontier}
For $X\ge2$ put $R(X)=\max\{r(a):a\le X\ \text{squarefree}\}$, where $r(a)=\sigma(a)\sigma(a')$. Let
$\omega_{\max}(X)$ be the largest $k$ with $\Pi_k\le X$, let $\sigma_{\max}(X)=T_{\omega_{\max}(X)}$,
and let $n(X)$ be the largest $n$ with $\Pi_n\le X^{2}\sigma_{\max}(X)$. Then
\[ R(X)\ \le\ \Bigl(\tfrac12\,T_{n(X)}\Bigr)^{2}. \]
\end{proposition}
\begin{proof}
Rigidity gives $\gcd(a,a')=1$, so $U=\operatorname{supp}(a)\sqcup\operatorname{supp}(a')$ has
$|U|=\omega(a)+\omega(a')$ distinct primes and $\prod_{p\in U}p\mid aa'$. By AM--GM,
$\sum_{p\in U}1/p=\sigma(a)+\sigma(a')\ge2\sqrt{\sigma(a)\sigma(a')}=2\sqrt{r}$, while the same sum
is at most $T_{|U|}$ by the extremality of the smallest primes; hence $r\le(T_{|U|}/2)^{2}$. For the
size, $\Pi_{|U|}\le\prod_{p\in U}p\le aa'=a^{2}\sigma(a)\le X^{2}\sigma_{\max}(X)$, so
$|U|\le n(X)$, and $T$ is increasing.
\end{proof}

The bound is unconditional, needs no squarefreeness of $a'$, and is close: against the measured
frontier of Proposition~\ref{prop:nearmiss} it reads
\[
\begin{array}{l|cccccc}
X & 10^{2} & 10^{3} & 10^{4} & 10^{5} & 10^{6} & 10^{7}\\\hline
\text{bound} & 0.4014 & 0.4920 & 0.5296 & 0.5879 & 0.6129 & 0.6342\\
\max r \text{ observed} & 0.333 & 0.345 & 0.502 & 0.502 & 0.536 & 0.554
\end{array}
\]
\textup{(\texttt{code/frontier\_bound.gp})}. Two consequences. First, since
$T_{n(X)}=\log\log\bigl(X^{2}\sigma_{\max}\bigr)+M+o(1)$ by Mertens, the frontier can climb only like
\[ R(X)\ \le\ \tfrac14\bigl(\log\log X^{2}+M+o(1)\bigr)^{2}, \]
a \emph{doubly logarithmic} law: this is why the measured frontier crawls, and why an extrapolation
of its per--decade advance is worthless. Second, the bound reaches $1$ exactly when $T_{n}\ge2$, that
is $n\ge59$, that is $X^{2}\sigma_{\max}\ge\Pi_{59}$, giving $X\ge2.0942\times10^{56}$ --- the
constant of Theorem~\ref{thm:barrier}, recovered here from the frontier side rather than from the
cycle equations. The two derivations are independent and agree to all printed digits.

\begin{corollary}[The cost of a near--miss; the barrier as a function of $r$]\label{cor:cost}
For $r>0$ put $n(r)=\min\{n:T_n\ge2\sqrt r\}$. Every squarefree $a$ with $r(a)\ge r$ has
$|U|\ge n(r)$, hence $aa'\ge\Pi_{n(r)}$ and
\[ \min(a,a')\ \ge\ \sqrt{\Pi_{n(r)}/T_{n(r)}}. \]
\end{corollary}
\begin{proof}
The first display of the proof of Proposition~\ref{prop:frontier} gives $2\sqrt r\le\sum_{p\in U}1/p
\le T_{|U|}$, so $|U|\ge n(r)$; the rest is Theorem~\ref{thm:barrier} verbatim with $59$ replaced by
$n(r)$.
\end{proof}

At $r=1$ this returns $n=59$ and $2.0929\times10^{56}$, so Theorem~\ref{thm:barrier} is the endpoint
of a one--parameter family rather than an isolated statement. The family is what makes the crawl
quantitative \textup{(\texttt{code/nearmiss\_cost.gp})}:
\[
\begin{array}{l|cccccccc}
r & 0.5 & 0.554 & 0.6 & 0.7 & 0.8 & 0.9 & 0.99 & 1\\\hline
n(r) & 8 & 9 & 11 & 16 & 24 & 37 & 56 & 59\\
\min(a,a')\ge & 2.6\!\cdot\!10^{3} & 1.2\!\cdot\!10^{4} & 3.6\!\cdot\!10^{5} & 4.4\!\cdot\!10^{9} &
 1.2\!\cdot\!10^{17} & 4.3\!\cdot\!10^{30} & 4.7\!\cdot\!10^{52} & 2.1\!\cdot\!10^{56}
\end{array}
\]
Each increment in $r$ is paid for exponentially in the product, which is the same double--logarithmic
law seen from the other side: $r$ enters the bound only through $2\sqrt r$, and $T_n$ grows like
$\log\log p_n$. It also explains the shape of the swept data. The record $r=0.5535$ needs
$\min(a,a')\ge1.2\times10^{4}$ and was found at $a=3{,}972{,}254$; the next decade of $r$ costs four
decades of $a$, and by $r=0.7$ the requirement has passed $10^{9}$, beyond any sweep that has been
run.

The restriction $r\le1$ is not used in the proof, and above $1$ the corollary says how expensive it
is for the second derivative to \emph{exceed} its argument:
\[
\begin{array}{l|cccccc}
r & 1 & 1.03944 & 1.25 & 1.5 & 2 & 9/4\\\hline
n(r) & 59 & 71 & 212 & 940 & 37{,}963 & 361{,}139\\
\min(a,a')\ge & 10^{56.3} & 10^{71.3} & 10^{274.8} & 10^{1579.7} & 10^{98275.7} & 10^{1127769.3}
\end{array}
\]
So $a''\ge2a$ forces a union support of $37{,}963$ primes and $\min(a,a')>10^{98275}$: the second
derivative cannot double its argument except at wholly inaccessible size. The witness of
Theorem~\ref{thm:lyapfalse}, with $r=1.03944$, needs $|U|\ge71$ and $\min\ge10^{71.3}$, and its $142$
digits supply both.

\begin{remark}[One inequality, six barriers]\label{rem:dictionary}
Corollary~\ref{cor:cost} and Propositions~\ref{prop:kuniform} and~\ref{prop:twoparam} are the same
statement read at different constants. Every barrier in this note is an instance of
\[ \sum_{p\in U}\frac1p\ \ge\ c \quad\Longrightarrow\quad |U|\ \ge\ n(c),\qquad \prod_{p\in U}p\ \ge\
\Pi_{n(c)}, \]
with $n(c)=\min\{n:T_n\ge c\}$, and the arithmetic enters only in producing $c$:
\[
\begin{array}{l|c|c|l}
\text{configuration} & c & n(c) & \text{recorded as}\\\hline
\text{near--miss ratio }r & 2\sqrt r & \text{varies} & \text{Corollary~\ref{cor:cost}}\\
r=1,\ \text{i.e.\ a two--cycle} & 2 & 59 & \text{Theorem~\ref{thm:barrier}}\\
k\text{--cycle} & k/\lfloor k/2\rfloor & 59\ (k\text{ even}) & \text{Proposition~\ref{prop:kuniform}}\\
\text{three--cycle} & 3 & 361{,}139 & \text{Proposition~\ref{prop:kcycles}}\\
\text{all--odd, }\min U\ge3 & 2\ \text{on }p\ge3 & 1{,}412 & \text{Theorem~\ref{thm:parity}}\\
a''\ge2a & 2\sqrt2 & 37{,}963 & \text{above}\\
\text{orbit of growth }G\text{ in }L\text{ steps} & \tfrac{L}{\lceil L/2\rceil}G^{1/L} &
 \text{varies} & \text{Proposition~\ref{prop:growth}}
\end{array}
\]
The coincidence $2\sqrt{9/4}=3=c(3)$ is why $r=9/4$ and the three--cycle carry the identical bound
$361{,}139$: they are the same value of $c$. What looked like six separate barriers is one
mass--to--size dictionary, and the only mathematics specific to each row is the derivation of its
$c$.
\end{remark}

\begin{proposition}[The mass window, and the forced core]\label{prop:window}
Let $(a,b)$ be a two--cycle with $|U|=n$ and let $t_n>1$ solve $t+1/t=T_n$. Then
\[ \frac1{t_n}\ \le\ \sigma(a),\,\sigma(b)\ \le\ t_n, \]
and every prime $p$ with $1/p\ge T_{n+1}-2$ lies in $U$.
\end{proposition}
\begin{proof}
The masses are positive with product $1$ and sum at most $T_n$, so each is between the roots of
$z^2-T_nz+1$. If $p\notin U$ then the $n$ primes of $U$ carry at most $T_{n+1}-1/p$, and
$\sigma(U)>2$ forces $1/p<T_{n+1}-2$.
\end{proof}

Both halves are sharp and both are informative at $n=59$. The window is
$[0.952682,\,1.049668]$: \emph{each member's mass is pinned within five per cent of $1$}. And the
forcing puts every prime up to $167$ in $U$ --- $39$ of the $59$ --- so two thirds of the support is
determined before any search begins. The window also recovers
Proposition~14.20 of the full version: the smaller member needs $T_k\ge1/t_n$, which gives $k\ge3$ for
$n\le68$, $k\ge2$ from $n=69$ and $k\ge1$ at $n=1413$, reproducing that proposition's three
thresholds exactly. The forcing decays with the level:
\[
\begin{array}{l|cccccccc}
n & 59 & 60 & 61 & 62 & 63 & 65 & 68 & 72\\\hline
\text{all }p\le & 167 & 103 & 73 & 61 & 47 & 37 & 23 & 19\\
\text{forced} & 39 & 27 & 21 & 18 & 15 & 12 & 9 & 8
\end{array}
\]
\textup{(\texttt{code/forced\_core.gp})}. This accounts for the size of the level--$59$ enumeration.
The $39$ forced primes carry mass $1.91162\ldots$, so the remaining $20$ must supply more than
$0.08838$ from primes above $167$; the twenty smallest such primes reach exactly that, at $277$. The
free part is therefore a choice of $20$ primes from a narrow band, which is why the admissible count
is $49{,}961$ and not astronomical. At level $60$ the forced core has fallen to $27$ and the free
part to $33$, and by level $68$ --- the top of the window Proposition~14.20 of the full version says a level
search must cross --- only $9$ primes remain forced.

\begin{proposition}[The free band; why exactly one level is finite]\label{prop:band}
Let $U$ be an admissible support with $|U|=n$ and let $q=\max U$. If $T_{n-1}<2$ then
\[ q\ <\ \frac1{2-T_{n-1}}. \]
For $n\le58$ no admissible support exists, since $T_n<2$. For $n=59$ the hypothesis holds,
$T_{58}=1.99874\ldots$, and $q<793.68$, so $q\le787$. For $n\ge60$ it fails, since
$T_{59}=2.00235\ldots>2$, and no bound on $q$ is available.
\end{proposition}
\begin{proof}
The $n-1$ primes of $U$ other than $q$ carry at most $T_{n-1}$, so $\sigma(U)>2$ forces
$1/q>2-T_{n-1}$.
\end{proof}

This is the mechanism of Proposition~\ref{prop:levelbarrier} in its sharpest form. Level--by--level
enumeration is finite at $59$ and infinite from $60$, and the switch is the single fact
\[ T_{58}\ <\ 2\ <\ T_{59}. \]
Below $59$ there is not enough mass to reach $2$ at all; at $59$ every prime is capped at $787$; at
$60$ the first $59$ primes already carry mass $2.00235>2$, so appending \emph{any} prime leaves an
admissible support and the count is infinite at once. There is nothing gradual about the transition,
and it happens at one level.

Together with Proposition~\ref{prop:window} this closes level $59$ from both sides. The support
contains every one of the $39$ primes $\le167$ and is contained in the $138$ primes $\le787$, so the
free part is a choice of $20$ primes from the $99$ candidates in $(167,787]$. Unconstrained that is
$\binom{99}{20}=4.288\times10^{20}$ choices; the mass condition $\sigma(U)>2$ cuts it to $49{,}961$,
a reduction by a factor $8.6\times10^{15}$ \textup{(\texttt{code/free\_band.gp})}. The level--$59$
computation of Proposition~\ref{prop:close59} is therefore not a brute force over an unbounded space
but an exact traversal of a set whose two ends are both pinned by the same inequality.

\begin{corollary}[The level--$59$ count, from the two pins]\label{cor:count59}
The admissible supports at level $59$ are exactly the sets
\[ \{p:p\le167\}\ \cup\ S, \qquad S\subset\{p:167<p\le787\},\quad |S|=20,\quad
   \sum_{p\in S}\frac1p\ >\ 2-T_{39}\,=\,0.0883759429\ldots, \]
and there are $49{,}961$ of them.
\end{corollary}
\begin{proof}
Membership of the core is Proposition~\ref{prop:window}, the ceiling is Proposition~\ref{prop:band},
and the displayed inequality is $\sigma(U)>2$ after subtracting the forced mass
$T_{39}=1.9116240570\ldots$. The count is a $20$--from--$99$ knapsack, evaluated in
\texttt{code/level59\_count.rs}.
\end{proof}

This is not a closed formula --- subset--sum counts do not have one --- but it is a closed
\emph{characterisation}, and it reduces the level to a search of $119{,}509$ nodes in place of the
original enumeration. That the two independent routes agree on $49{,}961$ is the point: the pins are
not merely necessary conditions that happen to hold, they cut out the level exactly, and nothing
outside them is doing any work. The evaluation also reports no subset sum within $10^{-12}$ of the
threshold, so the floating comparison is measured to be safe rather than assumed. For orientation,
the baseline $20$ smallest candidates carry $0.0907260944\ldots$, exceeding the threshold by
$T_{59}-2=0.0023501514\ldots$, and every admissible $S$ is a perturbation of that baseline within
this budget.

\begin{proposition}[How many rungs the ladder has]\label{prop:rungcount}
Let $(a,b)$ be a Pythagorean pair with $a<b$, $t=b/a$ and $|U|=n$. Positivity of the two masses gives
$-t<k\le0$, so $k$ takes at most $\lceil t\rceil$ values, and $t\le t_n$ where $t_n+1/t_n=T_n$. Hence
the number of rungs available at level $n$ is at most $\lceil t_n\rceil$, and
\[ t_n>m\iff T_n>m+\tfrac1m . \]
In particular the ladder has exactly the two rungs $k\in\{0,-1\}$ at every level $n\le1412$, and the
third rung $k=-2$ opens only at $n=1413$.
\end{proposition}

Two readings. First, the hypothesis $\sigma(ab)<5/2$ under which
Corollary~13.5 of the full version confines $k$ to $\{0,-1\}$ is not a restriction anywhere a search
operates: it is equivalent to $t<2$, hence to $n\le1412$, and at level $59$ one has $t=1.049668$ with
$\lceil t\rceil=2$. The confinement is unconditional throughout the range the rest of this note
inhabits, and the hypothesis only becomes real past level $1412$.

Second, $1413$ is not a new constant. It is the level at which $T_n$ reaches $5/2$, and that is the
same threshold as the all--odd case of Theorem~\ref{thm:parity}: an all--odd cycle needs mass $2$ on
primes $\ge3$, hence $T_n\ge2+\tfrac12$. The level at which the ladder acquires a third rung and the
level at which an all--odd cycle becomes possible are the same level, for the same reason
\textup{(\texttt{code/level\_structure.rs})}. The fourth rung would need $T_n>10/3$, beyond
$\pi(4\times10^{7})$.

\begin{proposition}[The admissible supports are not a tree]\label{prop:nottree}
There are admissible supports at level $60$ with no admissible subset at level $59$. Explicitly,
\[ U\;=\;\{p:p\le271\}\ \cup\ \{809,\,811\} \]
has $|U|=60$ and $\sigma(U)=2.0012091\ldots>2$, while every $59$--element subset has mass below $2$.
\end{proposition}
\begin{proof}
Deleting the largest prime loses the least mass, so if that deletion is inadmissible every deletion
is. Here $\sigma(U)-1/811=1.9999761\ldots<2$.
\end{proof}

Such supports are exactly those of the form $V\cup\{q\}$ with $|V|=59$, $\sigma(V)\le2$,
$q>\max V$ and $q<1/(2-\sigma(V))$, and they are not exceptional. Taking $V=\{p\le271\}\cup\{r\}$
with $r>794$, so that $\sigma(V)<2$, each $r$ admits every prime $q$ in
$\bigl(r,\ 1/(2-\sigma(V))\bigr)$: at $r=797$ that interval reaches $190{,}414$ and contains
$17{,}061$ primes, and summing over $r<4000$ gives $60{,}831$ such supports from this one shape alone
\textup{(\texttt{code/level\_structure.rs})}.

The consequence is for the organisation of a level search, and it is independent of
Proposition~\ref{prop:levelbarrier}. That proposition says level $60$ has infinitely many admissible
supports; this one says that even the part reachable with bounded primes is not exhausted by the
base--plus--tail organisation, because the admissible supports do not form a tree under adjunction.
Closing every one--new--prime family at level $60$, which is question~(i) of
Section~\ref{sec:status}, would therefore still not close level $60$. The searches of
Propositions~\ref{prop:tailkill} and~\ref{prop:close59} are stated over one--new--prime families
throughout and are unaffected; what changes is the reading of what completing them would buy.

\begin{proposition}[The level--minimal structure at any $c$]\label{prop:levelstructure}
Let $U$ satisfy $\sum_{p\in U}1/p\ge c$ with $|U|=n(c)=\min\{n:T_n\ge c\}$. Then every prime $p$ with
$1/p\ge T_{n(c)+1}-c$ lies in $U$, and $\max U<1/\bigl(c-T_{n(c)-1}\bigr)$, the denominator being
positive by minimality of $n(c)$.
\end{proposition}
\begin{proof}
Both are Propositions~\ref{prop:window} and~\ref{prop:band} with $2$ replaced by $c$; neither
argument uses the value.
\end{proof}

The slack $T_{n(c)}-c$ is the total mass a support may waste, and it collapses as $c$ grows, so the
support becomes more forced, not less \textup{(\texttt{code/level\_structure.rs})}:
\[
\begin{array}{l|ccccc}
c & 2 & 9/4 & 7/3 & 5/2 & 3\\\hline
n(c) & 59 & 231 & 400 & 1{,}413 & 361{,}139\\
T_{n(c)}-c & 2.35\!\cdot\!10^{-3} & 2.33\!\cdot\!10^{-4} & 2.44\!\cdot\!10^{-4} &
  2.07\!\cdot\!10^{-5} & 3.09\!\cdot\!10^{-8}\\
\text{forced primes} & 39 & 181 & 259 & 1{,}175 & 314{,}407\\
\text{forced fraction} & 66\% & 78\% & 65\% & 83\% & 87\%\\
\max U< & 793.7 & 2{,}194 & 8{,}284 & 15{,}592 & 6.19\!\cdot\!10^{6}
\end{array}
\]
The row $c=2$ is Propositions~\ref{prop:window} and~\ref{prop:band}, recovered. The row $c=3$ is the
three--cycle of Proposition~\ref{prop:kcycles}: at its minimal level a three--cycle contains
\emph{every} prime up to $4{,}476{,}608$, that is $314{,}407$ of its $361{,}139$, and no prime beyond
$6.19\times10^{6}$. Its support is almost completely determined, with a slack of $3.09\times10^{-8}$
to distribute among the remaining $46{,}732$ places.

The same reading applies with the mass restricted to primes $\ge P$, which is the setting of
Proposition~\ref{prop:twoparam}. At $P=3$, the all--odd case of Theorem~\ref{thm:parity}, the
minimal level is $n=1412$ with slack $2.069\times10^{-5}$; every odd prime up to $9{,}484$ is forced,
$1{,}174$ of the $1{,}412$ or $83\%$, no prime exceeds $15{,}592$, and the product is $10^{5040.1}$.
At $P=5$ the figures are $n=40{,}259$, $24{,}770$ forced, $\max U<1.61\times10^{6}$ and $10^{209582}$,
reproducing Proposition~\ref{prop:twoparam} from the structural side.

That the all--odd case is $83\%$ forced does not make it a feasible search, and it is worth
recording why, since the shape invites the comparison with level $59$. There the free part was $20$
primes from $99$ candidates and the count was $49{,}961$. Here it is $238$ from $643$: the forced
mass is $1.97751561\ldots$, so the free primes must carry more than $0.02248438\ldots$ against a
baseline of $0.02250507\ldots$, an excess of $2.069\times10^{-5}$. A traversal capped at
$4\times10^{8}$ nodes reaches $1.977\times10^{8}$ admissible supports without exhausting the space
\textup{(\texttt{code/level\_structure.rs})}. The slack per free place, not the forced fraction, is
what decides enumerability, and here it is $87$ times larger per place than at level $59$. Question
(ii) of this section does not close this way.

\begin{proposition}[The barrier stratified by the smaller support]\label{prop:split}
Let $(a,b)$ be a two--cycle and put $k=\min(\omega(a),\omega(b))$. Then
\[ k=1\ \Rightarrow\ |U|\ge1413, \qquad k=2\ \Rightarrow\ |U|\ge69, \qquad
   k\ge3\ \Rightarrow\ |U|\ge59, \]
and the last is attained. Consequently every two--cycle with $|U|\le68$ has
$\min(\omega(a),\omega(b))\ge3$.
\end{proposition}
\begin{proof}
Write $T=T_{|U|}$. The supports are disjoint, so $\sigma(a)+\sigma(b)\le T$, and
$\sigma(a)\le T_k$ for the member with $k$ primes. The product $x(T-x)$ is increasing for
$x\le T/2$, so $\sigma(a)\sigma(b)\le m(T-m)$ with $m=\min(T_k,T/2)$, and $\sigma(a)\sigma(b)=1$
forces $m(T-m)\ge1$. For $k\ge3$ one has $T_k\ge T_3=31/30>T/2$ in the relevant range, so $m=T/2$
and the condition is $T\ge2$, i.e.\ $|U|\ge59$. For $k\le2$ one has $T_k<T/2$, so $m=T_k$ and the
condition is $T\ge T_k+1/T_k$: at $k=1$, $T\ge2.5$ and $|U|\ge1413$; at $k=2$, $T\ge2.0333\ldots$ and
$|U|\ge69$.
\end{proof}

The stratification is sharp at each end. At $k=1$ the member is a prime $p$, so $\sigma(b)=p\ge2$ and
$b$ must carry mass $2$ on primes other than $p$, which needs $1412$ of them; at $k=2$ the best case
is $a=6$, whence $\sigma(b)\ge6/5$ on primes $\ge5$, which needs $67$. This contains and extends
Proposition~16 of \cite{Kovic2012}, that no two--cycle has both members a product of two primes: the
present statement rules out \emph{either} member having two prime factors unless $|U|\ge69$, and
either having one unless $|U|\ge1413$. Since Proposition~\ref{prop:close59} gives $|U|\ge60$, the
window $60\le|U|\le68$ that any level--by--level search must cross admits only splits with both
$\omega(a)\ge3$ and $\omega(b)\ge3$ \textup{(\texttt{code/split\_bound.gp})}.

\begin{proposition}[The derivative cannot outgrow its own support]\label{prop:growth}
Let $a_0\to a_1\to\dots\to a_L$ be an orbit segment of the arithmetic derivative with every $a_i$
squarefree, put $G=a_L/a_0$ and let $U=\bigcup_{i<L}\operatorname{supp}(a_i)$. Then
\[ \sum_{p\in U}\frac1p\ \ge\ \frac{L}{\lceil L/2\rceil}\,G^{1/L},
\qquad\text{equivalently}\qquad
G\ \le\ \Bigl(\frac{\lceil L/2\rceil}{L}\,T_{|U|}\Bigr)^{\!L}, \]
and for even $L$ the geometric growth rate satisfies $G^{1/L}\le T_{|U|}/2$.
\end{proposition}
\begin{proof}
Squarefreeness gives $\gcd(a_i,a_i')=1$, so consecutive members are coprime and
$S_p=\{i<L:p\mid a_i\}$ has no two consecutive elements: it is independent in the path on $L$
vertices, whence $|S_p|\le\lceil L/2\rceil$. Then $\sum_{i<L}\sigma(a_i)=\sum_{p\in U}|S_p|/p\le
\lceil L/2\rceil\sum_{p\in U}1/p$, while $\prod_{i<L}\sigma(a_i)=a_L/a_0=G$ gives
$\sum_{i<L}\sigma(a_i)\ge LG^{1/L}$ by AM--GM. The second form inverts the first using
$\sum_{p\in U}1/p\le T_{|U|}$.
\end{proof}

This is the general shape of which the rest of the dictionary is a specialisation: at $L=2$ it reads
$\sum_{p\in U}1/p\ge2\sqrt G$, which is Corollary~\ref{cor:cost} with $G=r$; at $G=1$ it returns
$c=L/\lceil L/2\rceil\to2$, which is Proposition~\ref{prop:kuniform} on the path in place of the
cycle. The statement it adds is dynamical rather than Diophantine, and it bounds the growth branch
of Ufnarovski--\AA hlander's Conjecture~2: an orbit may tend to infinity, but only slowly unless it
accumulates an enormous support \textup{(\texttt{code/growth\_bound.gp})}:
\[
\begin{array}{l|cccccc}
\text{rate }g & 1.05 & 1.10 & 1.25 & 1.50 & 1.75 & 2\\\hline
|U|\ \ge & 97 & 170 & 1{,}413 & 361{,}139 & 4.6\times10^{9} & 4.3\times10^{16}\\
\text{largest prime}\ \ge & 5.4\!\cdot\!10^{2} & 1.0\!\cdot\!10^{3} & 1.2\!\cdot\!10^{4} &
 5.2\!\cdot\!10^{6} & 1.2\!\cdot\!10^{11} & 1.8\!\cdot\!10^{18}
\end{array}
\]
A squarefree orbit that doubles at every step needs some $4.3\times10^{16}$ distinct primes in its
supports and a prime past $1.8\times10^{18}$: sustained rapid growth is not merely unobserved but
impossible on a small support. The hypothesis that every $a_i$ is squarefree is essential and is not
generic along orbits; it is automatic on cycles, by Ufnarovski--\AA hlander, which is why the cycle
rows of the dictionary carry no such caveat.

The region is not unreachable, only unsweepable. The witness $n_0$ of
Theorem~\ref{thm:lyapfalse} has $n_0$ and $n_0'$ both squarefree and
$r(n_0)=\sigma(n_0)\sigma(n_0')=1.03944\ldots$, so $n_0''>n_0$: the value $1$ that a solution must
hit is interior to the range of $r$, not extremal, and no inequality on $r$ alone can settle \#307.
The same construction approaches $1$ from below. Tuning $\sigma(a)$ against the forced
$\sigma(a')=\tfrac{31}{30}+\varepsilon$ gives
\[ a=3359\!\!\prod_{\substack{7\le p\le269\\p\notin\{229,271\}}}\!\!p,\qquad
   a'=2\cdot3\cdot5\cdot383\cdot113717\cdot P_{99}, \]
with $P_{99}$ prime, both $a$ and $a'$ squarefree, $108$ digits, and
\[ r(a)=0.99207796\ldots,\qquad |r(a)-1|=7.9220\times10^{-3}, \]
against the smallest relative defect $0.4465$ over the swept range
\textup{(\texttt{code/nearmiss\_tune.gp})}.

The same tuning with its acceptance window mirrored, $\sigma(a)$ just \emph{above} $30/31$ rather than
just below, reaches $1$ from the other side and closer. With
\[ a=3301\!\!\prod_{\substack{7\le p\le293\\ p\notin\{137,263\}}}\!\!p, \qquad
   a'=2\cdot3\cdot5\cdot P_{117}, \]
where $P_{117}$ is a $117$--digit prime, one has $a$ of $118$ digits and, exactly,
\[ r(a)=1.001014244\ldots,\qquad |r(a)-1|=1.0142\times10^{-3},\qquad a''>a. \]
This is exact rather than a lower bound, since $a'$ factors completely: the cofactor after trial
division is prime, so $\sigma(a')$ is known and not merely bounded below. It improves the record on
both sides at once, against $7.9220\times10^{-3}$ from below and $3.944\times10^{-2}$ from above,
the latter being $r(n_0)$ of Theorem~\ref{thm:lyapfalse}
\textup{(\texttt{code/nearmiss\_above.gp}, re--checked independently in
\texttt{code/nearmiss\_above\_verify.gp})}.

That record is not the natural limit, and seeing why identifies what the near--miss problem actually
is. In every witness of this shape $a'=2\cdot3\cdot5\cdot P$ with $P$ prime, so
$\sigma(a')=\tfrac{31}{30}+P^{-1}$ with $P^{-1}$ of size $10^{-118}$, and therefore
\[ |r(a)-1|=\tfrac{31}{30}\,\bigl|\sigma(a)-\tfrac{30}{31}\bigr| \]
to any precision that matters. The quantity $a'$ has dropped out. What remains is the Diophantine
question of how closely a sum of \emph{distinct prime reciprocals} approaches $30/31$ from above, and
the $10^{-3}$ above is only the granularity of adjusting by a single prime below $4000$.

Closing the deficit in stages removes that ceiling. With $d=\tfrac{30}{31}-\sigma(S)$ for
$S$ the primes in $[7,271]$, take $p_1$ the least prime exceeding $1/d$, then $p_2$ the least
exceeding $1/(d-p_1^{-1})$, and finally $p_3$ the greatest prime below the reciprocal of what
remains; the overshoot is then of order $\mathrm{gap}(p_3)/p_3^{2}$. Each rung squares the precision.
Sweeping the choice at the first two rungs and retaining those that satisfy the forcing congruence
and have $a'=2\cdot3\cdot5\cdot P$ gives, with
$(p_1,p_2,p_3)=(569,1949,15461)$,
\[ a=569\cdot1949\cdot15461\!\!\prod_{7\le p\le271}\!\!p \quad(58\text{ primes},\ 120\text{ digits}),
   \qquad a'=2\cdot3\cdot5\cdot P_{118}, \]
and, exactly,
\[ r(a)=1+6.278713\times10^{-9},\qquad a''>a. \]
This is $1.6\times10^{5}$ times closer to $1$ than the preceding record.

That record is on the wrong side of $1$, and identifying the family says why. Every witness here has
$a'=2\cdot3\cdot5\cdot P$ with $P$ prime, which is exactly the stratum $M=30$ of
Proposition~\textup{\ref{prop:oneprime}}: with $M'=31$ the equation $MN-M'N'=M^2$ reads
$30N-31N'=900$, and the proposition's $N=M+M'p$, $N'=Mp$ become $a=31P+30$, $a'=30P$. Two identities
follow, both exact: $a''=(30P)'=31P+30$, so $r=a''/a=(31P+30)/a$, and a solution of the stratum is
precisely $30a-31a'=900$. Hence
\[ \sigma(a)=\frac{30P}{31P+30}=\frac{30}{31}-\frac{900}{31a}\;<\;\frac{30}{31} \]
for every solution of this stratum, strictly. Approaches to $1$ from \emph{above}, which is what the
$6.28\times10^{-9}$ above and the $3.76\times10^{-9}$ of a wider sweep are, therefore lie on a side
the stratum cannot populate. They stand as records for the spectrum of $r$ and cannot converge on a
solution.

Aimed below instead, where the stratum's solutions must lie, the same ladder gives
\[ a=431\cdot68111\cdot2792213\!\!\prod_{7\le p\le271}\!\!p \quad(58\text{ primes},\ 123\text{ digits}),
   \qquad a'=2\cdot3\cdot5\cdot P_{122}, \]
and, exactly,
\[ r(a)=1-1.654663\times10^{-13},\qquad \sigma(a)<\tfrac{30}{31},\qquad a''<a, \]
which is $4.8\times10^{10}$ times closer to $1$ than the previous record from below. For this witness
$30a-31a'$ is a $112$--digit number where a solution needs $900$, which is the honest measure of the
distance remaining: the ratio is close, the integer is not.

Neither construction is exhausted. Over the same sweeps the smallest \emph{available} deficit was
$2.043\times10^{-18}$ from below and $1.209\times10^{-14}$ from above; what limits the reported values
is the rate at which $a'$ takes the form $2\cdot3\cdot5\cdot P$, not the tuning
\textup{(\texttt{code/nearmiss\_ladder.gp}, \texttt{code/nearmiss\_below.gp}, re--checked in
\texttt{code/nearmiss\_ladder\_verify.gp} and \texttt{code/nearmiss\_below\_verify.gp})}.

We stop here, and the reason is a measurement rather than a budget. Between the two from--below
witnesses the ratio improved by three orders of magnitude while $30a-31a'$ fell from $114$ digits to
$112$, against the $3$ a solution requires. Exhausting the available deficit would gain perhaps five
more digits of the $112$. The near--miss frontier measures $|r-1|$, a ratio; membership of the
stratum is the integer identity $30a-31a'=900$, and the two are not commensurable at this scale.
Driving the ratio further is therefore not evidence about the integer, and the line is closed as a
line of attack. What survives it is structural and is stated above: the family is the $M=30$ stratum
of Proposition~\ref{prop:oneprime}, its solutions satisfy $\sigma(a)<30/31$ strictly, and approaches
from above cannot reach one. Those three facts are machine--checked in
\texttt{lean/Erdos307/StratumM30.lean}.

The constant $30/31$ is not special. In the normal form of Proposition~\ref{prop:oneprime} a solution
of the $M$--stratum is $a=Mp$, $b=N=M+M'p$, so $\sigma(a)=\sigma(M)+1/p$ and, by the identity
$(\sigma(M)+1/p)\sigma(N)=1$ of part~(b), $\sigma(b)=1/\sigma(a)<1/\sigma(M)=M/M'$, with defect
exactly $M^2/(M'(M+M'p))$. Hence \emph{every} stratum is one--sided: for no squarefree $M$ does the
half--space $\sigma(b)\ge M/M'$ contain a solution, and $M=30$ recovers $30/31$. This is a corollary
of Proposition~\ref{prop:oneprime}(b) rather than a new result, but it is the constraint that governs
search direction, so it is machine--checked in \texttt{lean/Erdos307/StratumGeneral.lean}. The
consequence for the near--miss programme is that the mirrored tuning of the previous paragraph is not
merely unlucky at $M=30$; the corresponding tuning is empty in every stratum. Consistently with the two--sided statement above, $a$
lies past the barrier: $118>113$ digits, and $r$ sits inside the barred neighbourhood
$(0.998740,\infty)$.

What remains unmeasurable is the density of $r$ at $1$,
since constructed points do not estimate a density, and that is what the heuristic constant needs.
The frontier reported above is therefore superseded as a record and unchanged as an argument.

The situation is not the usual one. Heuristics of Bateman--Horn or Hardy--Littlewood
type are normally checked against data, and the agreement is what justifies confidence in them. Here
the prediction that solutions exist rests on a constant that cannot be estimated by sampling,
because the barrier and the sampleable range are disjoint. The expected-count heuristic for \#307 is
therefore unverifiable in the usual way.

It is not, however, beyond computation, and the reason is that the constant is the barrier.

\begin{proposition}[The heuristic constant is the barrier's own probability]\label{prop:f1}
Model a random squarefree $a$ by $\mathbb{P}(p\mid a)=1/(p+1)$, and $a'$ by the coprimality model
that Proposition~\textup{\ref{prop:anticorr}} validates, so that $\mathbb{P}(q\mid a')=1/q$ for
$q\nmid a$. Then each prime independently lies in $a$, in $a'$, or in neither, with probabilities
$\tfrac1{p+1},\tfrac1{p+1},1-\tfrac2{p+1}$, and
\[ \sigma(a)+\sigma(a')=\sum_{p\,\mid\,aa'}\frac1p,\qquad \mathbb{P}(p\mid aa')=\frac2{p+1}. \]
Since a two-cycle forces $\sigma(a)\sigma(a')=1$ and hence $\sigma(a)+\sigma(a')\ge2$ by
\textup{AM--GM}, the density $f$ of $r=\sigma(a)\sigma(a')$ is supported in the event
$B:\sigma(a)+\sigma(a')\ge2$, so $f(1)\le C\,\mathbb{P}(B)$ with $C$ the conditional density of $r$
at $1$ given $B$. Numerically
\[ \mathbb{P}(B)=2.74\times10^{-90}. \]
\end{proposition}

The point is the order and not the digits. $\mathbb{P}(B)$ is a \emph{tail} probability of a
one-dimensional sum, not a density, which is what makes it stable: it moves from $4.09$ to $3.01$ to
$2.74\times10^{-90}$ as the grid is refined from $2\times10^4$ to $8\times10^4$ bins and the prime
range from $10^4$ to $10^6$, while a direct attempt at $f(1)$ on a two-dimensional grid moves by four
orders of magnitude per refinement and converges to nothing. Two checks support it. The same
computation with $\mathbb{P}(p\mid n)=1/p$ returns $\mathbb{P}(\sigma(n)\ge1)=0.0420$, reproducing
the measured density of prime-abundant integers to four places; and the single configuration in which
$aa'$ absorbs exactly the first $59$ primes has probability $\prod_{p\le277}2/(p+1)=10^{-96.0}$,
below $\mathbb{P}(B)$ by the factor one expects from the many other prime sets of mass $2$
\textup{(\texttt{code/barrier\_probability.rs})}.

One more thing follows, and it disarms the near-miss spectrum entirely.

\begin{proposition}[The near-miss infimum is zero, so closeness is not evidence]\label{prop:infzero}
\[ \inf\Bigl\{\,\bigl|\sigma(P)\sigma(Q)-1\bigr|\ :\ P,Q\ \text{finite disjoint prime sets}\,\Bigr\}=0, \]
and \#307 asks exactly whether the infimum is attained.
\end{proposition}
\begin{proof}
Split the primes into two infinite classes, say by parity of index. Each class has terms tending to
$0$ and divergent reciprocal sum, so by the classical subseries theorem each carries a subseries
summing to exactly $1$. This produces disjoint \emph{infinite} $P_\infty,Q_\infty$ with
$\sigma(P_\infty)=\sigma(Q_\infty)=1$. Truncating both at height $y$ gives finite disjoint sets with
$\sigma(P_y)\sigma(Q_y)\to1$.
\end{proof}

The analytic content is formalised, and in a slightly stronger form than the proof above uses. For
the infimum one does not need the subseries theorem, only that each mass can be placed in
$(1-\delta,1)$: \texttt{Erdos307.exists\_block\_sum\_near} takes the first consecutive block whose
sum exceeds $T-\varepsilon$ and observes that it overshoots by less than one term, and
\texttt{infimum\_zero} assembles two such blocks into a product within $\varepsilon$ of $1$. Both
carry $0$ \texttt{sorry} and the three standard axioms
\textup{(\texttt{lean/Erdos307/InfZero.lean})}. The arithmetic input, that the primes split into
two infinite classes each with reciprocal terms tending to $0$ and unbounded partial sums, is cited
rather than reproved, as in Theorem~\ref{thm:halflyap}.

The greedy form of that argument reaches $|\sigma(P)\sigma(Q)-1|=6.50\times10^{-9}$ with $144$
primes, against $0.4465$ for the swept frontier of Proposition~\ref{prop:nearmiss}
\textup{(\texttt{code/infimum\_zero.gp})}. So the analytic half of \#307 is empty: the mass equation
can be approximated as closely as one likes, and no near-miss, however good, is evidence that a
solution exists.

What the near-misses do measure is a trade-off, and it is the honest picture of the obstruction. A
solution must kill two defects at once, the mass defect $|\sigma(P)\sigma(Q)-1|$ and the size defect
$\log_{10}\bigl(\prod Q\big/\sigma(P)\prod P\bigr)$, the latter vanishing exactly when
$\prod Q=N_P$, which is Theorem~\ref{thm:structure}. Each is reachable alone and neither
construction moves the other:
\[
\begin{array}{lcc}
\text{construction} & \text{mass defect} & \text{size defect}\\[2pt]
\text{greedy, } Q \text{ free} & 6.50\times10^{-9} & 10^{341.3}\\
Q=\text{primes}(a'),\ \text{\texttt{nearmiss\_tune.gp}} & 7.92\times10^{-3} & 0
\end{array}
\]
Fixing $Q$ as the prime set of $a'$ makes the size defect vanish identically and leaves the mass
defect at $10^{-3}$; choosing $Q$ freely drives the mass defect to $10^{-9}$ and costs $341$ orders
of magnitude in size. \#307 is the demand that both columns read zero simultaneously, and nothing in
this paper moves them together.

\begin{remark}[What the calibrated heuristic says, and it is not encouraging]\label{rem:calibrated}
With $\mathbb{E}(X)\sim f(1)\,\delta\log X$ and $\delta=0.3721$, Proposition~\ref{prop:f1} gives an
expected number of two-cycles below the proved barrier of about $10^{-88}$, and $\mathbb{E}(X)$ does
not reach $1$ until $X$ is of order $10^{4\times10^{89}}$. So the heuristic favours \textsc{yes} only
in the limiting sense that $\log X$ diverges; at every scale that any search or any barrier argument
will reach, it predicts nothing. Two consequences, both negative and both worth stating. A proof of
the negative direction would have to contradict a divergence, and a search for a solution is
hopeless not merely at present but at any scale expressible in the notation of this paper. The
constant is small for a reason that is not incidental: it \emph{is} the barrier, seen as a
probability rather than as a counting bound, which is the third independent appearance of the same
Mertens threshold in this work.

The label is \textup{HEURISTIC} and the model is doing the work. What is claimed is only that the
constant, said above to be unavailable to sampling, is available to computation once the coprimality
model of Proposition~\ref{prop:anticorr} is granted.
\end{remark}

\begin{remark}[What is left, on both sides]\label{rem:whatisleft}
The two directions can now be specified rather than merely called open.

\emph{Negative.} Theorem~\ref{thm:noinvariant} closes the certificate route on both branches, so a
proof must count rather than certify; and by Remark~\ref{rem:whatcounting} the count it must defeat
diverges. A proof of the negative direction would therefore have to contradict a divergent expected
count, which is to say that the direction is not merely unproved but probably false.

\emph{Positive.} Three routes are available in principle, and two are closed quantitatively.
\begin{enumerate}[label=\textup{(\alph*)},leftmargin=2.2em]
\item \emph{Exhibit a witness.} Excluded by Remark~\ref{rem:calibrated}: the first is expected near
$10^{4\times10^{89}}$.
\item \emph{Bound the count below by analytic means.} Excluded by size. At the proved barrier the
main term $f(1)\delta\log X$ lies $10^{200.6}$ below the trivial count $X$, and the saving a method
would need in order to keep the main term above its own error grows like $X$ itself. Sieve and
circle-method errors save a power of $X$ or a power of $\log X$. So an analytic proof here would
have to be \emph{exact}, with error identically zero, rather than sharp.
\item \emph{Construct exactly.} This is the only branch alive, and it exists: the Pell orbit of
Remark~\ref{rem:lehmer}. It is alive but narrower than it was, and the narrowing should not be
mistaken for proximity. Two cautions, both quantitative. The searches of
Remark~\ref{rem:notlehmer} came back empty over $3.99\times10^{11}$ candidates, but empty was the
\emph{predicted} outcome: the expected count was $2.4\times10^{-56}$, so the computation confirms the
heuristic model rather than deciding anything, and a confirmed model is not a decided problem. And by
Proposition~\ref{prop:levelbarrier} those searches cannot be continued upward at all, since the
admissible supports are infinite at every level past $59$. The unconditional exclusions of this note
stop where enumeration stops, which is now.
Proposition~\ref{prop:tailbound} bounds the tail prime by $4N$, so the orbit contributes nothing past
a bounded initial segment and the construction is a finite search over the divisors of $4D$ rather
than a Lehmer primality question; Remark~\ref{rem:notlehmer} then makes that finite search
heuristically empty on all $34$ immune families, by two independent estimates. What survives is the
branch itself, not its expectation: the construction must now either come from a base outside the
immune list or exploit the parametrisation $q=2(N\pm k)/m^2$ directly.
\end{enumerate}

The two exclusions interlock, and the interlock is the real statement. Exactness is what algebra
supplies, and Theorem~\ref{thm:ff} excludes every argument that survives the passage to $K[t]$,
where the conclusion is false. The Pell construction escapes that exclusion for one reason only:
its final input is the primality of a determined term, and primality does not survive to $K[t]$. So a
successful method must be \emph{archimedean and exact at the same time}, a combination which is
unusual but not unprecedented, continued fractions solving Pell being the standard instance, and
Remark~\ref{rem:lehmer} is exactly the analogue here.

The labels differ across the three branches and the difference matters. The exclusions of (a) and
(b) are \textup{HEURISTIC}: both rest on the value of $f(1)$ from Proposition~\ref{prop:f1}, hence on
the coprimality model, and a reader who rejects that model recovers both branches. The interlock
itself is not heuristic: Theorem~\ref{thm:ff} and Proposition~\ref{prop:nopoly} are
\textup{PROVED}, and so is the reduction of branch (c) to a primality question. So what follows is
stated at the strength of its weakest link, which is the model.

What remains of \#307 is therefore one question and not a programme: whether some base admits a
Pythagorean tail prime, together with the sign selection $k=0$, which Remark~\ref{rem:lehmer} shows
is a choice between $\{0,-1\}$ rather than a coincidence of measure zero. By
Proposition~\ref{prop:nopoly} the orbit parametrising the Pythagorean locus is forced to be
exponential rather than polynomial. That was previously read as placing the question in the Lehmer
class rather than the Schinzel class, hence open for every non-degenerate sequence; but by
Proposition~\ref{prop:tailbound} a prime tail satisfies $q\le4N$, so the exponential orbit supplies
no candidate past a bounded initial segment and the question is neither Lehmer nor Schinzel. It is a
finite search over the divisors $m\mid4D$ with $AB+Dm^2$ a square, which
Remark~\ref{rem:notlehmer} makes heuristically empty on all $34$ immune families. The residue is
thus smaller and better located than it was, and still undecided: the bound is proved, the
emptiness is not, and no base outside the immune list has been examined at all. That is why nothing
in this paper decides the problem.
\end{remark}

The measurement does expose a sharp structural fact, which turns out to be exactly the barrier seen
statistically.

\begin{proposition}[The barrier is statistically complete]\label{prop:anticorr}
$\mathbb{E}[\sigma(a')]$ falls by an order of magnitude as $\sigma(a)$ grows, from $0.2035$
unconditionally to $0.0207$ on $\sigma(a)\ge1.3$. The whole of this suppression is coprimality.
Writing the prediction for an integer chosen at random subject to $\gcd(n,a)=1$, namely
$\sum_{p\nmid a}p^{-2}$, the ratio of measured to predicted is
\[ 0.65,\quad 0.82,\quad 0.91,\quad 0.92,\quad 0.95,\quad 1.01 \]
at $\sigma(a)\ge0,\,0.5,\,0.9,\,1.0,\,1.2,\,1.3$: it converges to $1$ precisely in the regime a
solution would inhabit \textup{(\texttt{code/a10\_anticorr.rs})}.
\end{proposition}

The mean is not the moment that $f(1)$ depends on. Since $f(1)$ is an integral of the joint density
of $(\sigma(a),\sigma(a'))$ along $\sigma(a)\sigma(a')=1$, a model could match the conditional mean
and still misstate the density; the next moment is therefore the sharper test. The model makes a
prediction for it too, and a per--$a$ one: given $\operatorname{supp}(a)$, each $q\nmid a$ lands in
$a'$ independently with probability $1/q$, so
\[ \mathbb{E}[\sigma(a')]=\sum_{q\nmid a}\frac1{q^{2}},\qquad
   \operatorname{Var}[\sigma(a')]=\sum_{q\nmid a}\frac1{q^{3}}\Bigl(1-\frac1q\Bigr). \]
Measured against these over all $a\le10^{7}$ with $a,a'$ squarefree
\textup{(\texttt{code/f1\_variance.rs})}, the ratio of measured to predicted variance is
\[ 1.111,\ 0.968,\ 0.894,\ 0.956,\ 0.938,\ 0.860,\ 0.948,\ 1.030 \]
at $\sigma(a)\ge0,\,0.5,\,0.7,\,0.9,\,1.0,\,1.1,\,1.2,\,1.3$, against $0.648,\dots,1.010$ for the
mean. The variance is predicted as well as the mean, it shows no systematic drift, and the two
converge together. The top band rests on $438$ values and is correspondingly noisy, so we claim only
the absence of a trend, not the precise value. What the test was designed to detect --- a model
matching the first moment while its tail degrades, which would inflate $f(1)$ and could in principle
send the true value to $0$ --- does not occur.

So on the sampled range $a'$ is mass-poor for one reason only, that it is coprime to $a$ and $a$ has
absorbed the small primes, and there is no further effect to exploit. This is an independent
confirmation of Proposition~15.25 of the full version: the barrier is not merely near-optimal as a bound, it
already contains the entire statistical structure of the obstruction, and no refinement by
correlation is available.

The range matters, and the statement is about averages. Every $a\le4\times10^8$ with $\sigma(a)>1$
is even, so the measurement never sees an $a$ that omits the small primes. Such $a$ exist, since the
prime reciprocals diverge, and for them $a'$ is free to absorb $2$, $3$ and $5$ and is not mass-poor
at all: Theorem~\ref{thm:lyapfalse} exhibits one with $\sigma(a')=\tfrac{31}{30}$. Proposition~\ref{prop:anticorr}
bounds an expectation on a range, and nothing here bounds a supremum off it.

\subsection*{A Lyapunov candidate for the negative direction}

Theorem~\ref{thm:ff} rules out cycles over $\mathbb{F}_q[t]$ by exhibiting a function that provably
decreases, namely the degree. Proposition~13.11 of the full version shows no \emph{absolute value} on
$\mathbb{Q}$ plays that role. It does not exclude functions that are not absolute values, and the
following is the outcome of a search among such functions. It is a conjecture, not a theorem, and it
is recorded because it is falsifiable and because the mechanism protecting it is one of the paper's
own results.

\begin{definition}\label{def:lyap}
For $\lambda>0$ and squarefree $n$ set $\delta_\lambda(n)=\log n+\lambda\,\sigma(n)$.
\end{definition}

\begin{proposition}[A decreasing $\delta_\lambda$ would settle \#307 negatively]\label{prop:lyapworks}
If for some $\lambda>0$ one has $\delta_\lambda(n')<\delta_\lambda(n)$ for every squarefree $n$ with
$n'$ squarefree, then the arithmetic derivative has no cycle of any length among such $n$, and in
particular \#307 has answer \textsc{no}.
\end{proposition}
\begin{proof}
A cycle would give $\delta_\lambda(n)<\delta_\lambda(n)$.
\end{proof}

Since $n'=n\sigma(n)$, the requirement is the pointwise inequality
\[ \log\sigma(n)\;<\;\lambda\bigl(\sigma(n)-\sigma(n')\bigr). \tag{$\ast$} \]
Around a two-cycle the two instances of $(\ast)$ have increments summing to
$\log(\sigma(a)\sigma(b))=0$, so they cannot both hold: the admissible range of $\lambda$ closes
exactly at a cycle. The content is therefore whether that range is non-empty.

Below the barrier the criterion is not conjectural: it holds with the explicit value
$\lambda=\tfrac12$, by an elementary inequality.

\begin{theorem}[The half-mass Lyapunov function]\label{thm:halflyap}
Put $\delta(n)=\log n+\tfrac12\sigma(n)$. Let $n$ be squarefree with $n'$ squarefree. If
$\sigma(n\,n')<2$ then $\delta(n')<\delta(n)$.
\end{theorem}
\begin{proof}
By Theorem~\ref{thm:structure} $n$ and $n'$ are coprime, and both are squarefree, so
$\sigma(n\,n')=\sigma(n)+\sigma(n')$. The hypothesis therefore gives $\sigma(n')<2-\sigma(n)$, hence
\[ \tfrac12\bigl(\sigma(n)-\sigma(n')\bigr)\;>\;\sigma(n)-1\;\ge\;\log\sigma(n), \]
the last step by $\log x\le x-1$. If $\sigma(n)=1$ the two ends read $0$ and the first inequality is
still strict, since $\sigma(n)+\sigma(n')<2$ forces $\sigma(n')<1$. Now $n'=n\sigma(n)$ gives
$\delta(n')-\delta(n)=\log\sigma(n)-\tfrac12(\sigma(n)-\sigma(n'))<0$.
\end{proof}
Verified over all $7{,}480{,}605$ squarefree $n\le2\times10^7$ with $n'$ squarefree lying in the
region: $0$ violations, least slack $0.242$ \textup{(\texttt{code/invent\_lyapunov.rs})}. The analytic content is formalised:
\texttt{Erdos307.log\_lt\_half\_sub} proves $\log s<(s-t)/2$ for $s>0$ and $s+t<2$,
\texttt{delta\_decreasing} gives the increment form, and
\texttt{two\_le\_add\_of\_mul\_eq\_one} is the \textup{AM--GM} step of
Corollary~\ref{cor:lyapbarrier}; all three carry $0$ \texttt{sorry} and only the three standard
axioms \textup{(\texttt{lean/Erdos307/HalfLyap.lean})}. The arithmetic inputs, that $n$ and $n'$ are
coprime and that $n'=n\sigma(n)$, are Theorem~\ref{thm:structure} and the definition, so the
statement is \textup{VERIFIED} modulo those.

\begin{corollary}[A Lyapunov proof of the barrier]\label{cor:lyapbarrier}
The arithmetic derivative has no cycle, of any length, all of whose consecutive pairs satisfy
$\sigma(n\,n')<2$. In particular a two-cycle $a'=b$, $b'=a$ with $a\neq b$ has
$\sigma(ab)=\sigma(a)+\sigma(b)\ge2\sqrt{\sigma(a)\sigma(b)}=2$, with equality only if
$\sigma(a)=\sigma(b)=1$ and hence $a=b$; so $\sigma(ab)>2$ and therefore $ab\ge\Pi_{59}$.
\end{corollary}
\begin{proof}
A cycle would give $\delta(n)<\delta(n)$ by Theorem~\ref{thm:halflyap}. The second statement is
AM--GM together with the mass estimate of Theorem~\ref{thm:barrier}.
\end{proof}

There are exactly two ways to exploit the fact that increments cancel around a cycle, and the
second is foreclosed. Around a two-cycle the two increments of any function sum to zero, so an
obstruction must pin down either their \emph{sign}, which is the Lyapunov route above, or their
\emph{value in a finite group}: if some $\varphi$ satisfied
$\varphi(n')-\varphi(n)\equiv c\pmod m$ for all $n$ with $2c\not\equiv0$, a two-cycle would give
$0\equiv2c$. No such $\varphi$ exists, and not for want of searching:
Proposition~\ref{prop:localcomplete} proves the cycle system soluble modulo every $m$, so any
obstruction expressible as a congruence is empty a priori. Measured, the increments of
$\omega(n)$, of $\#\{p\mid n:p\equiv3\ (4)\}$, of $n$, and of $n+\omega(n)$ each hit every residue
class modulo $2$, $3$ and $4$ \textup{(\texttt{code/invent\_lyapunov.rs})}. The sign route is
therefore the only one of the two available, which is another form of the thesis that the
obstruction is ordered--archimedean rather than congruential (Remark~\ref{rem:ordered}).

The sign route is available but, as it turns out, not weaker than the problem. Both branches of the
dichotomy can now be closed at once.

\begin{theorem}[Both branches are closed]\label{thm:noinvariant}
Let $S$ be the set of squarefree $n\ge2$ with $n'$ squarefree.
\begin{enumerate}[label=\textup{(\arabic*)},leftmargin=2.2em]
\item \textup{(Finite-group branch.)} \textsc{Conjecture}, downgraded. The claim that there is no
$\varphi$ on $S$ and modulus $m$ with $\varphi(n')-\varphi(n)\equiv c\pmod m$ for all $n$ and
$2c\not\equiv0\pmod m$ is \emph{not} established. Proposition~\ref{prop:localcomplete} settles only
the moduli it covers, $m\le210$, and the general statement is equivalent to the target rather than
independent of it: on a functional graph with no periodic orbit one may build a constant-increment
$\varphi$ for any $(m,c)$ by assigning $\varphi$ along depth, so nonexistence for all $m$ is
equivalent to the existence of a periodic orbit. Branch (1) therefore has exactly the status of
branch (2), which Proposition~\ref{prop:lyapuniversal} already proves for the sign side. What is
proved is the mechanism recorded after the proof.
\item \textup{(Sign branch, not weaker than the target.)} A $\delta$ on $S$ with
$\delta(n')<\delta(n)$ throughout exists if and only if $n\mapsto n'$ has no periodic orbit in $S$;
that condition implies a negative answer to \#307 and is not implied by it.
\end{enumerate}
\end{theorem}
\begin{proof}
(1) is not proved; see the statement. Proposition~\ref{prop:localcomplete} makes the cycle system
soluble modulo every $m\le210$, which does not reach the quantifier the branch needs.
(2) is Proposition~\ref{prop:lyapuniversal}; a two-cycle is a periodic orbit, and periodic orbits of
length other than two are not excluded by a negative answer to \#307.
\end{proof}

The mechanism of branch (1) is formalised. \texttt{Erdos307.two\_nsmul\_eq\_zero\_of\_period\_two}
proves that a constant increment $c$ in any abelian group forces $2c=0$ at a point of period two,
\texttt{const\_increment\_iterate} gives the general $\varphi(D^{[k]}n)=\varphi(n)+kc$,
\texttt{no\_period\_two\_of\_two\_nsmul\_ne\_zero} is the certificate such a $\varphi$ would supply,
and \texttt{zmod\_two\_mul\_eq\_zero\_of\_period\_two} states it over $\mathbb{Z}/m$; all with $0$
\texttt{sorry} and on \texttt{propext} alone
\textup{(\texttt{lean/Erdos307/Congruential.lean})}. What is cited rather than formalised is the
\emph{emptiness} of the class, that no such $\varphi$ with $2c\neq0$ exists for the arithmetic
derivative, which is Proposition~\ref{prop:localcomplete} and is arithmetic. Branch (1) is therefore
\textup{VERIFIED} in its mechanism and \textup{PROVED} in its emptiness, and the two should not be
conflated.

The taxonomy that makes this a closure rather than two isolated facts is the displayed thesis above,
that an invariant whose increments cancel around a cycle can be exploited only through their sign or
their value in a finite group. That thesis is not a theorem, and an argument outside it is not
excluded by anything here.

This is a statement about methods, not about \#307, and it is the natural endpoint of the structural
programme: Proposition~13.11 of the full version removes the absolute values, Theorem~\ref{thm:ff} removes
the arguments that survive to $K[t]$, Proposition~\ref{prop:localcomplete} removes the congruences,
and Theorem~\ref{thm:noinvariant} removes what is left of the pointwise route. A proof of the
negative direction, if there is one, must be global: it must count, and not certify.

\begin{remark}[What a counting proof would have to defeat]\label{rem:whatcounting}
The count it would have to defeat diverges. A two-cycle is exactly $r(a)=1$ for
$r(a)=\sigma(a)\sigma(a')=a''/a$, so if $r$ has a limiting density $f$ on the class of squarefree $a$
with $a'$ squarefree, then $\mathbb{P}(a''=a)=\mathbb{P}\bigl(r\in[1,1+\tfrac1a)\bigr)\sim f(1)/a$ and
\[ \mathbb{E}\bigl(\#\{a\le X:a''=a\}\bigr)\ \sim\ f(1)\,\delta\log X,\qquad \delta=0.3721\ldots, \]
$\delta$ being the measured density of the admissible class \textup{(\texttt{code/region\_shape.rs},
$3{,}721{,}148$ admissible $a\le10^7$)}. This diverges, so the heuristic predicts infinitely many
two-cycles, and any proof of the negative direction must account for a divergent expected count
producing no solutions at all.

The same formula locates them, and shows why the location is not knowable. Setting the expectation to
$1$ gives a first solution near $\exp\bigl(1/(\delta f(1))\bigr)$, which is exponential in $1/f(1)$:
at $f(1)=10^{-2}$ it is $10^{116.7}$, at $f(1)=5\times10^{-3}$ it is $10^{233.4}$, at $f(1)=10^{-3}$
it is $10^{1167}$. Halving the constant squares the answer. Consistency with
Theorem~\ref{thm:barrier} needs $f(1)<0.01034$, which is evidence of the softest kind, since an
expectation near $1$ with no solution observed is an ordinary event and not a contradiction. What
Proposition~\ref{prop:nearmiss} shows is that $f(1)$ is the one quantity here that cannot be sampled
\textup{(\texttt{code/heuristic\_location.gp})}. That $f(1)>0$ is not established either; the
constructed witnesses of Theorem~\ref{thm:lyapfalse} and of the near-miss frontier show only that $r$
takes values on both sides of $1$ and within $8\times10^{-3}$ of it.
\end{remark}

Corollary~\ref{cor:lyapbarrier} recovers the barrier by a route independent of the minimisation
used in Theorem~\ref{thm:barrier}: no extremality argument, no enumeration, only $\log x\le x-1$ and
the coprimality of Theorem~\ref{thm:structure}. It is also stronger in one direction, since it
excludes cycles of every length rather than two-cycles alone. What it does not do is decide
\#307: the region it covers is exactly the region where the barrier already says nothing can
happen, and above $\Pi_{59}$ the hypothesis $\sigma(n\,n')<2$ fails, which is where any solution
must live. The general conjecture is what remains.

Over all $153{,}397{,}754$ such $n\le4\times10^8$, $(\ast)$ holds for every
$\lambda\in(0.275438,\,1.250828)$ and fails for no step
\textup{(\texttt{code/lyap\_battery.rs})}. The interval narrows with the range, of width
$1.654,\,1.267,\,1.237,\,1.116,\,1.092,\,1.004,\,0.975$ at cutoffs $10^4$ through $4\times10^8$;
the binding lower constraint is at $n=\Pi_9=223{,}092{,}870$. Six
candidate functions were tested and five died, among them the control $\delta=\log n$, which fails
first at $n=30$ where $\sigma=1.0333$; a battery that did not kill it would not be testing anything
\textup{(\texttt{code/invent\_lyapunov.rs})}.

That interval is an artefact of the range swept. No $\lambda$ works.

\begin{theorem}[The criterion fails, for every $\lambda$]\label{thm:lyapfalse}
There is no real $\lambda$ for which $(\ast)$ holds at every squarefree $n$ with $n'$ squarefree.
Two witnesses suffice. Let
\[ n_0=\!\!\!\prod_{\substack{7\le p\le373\\ p\notin\{307,\,317,\,359\}}}\!\!\!p, \]
a squarefree integer with $68$ prime factors and $142$ decimal digits. Then $n_0'=2\cdot3\cdot5\cdot C$
with $C$ a prime of $141$ digits, so $n_0'$ is squarefree and $\sigma(n_0')=\tfrac{31}{30}$ exactly,
while $\sigma(n_0)=1.0059067\ldots$. Since $\sigma(n_0)>1$ and $\sigma(n_0')>\sigma(n_0)$,
\[ \log\sigma(n_0)\;>\;0\;>\;\lambda\bigl(\sigma(n_0)-\sigma(n_0')\bigr)\qquad(\lambda>0), \]
so $(\ast)$ fails at $n_0$ for every positive $\lambda$. At $n=30$, with $30'=31$ prime,
$\sigma(30)=\tfrac{31}{30}>1$ and $\sigma(30)-\sigma(31)=\tfrac{931}{930}>0$, so $\log\sigma(30)>0\ge
\lambda(\sigma(30)-\sigma(31))$ for every $\lambda\le0$. Together the two exclude every real
$\lambda$.
\end{theorem}
\begin{proof}
Both statements are finite verifications in exact rational arithmetic. For $n_0$: $\sigma(n_0)$ and
$\sigma(n_0')$ are computed as rationals, $\gcd(n_0,n_0')=1$ and $\sigma(n_0)=n_0'/n_0$ confirm
Theorem~\ref{thm:structure}, the factorisation $n_0'=2\cdot3\cdot5\cdot C$ multiplies back to $n_0'$
and has all exponents $1$, and $C$ carries an Atkin--Morain certificate, generated and validated
\textup{(\texttt{certs/lyap\_refute\_cofactor\_ecpp.txt})}. The comparisons $\sigma(n_0)>1$ and
$\sigma(n_0')>\sigma(n_0)$ are exact, and $\sigma(n_0')=\sigma(n_0)$ is impossible in any case by
Lemma~15.2 of the full version. The case $n=30$ is arithmetic
\textup{(\texttt{code/lyap\_refute\_verify.gp})}.
\end{proof}

The argument is formalised. \texttt{Erdos307.criterion\_fails\_of\_pos} and
\texttt{criterion\_fails\_of\_nonpos} give the two sign clashes,
\texttt{no\_lambda\_works} combines them, and \texttt{lyapunov\_criterion\_false} instantiates the
combination at $\sigma(n_0)$, carried as the exact quotient $n_0'/n_0$, and at the pair
$(\tfrac{31}{30},\tfrac1{31})$ of $n=30$; the two numeric comparisons $1<\sigma(n_0)$ and
$\sigma(n_0)<\tfrac{31}{30}$ are checked on the $142$-digit numerals themselves. All six
declarations carry $0$ \texttt{sorry} and only \texttt{propext}, \texttt{Classical.choice},
\texttt{Quot.sound}, with no \texttt{native\_decide}
\textup{(\texttt{lean/Erdos307/LyapFalse.lean})}. What Lean does not see is the factorisation
$n_0'=2\cdot3\cdot5\cdot C$ and the primality of $C$, which is the Atkin--Morain certificate; the
status is therefore \textup{VERIFIED} modulo that certificate, in the same sense that
Theorem~\ref{thm:halflyap} is \textup{VERIFIED} modulo Theorem~\ref{thm:structure}.

The witness is built, not found. For squarefree $n$ and a prime $w\nmid n$ one has
$n'\equiv n\sum_{p\mid n}p^{-1}\pmod w$, hence
\[ w\mid n'\quad\Longleftrightarrow\quad \sum_{p\mid n}p^{-1}\equiv0 \pmod w. \tag{F} \]
Imposing (F) for $w\in\{2,3,5\}$ forces $\sigma(n')\ge\tfrac{31}{30}$, and building $n$ from primes
$\ge7$ with $\sigma(n)$ just above $1$ gives $\sigma(n')>\sigma(n)$. Reaching mass $1$ without
$2,3,5$ needs primes only to about $360$, which is what keeps $n_0'$ small enough to factor
completely; the squarefreeness of $n_0'$ is therefore verified rather than assumed. Three of the
larger pool primes are dropped to satisfy (F), at negligible cost in mass
\textup{(\texttt{code/lyap\_refute.gp})}.

The sweep could not have seen this. Every $n\le4\times10^8$ with $\sigma(n)>1$ is even, so the
battery measured only $n$ containing $2$, whereas the failure needs $n$ to omit $2$, $3$ and $5$ so
that $n'$ may absorb them. By Theorem~\ref{thm:rhobarrier} below, any failure must have
$n\,n'\ge\Pi_{59}$, and $n_0n_0'$ has $284$ digits: the region was out of reach of any sweep, and
only a construction enters it.

One route to the criterion had looked provable rather than measured, and it fails for the same
reason. Write $\rho=\sup\{\sigma(n')/\sigma(n):\sigma(n)>1\}$, the supremum being over the same
class of $n$. If $\rho\le c<1$ then $\sigma(n)-\sigma(n')\ge(1-c)\sigma(n)$, and since
$\log x/x\le1/e$,
\[ L:=\sup_{\sigma(n)>1}\frac{\log\sigma(n)}{\sigma(n)-\sigma(n')}\;\le\;\frac1{e(1-c)} . \]
Measured over $n\le4\times10^8$, $\rho=0.519746$, which would give $L\le0.766$ against a measured
upper demand of $1.274$; a proof of $\rho<1-1/(eU)$, for $U$ a lower bound on the upper demand,
would then have given $(\ast)$. The measurement is not the truth. Already
\[ a_0=\!\!\prod_{5\le p\le163}\!\!p,\qquad
   a_0'=2\cdot3\cdot173\cdot227\cdot1973\cdot82890547\cdot68126695363\cdot P_{36}, \]
with $P_{36}$ prime, is squarefree with squarefree derivative, $\sigma(a_0)=1.0723027\ldots>1$ and
$\sigma(a_0')=0.8440258\ldots$, so $\rho\ge0.7871152\ldots$; the bound the route would yield is
$1/(e(1-\rho))\ge1.7280$, already above the measured upper demand, so it constrains nothing. The
witness $n_0$ of Theorem~\ref{thm:lyapfalse} gives $\rho\ge1.02726$, so $\rho<1$ is false outright
\textup{(\texttt{code/rho\_verify.gp})}.

The mechanism advanced for $\rho<1$ was that $\sigma(n)>1$ forces $n$ to absorb the small primes
while $n'$ stays coprime to $n$. It does not: the reciprocals of the primes diverge, so $\sigma(n)>1$
is reachable from primes $\ge7$ alone, and then $n'$ is free to take $2$, $3$ and $5$. Both $a_0$ and
$n_0$ are of that shape. The earlier probe of this regime evaluated $\sigma(n')/\sigma(n)$ at the
least odd $n$ with $\sigma(n)>1$, obtaining $0.0312$, and read the regime off that single point;
the quantity to be bounded is a supremum, and the least member of a family does not report it. This
is the error recorded in Remark~15.26 of the full version, in a second instance.

What survives is the location of the failure, and it is exactly where the barrier puts it.

\begin{theorem}[$\rho<1$ below the barrier]\label{thm:rhobarrier}
Let $n$ be squarefree with $n'$ squarefree and $\sigma(n)>1$. If $\sigma(n')\ge\sigma(n)$ then
$n\,n'\ge\Pi_{59}>8.77\times10^{112}$. Consequently $\rho<1$ for every $n$ with
$n\,n'<\Pi_{59}$, and since $\sigma(n)>1$ gives $n'>n$, for every $n<\sqrt{\Pi_{59}}$, that is for
every $n<10^{56.5}$.
\end{theorem}
\begin{proof}
$\sigma(n')\ge\sigma(n)>1$ gives $\sigma(n)+\sigma(n')>2$. By Theorem~\ref{thm:structure} $n$ and
$n'$ are coprime, and both are squarefree, so $\sigma(n\,n')=\sigma(n)+\sigma(n')>2$. A squarefree
integer of mass at least $2$ has at least $59$ prime factors and is therefore at least $\Pi_{59}$,
since $\sigma(\Pi_{58})=1.998740<2\le2.002350=\sigma(\Pi_{59})$; this is the mass estimate
underlying Theorem~\ref{thm:barrier}. Finally $\sigma(n)>1$ gives $n'=n\sigma(n)>n$, so
$n^2<n\,n'$.
\end{proof}

The counting step is formalised. \texttt{Erdos307.card\_ge\_59\_of\_recipSum\_ge\_two} isolates the
mass estimate, that a set of primes of reciprocal sum at least $2$ has at least $59$ elements, and
\texttt{Erdos307.card\_union\_ge\_59\_of\_masses} applies it to two disjoint sets each of mass
exceeding $1$; both carry $0$ \texttt{sorry}
\textup{(\texttt{lean/Erdos307/RhoBarrier.lean})}. Their axiom lists are those of
Theorem~\ref{thm:barrier} itself, the three standard axioms together with the
\texttt{native\_decide} evaluation of the first $58$ primes. The same core proves
\texttt{card\_ge\_59} of the barrier, which assumes $\sigma(P)\sigma(Q)=1$ instead; the two
hypotheses differ and the shared step has been extracted rather than duplicated.

Theorem~\ref{thm:rhobarrier} therefore locates the failure before exhibiting it: $\rho\ge1$ forces
$n\,n'\ge\Pi_{59}$, so no sweep can meet a counterexample, and any counterexample must have $n$ of
at least $57$ digits. The witness $n_0$ has $142$, and $n_0n_0'$ has $284$. The regime in which
$\rho\ge1$ is possible is the regime in which a two-cycle is possible, since both need two coprime
squarefree integers of total mass at least $2$; what Theorem~\ref{thm:lyapfalse} shows is that the
first of those is genuinely non-empty, which says nothing about the second.

\begin{corollary}[The half-mass hypothesis cannot be relaxed]\label{cor:halfsharp}
For every $\lambda>0$ and every $c>2.03925$, the implication
``$\sigma(n\,n')<c$ implies $\delta_\lambda(n')<\delta_\lambda(n)$'' is false. Indeed
$\sigma(n_0\,n_0')=2.039240\ldots$ while
$\delta_\lambda(n_0')-\delta_\lambda(n_0)=\log\sigma(n_0)+\lambda(\sigma(n_0')-\sigma(n_0))>0$ for
every $\lambda>0$.
\end{corollary}

Two is also the exact threshold for the argument of Theorem~\ref{thm:halflyap}, in the following
sense, which concerns the inequality alone and not the arithmetic of $\sigma$.

\begin{proposition}[The mass-sum relaxation stops at $2$]\label{prop:masssumthreshold}
Let $c>0$. There is $f:(0,\infty)\to\mathbb{R}$ with $f(x)-f(y)>\log x$ for all $x,y>0$ satisfying
$x+y<c$ if and only if $c\le2$. For $c=2$ the linear solutions are exactly $f(x)=\tfrac12x+\text{const}$.
\end{proposition}
\begin{proof}
The region $x+y<c$ is symmetric, so a solution gives both $f(x)-f(y)>\log x$ and $f(y)-f(x)>\log y$;
adding, $0>\log(xy)$. Taking $x=y\uparrow c/2$ forces $c^2/4\le1$, that is $c\le2$. For $c\le2$ take
$f(x)=\tfrac12x$: if $x+y<2$ then $\tfrac12(x-y)>x-1\ge\log x$. If $f=\lambda x$ then letting
$y\uparrow2-x$ gives $2\lambda(x-1)\ge\log x$ for all $x\in(0,2)$, and dividing by $x-1$ and letting
$x\to1$ from each side gives $2\lambda\ge1$ and $2\lambda\le1$.
\end{proof}

Both halves are formalised: \texttt{Erdos307.half\_works} for $c\le2$ and
\texttt{no\_f\_above\_two} for $c>2$, the latter by the diagonal point $x=y$ with $1<x<c/2$, where
the right side vanishes and the left is positive; \texttt{mass\_sum\_threshold} states the two
together. The uniqueness of the constant is formalised as well: \texttt{lambda\_eq\_half} shows that a linear solution at $c=2$ forces $\lambda=\tfrac12$, by the local-minimum route rather than a two-sided limit, since $g(x)=2\lambda(x-1)-\log x$ is nonnegative on $(0,2)$ and vanishes at $1$, so its derivative $2\lambda-1$ vanishes there \textup{(\texttt{lean/Erdos307/NoInvariant.lean})}.

Any strengthening must therefore use an input about $n$ beyond the pair
$(\sigma(n),\sigma(n'))$, and by the next proposition there is no room in the choice of function
either.

\begin{proposition}[The ansatz is universal, so the search was the problem]\label{prop:lyapuniversal}
Let $S$ be the set of squarefree $n\ge2$ with $n'$ squarefree. The following are equivalent:
\begin{enumerate}[label=\textup{(\arabic*)},leftmargin=2.2em]
\item $n\mapsto n'$ has no periodic orbit inside $S$;
\item there is $\delta:S\to\mathbb{R}$ with $\delta(n')<\delta(n)$ whenever $n,n'\in S$;
\item there is $f$ on $\sigma(S)$ with $f(\sigma(n))-f(\sigma(n'))>\log\sigma(n)$ for all such $n$.
\end{enumerate}
\end{proposition}
\begin{proof}
(2) implies (1), since a periodic orbit would give $\delta(n)<\delta(n)$. For (1) implies (2), the
transitive closure of $n\mapsto n'$ is irreflexive by (1) and transitive by construction, hence a
strict partial order on the countable set $S$; it has a linear extension, and a countable linear
order embeds order-preservingly in $\mathbb{Q}$, which supplies $\delta$. For (2) and (3), $\sigma$
is injective on squarefree integers by Lemma~15.2 of the full version, so every $\delta$ on $S$ is
$\log n+f(\sigma(n))$ for exactly one $f$, namely $f(\sigma(n))=\delta(n)-\log n$; and $n'=n\sigma(n)$
turns $\delta(n')<\delta(n)$ into the inequality of (3).
\end{proof}

The collapse of the ansatz is formalised. \texttt{Erdos307.exists\_f\_of\_injective} proves that
injectivity of $\sigma$ forces \emph{every} $\delta$ into the shape $\log n+f(\sigma(n))$,
\texttt{increment\_iff} converts $\delta(n')<\delta(n)$ into $(\ast)$ verbatim, and
\texttt{no\_two\_cycle\_of\_decreasing} is the implication of
Proposition~\ref{prop:lyapworks}; all with $0$ \texttt{sorry} and the three standard axioms
\textup{(\texttt{lean/Erdos307/NoInvariant.lean})}.

So Definition~\ref{def:lyap} restricts nothing. The tournament that produced $\delta_\lambda$ was
searching the space of all functions on $S$, and by Proposition~\ref{prop:lyapuniversal} the
existence of any member of that space is equivalent to the absence of cycles, which is strictly
stronger than a negative answer to \#307. The criterion was never a weakening of the problem, and
Theorem~\ref{thm:lyapfalse} settles the one form of it that had been left open. What remains true,
and is the residue of the whole construction, is Theorem~\ref{thm:halflyap} together with
Corollary~\ref{cor:lyapbarrier}: an elementary Lyapunov proof of the barrier, valid exactly up to
mass $2$ and, by Corollary~\ref{cor:halfsharp}, not one step further.

\begin{remark}[Four independent blocks on $\mathrm{A}10$]\label{rem:threeblocks}
The existence route is therefore obstructed at four separate levels, and no two of them are the
same obstruction. Proposition~\ref{prop:nofreeslot} removes the freedom: the construction carries
one parameter and cannot be relaxed into a denser family. Theorem~\ref{thm:ff} removes the
method: no argument that survives the passage to $K[t]$ can prove existence, because there the
statement is false. Proposition~\ref{prop:untestable} removes the computation: even the $34$ most
reduced objects in the whole problem, the ones every sieve has failed to kill, cannot be examined.
Proposition~\ref{prop:multiplicityfail} removes the last refuge, the hope that a large family of
sequences is easier than one: the union stays logarithmically thin however many sequences it has.
What is left is a lower bound for primes in a sparse exponential sequence, which exists for no
non-degenerate linear recurrence at all. We state plainly that the present work does not decide the
existence direction and does not contain a route to deciding it.
\end{remark}

\begin{thebibliography}{9}
\bibitem{Barbeau1961} E.~J.~Barbeau, \emph{Remarks on an arithmetic derivative}, Canad.\ Math.\
Bull.\ \textbf{4} (1961), 117--122.
\bibitem{Barbeau1976} E.~J.~Barbeau, \emph{Computer challenge corner: Problem 477}, J.\ Recreational
Math.\ \textbf{9} (1976/77), 30.
\bibitem{ErdosGraham1980} P.~Erd\H{o}s and R.~L.~Graham, \emph{Old and New Problems and Results in
Combinatorial Number Theory}, Monogr.\ Enseign.\ Math.\ \textbf{28}, Geneva, 1980, p.~38.
\bibitem{GrahamEgyptian} R.~L.~Graham, \emph{Paul Erd\H{o}s and Egyptian fractions}, in \emph{Erd\H{o}s
Centennial}, Bolyai Soc.\ Math.\ Stud.\ \textbf{25}, Springer, 2013, pp.~289--309.
\bibitem{Watanabe2020} T.~Watanabe, \emph{New examples of the representation of $1$ by the sum of
reciprocals of semiprime numbers}, arXiv:2009.03275 (2020).
\bibitem{BloomShifted} T.~F.~Bloom, \emph{Unit fractions with shifted prime denominators}, Proc.\
Roy.\ Soc.\ Edinburgh Sect.\ A (2024); arXiv:2305.02689.
\bibitem{UfnarovskiAhlander2003} V.~Ufnarovski and B.~\AA hlander, \emph{How to differentiate a
number}, J.\ Integer Seq.\ \textbf{6} (2003), Article 03.3.4.
\bibitem{Kovic2012} J.~Kovi\v{c}, \emph{The arithmetic derivative and antiderivative}, J.\ Integer
Seq.\ \textbf{15} (2012), Article 12.3.8.
\bibitem{MO320838} MathOverflow, \emph{Can the product of two sums of prime reciprocals be $1$?},
question 320838 (comments by R.~Israel and J.~Rosen, answer by ``Woett''), 2019.
\url{https://mathoverflow.net/questions/320838}.
\bibitem{ErdosProblems307} T.~F.~Bloom, \emph{Erd\H{o}s problem \#307},
\url{https://www.erdosproblems.com/307}.
\bibitem{OEIS_A003415} OEIS Foundation, \emph{Sequence A003415 (arithmetic derivative)},
\url{https://oeis.org/A003415}.
\bibitem{teRiele1984} H.~J.~J.~te~Riele, \emph{Computation of all the amicable pairs below
$10^{10}$}, Math.\ Comp.\ \textbf{47} (1986), 361--368.
\bibitem{BHP2001} R.~C.~Baker, G.~Harman, and J.~Pintz, \emph{The difference between consecutive
primes, II}, Proc.\ London Math.\ Soc.\ (3) \textbf{83} (2001), 532--562.
\bibitem{Iwaniec1978} H.~Iwaniec, \emph{Almost-primes represented by quadratic polynomials},
Invent.\ Math.\ \textbf{47} (1978), 171--188.
\bibitem{LemkeOliver2012} R.~J.~Lemke Oliver, \emph{Almost-primes represented by quadratic
polynomials}, Acta Arith.\ \textbf{151} (2012), 241--261.
\bibitem{HardyWright} G.~H.~Hardy and E.~M.~Wright, \emph{An Introduction to the Theory of Numbers},
6th ed., Oxford Univ.\ Press, 2008.
\bibitem{Tenenbaum} G.~Tenenbaum, \emph{Introduction to Analytic and Probabilistic Number Theory},
3rd ed., Grad.\ Stud.\ Math.\ \textbf{163}, Amer.\ Math.\ Soc., 2015 (Erd\H{o}s--Wintner theorem).
\bibitem{Stothers} W.~W.~Stothers, \emph{Polynomial identities and Hauptmoduln}, Quart.\ J.\ Math.\
Oxford (2) \textbf{32} (1981), 349--370; R.~C.~Mason, \emph{Diophantine Equations over Function
Fields}, Cambridge Univ.\ Press, 1984.
\bibitem{Seva323194} Seva (\texttt{mathoverflow.net} user), \emph{A recursion with a number--theoretic
function}, MathOverflow question 323194 (2019). \url{https://mathoverflow.net/questions/323194}.
\bibitem{Bado2026} I.~O.~Bado, \emph{Structural constraints for an Erd\H{o}s unit--fraction problem
over primes}, preprint, Aix--Marseille University (May 2026).
\url{https://doi.org/10.13140/RG.2.2.32245.54245}.
\bibitem{Giuga1950} G.~Giuga, \emph{Su una presumibile propriet\`a caratteristica dei numeri primi},
Ist.\ Lombardo Sci.\ Lett.\ Rend.\ Cl.\ Sci.\ Mat.\ Nat.\ (3) \textbf{14} (1950), 511--528.
\bibitem{BBBG1996} D.~Borwein, J.~M.~Borwein, P.~B.~Borwein, and R.~Girgensohn, \emph{Giuga's conjecture
on primality}, Amer.\ Math.\ Monthly \textbf{103} (1996), no.~1, 40--50.
\bibitem{BorweinWong1997} J.~M.~Borwein and E.~Wong, \emph{A survey of results relating to Giuga's
conjecture on primality}, CRM Proc.\ Lecture Notes \textbf{11} (1997), 13--27.
\bibitem{BJM2000} W.~Butske, L.~M.~Jaje, and D.~R.~Mayernik, \emph{On the equation
$\sum_{p\mid N}1/p+1/N=1$, pseudoperfect numbers, and perfectly weighted graphs}, Math.\ Comp.\
\textbf{69} (2000), no.~229, 407--420.
\bibitem{GrauOllerMarcen2012} J.~M.~Grau and A.~M.~Oller-Marc\'en, \emph{Giuga numbers and the
arithmetic derivative}, J.\ Integer Seq.\ \textbf{15} (2012), Article 12.4.1; arXiv:1103.2298.
\bibitem{SondowMacMillan2017} J.~Sondow and K.~MacMillan, \emph{Primary pseudoperfect numbers,
arithmetic progressions, and the Erd\H{o}s--Moser equation}, Amer.\ Math.\ Monthly \textbf{124} (2017),
no.~3, 232--240; arXiv:1812.06566.
\bibitem{GrauOllerMarcenSadornil2021} J.~M.~Grau, A.~M.~Oller-Marc\'en, and D.~Sadornil,
\emph{On $\mu$-Sondow numbers}, arXiv:2111.14211 (2021); Acta Math.\ Hungar.\ (2022).
\bibitem{Cox} D.~A.~Cox, \emph{Primes of the Form $x^2+ny^2$: Fermat, Class Field Theory, and Complex
Multiplication}, 2nd ed., Wiley, 2013.
\bibitem{OEIS_A054377} OEIS Foundation, \emph{Sequence A054377 (primary pseudoperfect numbers)},
\url{https://oeis.org/A054377}.
\bibitem{OEIS_A007850} OEIS Foundation, \emph{Sequence A007850 (Giuga numbers)},
\url{https://oeis.org/A007850}.
\bibitem{OEIS_A396915} OEIS Foundation, \emph{Sequence A396915 (composite squarefree $k$ with
$k'+2k$ a perfect square)}, \url{https://oeis.org/A396915}.
\bibitem{Evertse1984} J.-H.~Evertse, \emph{On sums of $S$-units and linear recurrences}, Compositio
Math.\ \textbf{53} (1984), 225--244.
\bibitem{FI1998} J.~Friedlander and H.~Iwaniec, \emph{The polynomial $X^2+Y^4$ captures its
primes}, Ann. of Math. \textbf{148} (1998), 945--1040.
\bibitem{HB2001} D.~R.~Heath-Brown, \emph{Primes represented by $x^3+2y^3$}, Acta Math.
\textbf{186} (2001), 1--84.
\bibitem{Stewart2013} C.~L.~Stewart, \emph{On divisors of Lucas and Lehmer numbers}, Acta Math.
\textbf{211} (2013), 291--314.
\bibitem{BCZ2003} Y.~Bugeaud, P.~Corvaja, and U.~Zannier, \emph{An upper bound for the G.C.D. of
$a^n-1$ and $b^n-1$}, Math. Z. \textbf{243} (2003), 79--84.
\bibitem{Merikoski2019} J.~Merikoski, \emph{On the largest prime factor of $n^2+1$}, J. Eur. Math.
Soc. \textbf{25} (2023), 1253--1284.
\bibitem{BGS2010} J.~Bourgain, A.~Gamburd, and P.~Sarnak, \emph{Affine linear sieve, expanders,
and sundry applications}, Invent. Math. \textbf{179} (2010), 559--644.
\bibitem{ESS2002} J.-H.~Evertse, H.~P.~Schlickewei, and W.~M.~Schmidt, \emph{Linear equations in
variables which lie in a multiplicative group}, Ann.\ of Math.\ (2) \textbf{155} (2002), 807--836.
\bibitem{Erdos1952} P.~Erd\H{o}s, \emph{On the sum $\sum_{k=1}^x d(f(k))$}, J.\ London Math.\ Soc.\
\textbf{27} (1952), 7--15.
\bibitem{BHV2001} Y.~Bilu, G.~Hanrot, and P.~M.~Voutier, \emph{Existence of primitive divisors of
Lucas and Lehmer numbers} (with an appendix by M.~Mignotte), J.\ Reine Angew.\ Math.\ \textbf{539}
(2001), 75--122.
\end{thebibliography}

\end{document}
